\documentclass[12pt,letterpaper]{article}

\usepackage{amsmath,amssymb,amsfonts,mathtools,nccmath}
\usepackage{array,booktabs,multirow,threeparttable,adjustbox,hhline}
\usepackage{graphicx,svg,pdflscape,float}
\usepackage{cancel,eurosym,nicefrac}
\usepackage{sectsty,comment,footmisc,epigraph,enumerate,pifont}
\usepackage{diagbox,istgame}
\usepackage[table]{xcolor}
\usepackage{subcaption,caption}
\usepackage{palatino}
\usepackage{microtype}
\emergencystretch=2em
\usepackage[margin=1in]{geometry}
\usepackage[nodisplayskipstretch]{setspace}
\usepackage[indent]{parskip}
\usepackage{etoolbox}
\usepackage{natbib}
\bibpunct{(}{)}{;}{a}{}{,}
\usepackage{xurl}
\usepackage{hyperref}
\usepackage{cleveref}
\usepackage[section]{placeins}
\usepackage{needspace}

\usepackage{tikz}
\usepackage{pgfplots}
\usepackage{tikz-timing}
\usetikzlibrary{arrows.meta,intersections,fit,calc,decorations.pathreplacing,calligraphy}
\usetikzlibrary{patterns,patterns.meta}
\usetikzlibrary{
arrows.meta, intersections, fit, calc, decorations.pathreplacing, calligraphy }

\usepgfplotslibrary{fillbetween}
\pgfplotsset{compat=1.18}

\hypersetup{
colorlinks, linkcolor={red!50!black}, citecolor={blue!50!black}, urlcolor={blue!80!black} }

\onehalfspacing
\renewcommand{\arraystretch}{0.75}
\makeatletter
\patchcmd{\env@cases}{1.2}{0.75}{}{}
\makeatother

\newcommand{\cmark}{\text{\ding{51}}}
\newcommand{\xmark}{\text{\ding{55}}}
\newcommand{\wclear}{\succcurlyeq_{\text{clear}}}
\newcommand{\wunclear}{\preccurlyeq_{\text{clear}}}
\newcommand{\clear}{\succ_{\text{clear}}}
\newcommand{\unclear}{\prec_{\text{clear}}}
\newcommand{\bio}{\succcurlyeq_{\text{B}}}
\newcommand{\bioed}{\preccurlyeq_{\text{B}}}
\newcommand{\ac}{A}
\newcommand{\tp}{\theta^\mathcal{P}}
\newcommand{\btp}{\mathbf{t}^\mathcal{P}}
\newcommand{\DV}{\Delta V}
\newcommand{\DW}{\Delta P}
\newcommand{\Dg}{\Delta \gamma}
\newcommand{\DU}{\Delta U}
\newcommand{\DL}{\Delta \Lambda}
\newcommand{\DK}{\Delta K}
\newcommand{\Dk}{\Delta k}

\usepackage{ntheorem}
\theorempreskip{10pt}
\newtheorem{theorem}{Theorem}
\newtheorem{result}{Result}
\newtheorem{corollary}{Corollary}
\newtheorem{assumption}{Assumption}
\newtheorem{condition}{Condition}
\newtheorem{proposition}{Proposition}
\newtheorem{example}{Example}
\newtheorem{observation}{Observation}
\newtheorem{claim}{Claim}
\newtheorem{conjecture}{Conjecture}
\newtheorem*{definition}{Definition}
\newtheorem{remark}{Remark}
\newtheorem{lemma}{Lemma}
\newtheorem*{CLT}{The Lindeberg--L\'evy Central Limit Theorem}
\newtheorem*{RNB}{Theorem 1 in \citet{Mu2021FromAgain}}
\newtheorem*{FSD}{Theorem 1 in \citet{Pomatto2020StochasticNoise}}
\newtheorem*{CMR}{Corollary 1 in \citet{Milgrom1994ComparingEquilibria}}

\newenvironment{proof}[1][Proof]{\noindent\textbf{#1.} }{\ \rule{0.5em}{0.5em}}
\newcommand*{\Scale}[2][4]{\scalebox{#1}{$#2$}}%
\newcommand*{\Resize}[2]{\resizebox{#1}{!}{$#2$}}%
\DeclareMathOperator{\inte}{int}
\newcommand{\E}{\mathrm{E}}
\newcommand{\Var}{\mathrm{Var}}
\newcommand{\To}{\text{To }}
\newcommand{\st}{\text{s.t. }}
\newcommand{\dd}{\mathrm{d}}
\newcommand{\bt}{{t}}
\newcommand{\KL}{\mathrm{KL}}
\newcommand{\LR}{\mathrm{LR}}
\newcommand{\mbf}[1]{\mathbf{#1}}
\newcommand{\mbb}[1]{\mathbb{#1}}
\newcommand{\argmax}{\operatorname*{arg\,max}}
\newcommand{\argmin}{\operatorname*{arg\,min}}
\newcommand\scalemath[2]{\scalebox{#1}{\mbox{\ensuremath{\displaystyle #2}}}}

\newtheorem{hyp}{Hypothesis}
\newtheorem{subhyp}{Hypothesis}[hyp]
\renewcommand{\thesubhyp}{\thehyp\alph{subhyp}}
\newcommand{\red}[1]{{\color{red} #1}}
\newcommand{\blue}[1]{{\color{blue} #1}}
\newcolumntype{L}[1]{>{\raggedright\arraybackslash}m{#1}}
\newcolumntype{C}[1]{>{\centering\arraybackslash}m{#1}}
\newcolumntype{R}[1]{>{\raggedleft\arraybackslash}m{#1}}


% Compatibility for the imported strategic proofs.
\providecommand{\Tilde}{\widetilde}
\providecommand{\qed}{\hfill\ensuremath{\blacksquare}}

\usepackage{graphicx}
\newcommand{\epiquote}[2]{\textit{#1}\par\vspace{0.4ex}{\small #2}\par}
\newenvironment{sectionepigraph}{\begin{flushright}\begin{minipage}{\textwidth}\raggedleft}{\end{minipage}\end{flushright}\vspace{1ex}}
\begin{document}
\addtocontents{toc}{\protect\setcounter{tocdepth}{-1}}
\begin{titlepage}
\title{\Large Politics in an Unpredictable World: Realignment and Radicalization}
\author{Kun Heo\footnote{Post-doctoral Fellow, Department of Government, London School of Economics and Political Science, b.heo@lse.ac.uk}}
\date{\today\\[1.5ex] {\normalsize Word count: 15,211}\\[0.3ex] {\footnotesize Main text and footnotes; excluding references, tables, figures, displayed equations, and appendix}}
\maketitle
\begin{center}\textbf{MODEL DRAFT}\end{center}
\begin{abstract}
\noindent Since the early 2010s, radical politics has risen across democracies, often with the votes of former supporters of mainstream parties, while polarization has grown and partisan attachments have hardened. I offer a simple model in which all of these follow from one change, without manipulation or a change in voters' preferences. Citizens' preferences have two aspects: the policy that they want on average, and how they want policy to change when conditions change. Greater uncertainty or better information for the government makes the second aspect matter more. Voters realign, and a party whose average policy is farther from the median voter's can win because its response is closer to hers. The expected gap between the candidates' policies widens even though parties nominate more moderate candidates. Losing office becomes more costly, so voters become less willing to accept defeat, and scandals can move only few voters.

\medskip\noindent\textit{Keywords:} radical politics; polarization; realignment; uncertainty; information
\end{abstract}
\thispagestyle{empty}
\setcounter{page}{0}
\end{titlepage}


\section{Introduction}

\begin{sectionepigraph}
\epiquote{Constant revolutionising of production, uninterrupted disturbance of all social conditions,\\ everlasting uncertainty and agitation \ldots\ All that is solid melts into air.}{Karl Marx and Friedrich Engels, \textit{The Communist Manifesto} (1848)}
\end{sectionepigraph}
\vspace{1ex}
We live in a less predictable world, in which we do not know what we will see tomorrow as the state of the world evolves. Climate change and artificial intelligence (AI) are two examples. The climate is changing, and each year brings new climate disasters and record temperatures around the world. We know that things will get worse, but we do not know how far or how fast the change will go. The year 2024 was the warmest on record, and the heat of 2023 exceeded what climate models had predicted \citep{WMO2025,Schmidt2024}. Likewise, we are told that AI and other technological transformations will change the world, but how far and in which directions remains contested. Projections of their effect on the economy over the next decade range from a rise in productivity of less than one percent to a doubling of output \citep{Acemoglu2025}. Some worry that AI could cause calamitous harm, such as a collapse of human knowledge \citep{AcemogluKongOzdaglar2026}. Others expect AI to free people from hard labour and resource constraints \citep{SOURCE}.

I study how unpredictability about the future shapes politics. I use a spatial model of electoral competition in which two parties nominate candidates and voters choose between them. Parties care about policies and office benefits. The election outcome is stochastic for parties because voters are affected by valence shocks. Candidates cannot commit to policies, so voters judge each candidate by her characteristics, which determine what she will do in office.

I introduce two changes to the conventional framework. First, as in \citet{KVW2017}, the election winner observes informative events about the unknown state of the world after the election and before she chooses her policy. Voters and parties therefore base their decisions on how each candidate would respond to the events that may occur, given the distribution of these events.

Formally, unpredictability is the spread of the officeholder's forecast of the state across possible futures. The forecast is the summary of her beliefs that determines her policy. When the spread is larger, events that move her forecast further from the prior become more likely. Defined this way, unpredictability is equivalent to suspense in \citet{ElyFrankelKamenica2015}. Unpredictability is the product of two components: the variance of the underlying state and the explanatory power of the events about the state ($R^2$ in empirical analysis). The more heterogeneous the state is and the more informative the events are to the officeholder, the less predictable the world becomes. The statement that the informativeness of future events makes the world less predictable seems paradoxical. But unpredictability is defined ex ante, before the events are realized. If the events are not informative, we will see similar events under any state, so the events will barely move our beliefs. If the events are more informative, the events that we would see under different states overlap less, so what we see can move our beliefs far from the prior.

The second change is that each player (party, candidate, or voter) has two-dimensional characteristics (types). The first characteristic is her bliss point at the expected state (her baseline). The second is how she would respond to a change in the state (her responsiveness). Some people respond more strongly than others, so their preferred policy changes more as the state evolves. Two people who respond differently therefore disagree more, on average, after events that are far from what was expected.

An extension endogenizes unpredictability and allows the incumbent officeholder to raise or reduce it for the next term.

I argue that politics changes in the following ways when we are highly uncertain about what we will see, compared with a world in which we have a better sense of what we will see. First, voters realign across the conventional left-right spectrum: some former supporters of the left may support the right, and vice versa. The median voter may also support a more radical candidate, whose expected policy is farther from her preferred policy.

The reason is that players care more about how each candidate will respond to the state, relative to the policy that she is expected to choose. The expected policy does not depend on unpredictability. As a result, voters realign along the dimension of responsiveness, and their baseline preference over the expected policy becomes less important. The median voter can therefore prefer a radical politician whose expected policy is farther from hers than a moderate rival's, because the radical's responses to a changing world are closer to her own.

Second, the policy stakes of winning office rise for voters, parties, and candidates alike, holding the candidates fixed.

The higher stakes have two sources. The first is the assumption that players are risk averse, with quadratic losses. However, risk aversion alone is not enough. If the two candidates responded to future events in the same way, the loss from losing office would depend only on the difference in their baselines, which does not change with unpredictability. Only when the candidates respond differently do less predictable events add a loss from the gap in their responses on top of the baseline difference.

Higher stakes from unpredictability mean that voters, parties, and candidates lose more from defeat, so they become less willing to accept the defeat of their preferred candidate. At the same time, the higher stakes thin the centre near the median: the distribution of voters' preferences over the candidates becomes more bimodal, and the centre collapses. The thinner centre makes voters' decisions less sensitive to valence shocks that are orthogonal to policy, such as gaffes or scandals, because few voters remain close enough to indifference for such a shock to move them. News about how a candidate would respond moves more voters than before. The effects above are similar when policy becomes more important for every player by a multiplicative factor that does not involve unpredictability.

However, higher stakes from unpredictability have a distinctive effect on equilibrium nominations and policymaking. First, parties respond to the higher stakes by nominating more convergent candidates whose types are similar to the median. The same convergence occurs when policy becomes more important by a multiplicative factor. But the convergence does not make the candidates' policies more similar in expectation. In fact, their policies will diverge more despite the convergence in types. When the policy stake increases by a multiplicative factor, the policies do not diverge more. The policies diverge more because an unpredictable world is more likely to produce events far from what was expected, and the two candidates' policies differ more after such events. In short, a mechanical effect increases the policy divergence with the candidates held fixed, and the parties' strategic response to that mechanical effect (i.e., nominating candidates more similar to the median) reduces the divergence. The strategic effect falls short of compensating for the mechanical effect.

In the extension, suppose that an incumbent trails because her expected policy is farther from the median voter's than the challenger's. The incumbent can then win with probability greater than one half by raising unpredictability for the next term, provided that her responses are closer to those of decisive voters and she values office enough. The higher unpredictability lowers every player's policy payoff, and electing the challenger would not reduce the unpredictability.

The model's results are relevant to a sequence that recurs when a shock makes the future less predictable, and to the direction of the realignment within the sequence.

The sequence begins with a realignment that occurs in the midst of uncertainty and does not fit the conventional left-right divide. In France, harvest failures and a financial crisis in 1846 and 1847 preceded the February revolution of 1848, and \citet{Tocqueville2016} wrote that ``the collapse of commerce, ubiquitous hostility, and fear of socialism increasingly aroused hatred of the Republic.'' A new coalition, the Bonapartists, supported Louis Napoleon Bonaparte and did not fit the existing divide between left and right. Promising both order and an end to poverty, Bonaparte drew peasants, workers, and notables from both camps \citep{Price2004}. \citet{Lipset1960} described fascism as ``extremism of the centre.'' The supporters of fascism consisted mostly of the liberal middle class, small business owners and professionals \citep{King2008}. A similar pattern appears in recent events. In the aftermath of the 2008 financial crisis, the new coalition after the realignment includes many working-class and young people who used to be considered supporters of progressive politics.

The new cleavage that replaces class follows identities such as religion, gender, and education. In France in 1849, the new divide ran partly along religion: conservatives allied with the Catholic Church, and the democratic socialists were anticlerical \citep{Price2004}. In Weimar Germany, religion, mostly the divide between Catholics and Protestants, was the strongest predictor of the vote, ``dwarfing the explanatory power of any other demographic or socioeconomic variable'' \citep{SpenkuchTillmann2018}. Today, education is one of the strongest, if not the strongest, predictors of voters' choices in developed democracies \citep{Gethin2022,GrossmannHopkins2024,HoogheMarks2018}. Changes in financial well-being explained little of American voters' choices in 2016 \citep{Mutz2018}. In South Korea, young voters split by gender in the 2022 presidential election: most young men supported the conservative candidate and most young women the progressive one \citep{KimHanna2022}.

The centre then collapses, and positions harden on each side. The moderate republicans who had won the French Constituent Assembly in April 1848 collapsed in the legislative election of May 1849, and the new Assembly divided between the conservative Party of Order and the democratic socialists. In Weimar Germany, the liberal parties collapsed as the Nazis and the Communists grew. In the United States, fewer voters have been swayed by scandals over time.

A radical politician then wins with the coalition that the realignment produced. Bonaparte won the presidential election of December 1848 on a populist platform. Donald Trump won the presidency in 2016 with the support of working-class voters in the Rust Belt, many of whom had voted for Democrats before. In South Korea, Yoon Suk Yeol, a former prosecutor who had served under the left-leaning government, won the presidency in 2022 as the conservative candidate with many former supporters of the left.

Finally, the winners and their supporters become less willing to accept defeat. When the National Assembly refused in 1851 to amend the constitution that barred Bonaparte from a second term, he dissolved the Assembly in a self-coup. After Trump lost in 2020, he and many of his supporters refused to accept the result. In December 2024, Yoon declared martial law, and many of his supporters defended the declaration and opposed his removal.\footnote{See \citet{Price2004}, \citet{NohlenStover2010}, and \citet{Tocqueville2016} for the French Second Republic; \citet{King2008} and \citet{Evans2003} for Weimar Germany; \citet{Skelley2017} and \citet{Rottinghaus2025} for the United States; and \citet{Jang2025} and \citet{Jeong2025} for South Korea.}

In the model, every step follows from one change. Many values are rules about how one should act when circumstances change, so they are correlated with responsiveness and can serve voters as indicators of how a politician would respond to an evolving world. As the future becomes less predictable, voters weigh responses more than baselines, and cleavages that follow values, such as religion and education, gain at the expense of class. The same change hardens attachments, can give the lead to a radical candidate whose responses are closer to those of decisive voters, and raises the cost of losing. The realignment can be gradual, as the world becomes steadily less predictable, or sudden, when a shock makes the world less predictable at once.

The model's explanation of voter realignment differs from the existing one, which focuses on how structural change affected voters' material conditions, through a financial shock \citep{King2008,Funke2016,GurievPapaioannou2022} or a trade shock \citep{Autor2020,ColantoneStanig2018}. On that account, globalization, automation, and deindustrialization hollowed out the middle class, and mainstream parties, converging on similar programmes, failed to represent it. The failure of mainstream parties opened the door to populists of the left and the right who promised to protect these voters.

The existing account explains why some voters moved. It does not explain why the entire equilibrium shifted. \citet{Margalit2019} argues that the losers from recent structural changes are numerous enough to tip a close election but explain only a modest share of the populist vote. For instance, the effects on vote shares estimated by \citet{Autor2020} and \citet{ColantoneStanig2018} were a few percentage points.

The model's second result relates to the literature on polarization, which traces it to inequality, elite sorting, competition for majority control, activism, and partisan media \citep{McCarty2016,Bonica2014,Barber2016,Canen2020,Lee2016,Hall2019,Aldrich1983,Fiorina1999,MartinYurukoglu2017}. On top of these sources, uncertainty can widen the expected distance between policies, because politicians respond differently to heterogeneous states.

The model's third result concerns the stakes of office, the willingness to accept a loss, and the reduced effect of valence shocks. This result relates to the literatures on democratic stability and on voters' reactions to scandals. Existing work connects the willingness to accept a loss with polarization: \citet{Svolik2019} argues that polarization makes voters more willing to sacrifice democratic norms for their side's victory. Voters' failure to punish misconduct is usually attributed to their partisanship. In the model, misconduct is a valence shift orthogonal to policy. When uncertainty is high, misconduct can move few voters, because the gap in responses dominates voters' calculation.

\paragraph{Related formal literature.}
The paper builds on models of electoral competition between parties that value policy and office and that diverge when the electoral outcome is uncertain \citep{Wittman1983,Calvert1985,Roemer2001}, and on probabilistic voting \citep{Lindbeck1987}. In models of electoral competition, parties also diverge because of the entry of citizen candidates \citep{OsborneSlivinski1996,BesleyCoate1997}, valence advantages \citep{Groseclose2001,Schofield2007}, voters' uncertainty about how the economy works \citep{Roemer1994}, parties' private information about the state \citep{Schultz1996}, or positions targeted to supporters \citep{GlaeserPonzettoShapiro2005}. Here beliefs are common, and disagreement concerns how policy should respond to a state that the officeholder learns about after the election. The closest model is \citet{KVW2017}, in which information also reaches the winner after the election. In that model, preferences differ by a constant bias, so learning raises the value of discretion and leaves disagreement unchanged. In this paper, learning raises disagreement itself.

Candidates here have two-dimensional types. In two dimensions, a majority-rule equilibrium generally fails to exist \citep{Plott1967,McKelvey1976}, and a common point on every halving line requires special conditions \citep{DavisDeGrootHinich1972,McKelvey1986}. Each voter's comparison here reduces to a single index, and a unique equilibrium follows under party symmetry and a median point, or under local regularity (Section~\ref{sec:model}). Two-dimensional candidates also appear in models of competence and ideology \citep{BernhardtCamaraSquintani2011}, specialized candidates \citep{KrasaPolborn2010}, and valence \citep{BuisseretVanWeelden2022}. Here the second dimension is the slope of a preference over policy rules. The weight on the second dimension rises with uncertainty and information, and the electoral advantage comes from closeness to the decisive voter in responsiveness.

In a few models, voters realign. Voters move toward the candidates that they support \citep{CallanderCarbajal2022}, learn about their own positions, with each voter's comparison reduced to a single index as here \citep{Vaeth2024}, or adopt the parties' models of the world, including the slope of the relation between policy and outcomes \citep{IzzoMartinCallander2023}. In \citet{BonomiGennaioliTabellini2021} and \citet{GennaioliTabellini2025}, voters change their social identity. Each voter identifies with whichever predefined group, economic or cultural, stands in stronger policy conflict with its outgroup, and her beliefs tilt toward the group's stereotype. The electorate switches from class to cultural identity when the importance of cultural policy rises or a trade shock divides cultural groups. The groups stay fixed, voters move between them together, and the preferences that voters express change. On the supply side, activists produce realignment through shifting coalitions that rotate and separate the parties \citep{MillerSchofield2003}. In these models, voters' positions, beliefs, or identities change. Here a common shock changes the weight that voters place on responsiveness. Coalitions therefore rotate even with candidates held fixed, and policies diverge without a centrifugal force such as activists. Preferences and beliefs stay fixed, and no groups are defined in advance. Each voter switches when uncertainty passes her own threshold, and the coalitions that result need not follow any predefined line. Identity plays no causal role, and a shared identity matters at most as a marker of shared responses.

Most models of information in elections study voters' information \citep{Ashworth2012,BesleyBurgess2002,SnyderStromberg2010,Izzo2023,IzzoInvernizziCrosson2025}. Here voters' beliefs are fixed, and only the officeholder's policy information changes, as in accounts of state capacity \citep{Scott1998,Brambor2020}. The effect of information on disagreement depends on its source, and different predictions follow from heterogeneous priors \citep{KartikLeeSuen2021}, common observations \citep{DixitWeibull2007}, constant biases \citep{CrawfordSobel1982}, and designed public information \citep{AlonsoCamara2016}. Here better information raises expected policy disagreement because preferences differ in responsiveness.

Models of delegation and executive discretion study disagreement in level and slope when a principal can restrict an informed agent \citep{MelumadShibano1991,AlonsoMatouschek2008,AmadorBagwell2013}, or weigh responsiveness against stability \citep{HowellWolton2018,LoeperDziuda2024}. Here no one can restrict the officeholder, and a party selects on responsiveness, as a government selects a central banker for her weights on objectives \citep{Rogoff1985}.

In \citet{AlesinaCukierman1990}, an incumbent chooses uncertainty about her own policy preferences, so that voters cannot infer them. Here voters know the incumbent's type, and the incumbent raises uncertainty about the state. In \citet{DownsRocke1994}, leaders prolong wars to avoid removal from office. Models of democratic backsliding study incumbents who manipulate elections \citep{GrattonLee2024,HelmkeKroegerPaine2022,LuoPrzeworski2023} or remove constraints on their power \citep{HowellShepsleWolton2023,GrilloPrato2023,ChioprisNalepaVanberg2025}. All of these models take the incumbent's motive to stay in office as given. Here uncertainty and information strengthen that motive by making losing office more costly. In Section~\ref{sec:pre-election}, the incumbent leaves elections and constraints in place and raises uncertainty instead. Voters then prefer her because she is closer to decisive voters in responsiveness, and electing the challenger would not remove the uncertainty.

Section~\ref{sec:model} presents the model. Section~\ref{sec:policy} solves the policy choice and shows that polarization rises with fixed candidates. Section~\ref{sec:realign} derives voters' induced preferences over officeholders and studies realignment and hardening. Section~\ref{sec:competition} describes the parties' nominations in equilibrium. Section~\ref{sec:effects} studies comparative statics of the equilibrium in uncertainty, information, and office benefits, examines welfare, and ends with two corollaries on learning. Section~\ref{sec:pre-election} studies strategic uncertainty. Section~\ref{sec:evidence} reinterprets recent realignments in light of the model. Proofs and generalizations appear in the Supplementary Information (SI), whose Section~\ref{app:requirements} shows how the results build on one another, and SI Section~\ref{app:tests} proposes empirical tests.

\section{Benchmark model}\label{sec:model}

In the benchmark model, the officeholder learns about the state only after the election. In Section~\ref{sec:pre-election}, an incumbent can raise uncertainty strategically. At the end of Section~\ref{sec:effects}, the officeholder chooses how much to learn, and some learning occurs before nomination.

\paragraph{Players and preferences.}
There are two parties, $L$ and $R$, a continuum of voters, and candidates whom the parties nominate. In general, there are $N$ independent states and $K$ policy dimensions, and a type consists of a baseline for each policy and a responsiveness of each policy to each state. In the main text, $N=K=1$, and SI Section~\ref{app:multi} shows which results hold qualitatively in the general case. Every player $i$ has a type $\theta_i$, which determines her baseline and her response to the state, both defined below. In the main text, $\theta_i=(\alpha_i,\beta_i)\in\mathbb{R}^2$, where $\alpha_i$ is her responsiveness to the state and $\beta_i$ is her baseline. Voter types follow a distribution $G$ with compact, convex support. The density $g$ of $G$ is continuously differentiable and bounded away from zero on $\operatorname{supp}(G)$.

After the election, the winner chooses a policy $x\in\mathbb{R}$. The state $\omega\in\mathbb{R}$ has prior distribution $F_0$ with mean zero and finite variance $\operatorname{Var}(\omega)\equiv\operatorname{Var}_{F_0}(\omega)$. Player $i$'s ideal policy in state $\omega$ is
\[
\tau(\omega,\theta_i)=\beta_i+\alpha(\omega,\theta_i),\qquad \alpha\big(\mathbb{E}[\omega],\theta_i\big)=0,
\]
and her policy payoff is $-\big(x-\tau(\omega,\theta_i)\big)^2$. The baseline $\beta_i=\tau(\mathbb{E}[\omega],\theta_i)$ is the policy that she would want if tomorrow were exactly like today, that is, if the state took its expected value, and the response $\alpha(\omega,\theta_i)$ describes how her ideal policy departs from the baseline as the state moves away from that value. For exposition, the main text uses linear responses, $\alpha(\omega,\theta_i)=\alpha_i\,\omega$, so that
\[
\tau(\omega,\theta_i)=\beta_i+\alpha_i\,\omega .
\]
Baselines and responsiveness can be distributed across players in any way, and each officeholder's policy depends on the events that she observes only through her forecast of the state (Section~\ref{sec:policy}). Section~\ref{sec:choices} shows which results extend to other responses. Parties place a common weight $P_J>0$ on policy payoffs and voters weight $P_V>0$. The weights matter only relative to the office benefit and the electoral noise introduced below, and I set $P_J=P_V=1$ except in Proposition~\ref{prop:symmetric}(e).

\paragraph{Baseline and responsiveness.}
Consider public ownership of banks. One market liberal opposes nationalization in all circumstances. Another wants little public ownership in normal conditions but extensive nationalization during a financial collapse. The two market liberals share a low baseline and differ in responsiveness. One socialist wants extensive public ownership in all circumstances, and another wants public ownership to expand substantially as financial distress worsens. Comparing baselines groups the two market liberals together and the two socialists together. Comparing responses can bring the second market liberal closer to the second socialist.

Table~\ref{tab:examples} gives further examples. Voters who agree on limited executive authority in ordinary conditions may disagree about emergency powers under a severe threat. A voter who supports equal access to welfare for citizens and non-citizens in normal conditions may favour a large gap under scarcity. Asking only about current policy would miss this conditional preference for exclusion.

In each example, the baseline is the preferred policy at the expected state, the policy that one wants if tomorrow is like today. With linear responses, it is also the preferred policy on average. Responsiveness describes how the preferred policy changes as the state moves away from its expected value.

\begin{table}[t]
\centering
\small
\renewcommand{\arraystretch}{1.3}
\begin{tabular}{@{}L{3.6cm}L{3.2cm}L{3.8cm}L{3.8cm}@{}}
\toprule
{\bfseries\boldmath Policy} & {\bfseries\boldmath State} & {\bfseries\boldmath Baseline} & {\bfseries\boldmath Responsiveness} \\
\hline\hline
Public ownership of banks & Financial distress & Ownership at the reference level & Rise as distress worsens \\ \midrule
Executive authority & Threats to order & Powers at the reference level & Expansion as threats grow \\ \midrule
Citizen priority in welfare or jobs & Scarcity & Citizen--non-citizen gap at the reference level & Widening as scarcity grows \\ \midrule
Immigration restrictions & Migration pressure & Restrictions at the reference level & Tightening as pressure grows \\ \midrule
Tariffs or domestic-content rules & External competition & Protection at the reference level & Rise with competition \\ \midrule
Public investment in restructuring & Technological disruption & Involvement at the reference level & Expansion with disruption \\
\bottomrule
\end{tabular}
\caption{Examples of baseline and responsiveness. Each row is a separate domain with its own state and units; the sign of responsiveness depends on coding.}
\label{tab:examples}
\end{table}

\paragraph{Information.}
No information about $\omega$ arrives before the election. After winning, the officeholder observes a signal. The signal can be read as the events that she observes during her term, such as a storm season, a wave of layoffs, or a surge at the border, each of which is informative about a latent state that policy must address. Let $\mu$ be her posterior over $\omega$, and let $F_\mu$ be the distribution of posteriors that her signal induces, which all players know.

\paragraph{Parties and candidates.}
Party $J\in\{L,R\}$ values policy and winning office. Party $J$ has ideal type $\theta_J=(\alpha_J,\beta_J)\in\operatorname{supp}(G)$ and can nominate a citizen of any type in that support. A party can also be read as a leader who shapes nominations. Trump appears to have played this role through his endorsements in the 2022 Republican primaries, where endorsed candidates received about ten percentage points more of the vote \citep{BlumCowburnMasket2024}. Write $\theta_l$ and $\theta_r$ for the nominees of $L$ and $R$. The nominees' types are common knowledge. The types represent what voters learn from the nominees' records and communications, and strategic communication about types is left for future work. The nominees' midpoint has coordinates $\alpha_m=(\alpha_l+\alpha_r)/2$ and $\beta_m=(\beta_l+\beta_r)/2$.

Each party also receives an office benefit $B\ge0$ when its nominee wins. The benefit is independent of policy, the state, and the nominee's type, and represents benefits of governing beyond the policies implemented. The parties receive the same benefit in the benchmark model; Section~\ref{sec:pre-election} allows different benefits.

\paragraph{Elections.}
Before voting, all voters experience common electoral noise $\eta\sim\operatorname{Unif}\big[-\tfrac{1}{2\phi},\tfrac{1}{2\phi}\big]$, where $\phi>0$, which shifts every voter's comparison of the candidates by the same amount. Voter $i$ votes for $l$ if and only if her expected policy payoff from electing $l$ exceeds her expected policy payoff from electing $r$ by at least $\eta$, and the candidate with a majority wins. Let $\pi\equiv\Pr(l\text{ wins})$, and let $x_j$ be the policy that candidate $j$ implements if she wins. The parties evaluate nominations by
\begin{align*}
\Pi_L(\theta_l,\theta_r)&=\pi\Big(\mathbb{E}\big[-\big(x_l-\tau(\omega,\theta_L)\big)^2\big]+B\Big)+(1-\pi)\,\mathbb{E}\big[-\big(x_r-\tau(\omega,\theta_L)\big)^2\big],\\
\Pi_R(\theta_l,\theta_r)&=\pi\,\mathbb{E}\big[-\big(x_l-\tau(\omega,\theta_R)\big)^2\big]+(1-\pi)\Big(\mathbb{E}\big[-\big(x_r-\tau(\omega,\theta_R)\big)^2\big]+B\Big).
\end{align*}

\paragraph{Timing.}
\begin{enumerate}
  \item Parties simultaneously nominate candidates $\theta_l,\theta_r\in\operatorname{supp}(G)$.
  \item The noise $\eta$ is realized, voters vote, and the candidate with a majority wins.
  \item The state $\omega$ is realized, the winner observes the signal, forms posterior $\mu$, and implements $x^*_j(\mu)$.
  \item Payoffs are realized.
\end{enumerate}
I solve the game by backward induction: the policy choice (Section~\ref{sec:policy}), the vote (Section~\ref{sec:realign}), and the nominations (Section~\ref{sec:competition}).

\paragraph{Equilibrium.}
The solution concept is weak perfect Bayesian equilibrium, with sincere voting. The winner chooses the policy that maximizes her expected payoff given her posterior, which she forms by Bayes' rule. Each voter votes for the candidate whose election gives her the higher expected payoff. Each party's nomination maximizes its expected payoff given the other party's nomination. I refer to an equilibrium by its nominations $(\theta_l^*,\theta_r^*)$. In the cases that the paper studies, the equilibrium is unique; Propositions~\ref{prop:symmetric} and~\ref{prop:asymmetric} state the conditions.
Table~\ref{tab:notation} summarizes the notation of the main text. Table~\ref{tab:notation-si} in the SI lists all other notation, including equilibrium quantities, so every symbol can be found in one of the two tables.

\begin{table}[t]
\centering
\footnotesize
\renewcommand{\arraystretch}{1.05}
\setlength{\aboverulesep}{0pt}\setlength{\belowrulesep}{0pt}\setlength{\extrarowheight}{2pt}
\begin{adjustbox}{max width=0.9\textwidth,center}
\begin{tabular}{@{}L{2.9cm}L{4.4cm}|L{2.7cm}L{4.4cm}@{}}
\toprule
{\bfseries\boldmath Symbol} & {\bfseries\boldmath Meaning} & {\bfseries\boldmath Symbol} & {\bfseries\boldmath Meaning} \\
\hhline{==|==}
\multicolumn{2}{l|}{\textbf{Preferences and the state}} & \multicolumn{2}{l}{\textbf{Players and payoffs}} \\
\hhline{==|==}
$\theta_i=(\alpha_i,\beta_i)$ & Type of player $i$ & $G$, $g$ & Distribution and density of voter types \\
\hhline{--|--}
$\beta_i$ & Baseline: preferred policy at the reference state and on average & $\theta_L,\theta_R$; $\theta_l,\theta_r$ & Party ideal types; nominees \\
\hhline{--|--}
$\alpha_i$ & Responsiveness: change in preferred policy with the state & $\theta_m$; $\Delta\alpha$, $\Delta\beta$ & Nominees' midpoint; differences $\alpha_l-\alpha_r$, $\beta_l-\beta_r$ \\
\hhline{--|--}
$\tau(\omega,\theta_i)=\beta_i+\alpha_i\omega$ & Ideal policy of player $i$ in state $\omega$ & $U_i^j$ & $i$'s expected policy payoff if $j$ governs \\
\hhline{--|--}
$\omega$ & State, mean zero, realized after the election & $\Delta U_i=U_i^l-U_i^r$ & $i$'s payoff gap \\
\hhline{--|--}
$\operatorname{Var}(\omega)$ & Uncertainty & $\Lambda_i^j$; $D$ & Cost of delegation: expected squared policy gap; polarization $D=\Lambda_l^r$ \\
\hhline{--|--}
\multicolumn{2}{l|}{\textbf{Information}} & $B$; $V_J$ & Office benefit; total value of office \\
\hhline{==|--}
$F_0$; $\mu$; $F_\mu$ & Prior distribution of the state; posterior; distribution of posteriors & \multicolumn{2}{l}{\textbf{Elections}} \\
\hhline{--|==}
$S=\operatorname{Var}_{F_\mu}(\mathbb{E}_\mu[\omega])$ & Unpredictability of future events: variance of the officeholder's forecast & $\eta$; $\phi$ & Electoral noise; electoral sensitivity \\
\hhline{--|--}
$R^2=S/\operatorname{Var}(\omega)$ & Revealed share: $R^2$ of the officeholder's forecast & $M$; $\pi$ & Electoral margin (tie value of the noise); $l$'s winning probability \\
\hhline{--|--}
$P_J$; $P_V$ & Weights on policy of parties; voters & $\Pi_J$ & Party $J$'s payoff \\
\bottomrule
\end{tabular}
\end{adjustbox}
\caption{Notation.}
\label{tab:notation}
\end{table}

\subsection{Modeling choices}\label{sec:choices}

\paragraph{Quadratic loss.} Quadratic loss reduces the information structure to a single index, the unpredictability of future events $S$ (Section~\ref{sec:policy}). The effect of information needs no further assumption. With quadratic loss, each officeholder's policy is the posterior mean of her target, so the disagreement between two players is the expected value of a convex function of the posterior. Better information raises that expected value for any targets (Lemma~\ref{lem:targets}).

\paragraph{Linear responses.} Linear responses are for exposition. With linear responses, the policy that a player wants at the expected state is also the policy that she wants on average. The results do not depend on this coincidence. With nonlinear responses, the two policies differ by $\mathbb{E}[\alpha(\omega,\theta_i)]$, and comparisons between candidates depend on the expected policies. If responses share a shape and differ in intensity, $\alpha(\omega,\theta_i)=\alpha_i h(\omega)$ with $h(\mathbb{E}[\omega])=0$, every result holds with the expected policy $\beta_i+\alpha_i\mathbb{E}[h(\omega)]$ in place of the baseline, $S$ replaced by $\operatorname{Var}_{F_\mu}\big(\mathbb{E}_\mu[h(\omega)]\big)$, the unpredictability of the policy-relevant content of events, and $\operatorname{Var}(\omega)$ replaced by $\operatorname{Var}\big(h(\omega)\big)$ (SI). A change in uncertainty can then also move the expected policies, an effect on levels that linear responses rule out. The expected squared change in an officeholder's policy from the policy that she would choose if tomorrow were like today then splits into a predictable drift, $\alpha_j^2\big(\mathbb{E}[h(\omega)]\big)^2$, and an unpredictable spread, $\alpha_j^2\operatorname{Var}_{F_\mu}\big(\mathbb{E}_\mu[h(\omega)]\big)$, and unpredictability is the second. With responses of different shapes, information still raises every cost of delegation (Lemma~\ref{lem:targets}), and each voter's comparison of the candidates still depends on her response only through one projection, its covariance with the candidates' response gap. The information environment is then no longer summarized by a single number, and the clean properties of realignment, such as a single switch for each voter, need not hold.

\paragraph{Affine targets.} Affine targets separate baseline from responsiveness, and the effect of uncertainty depends on them. Disagreement between two players is the expected square of the gap between their policies (Figure~\ref{fig:delegation}). Under full information, greater uncertainty raises it whenever the squared gap between the targets is convex in the state, because a mean-preserving spread then raises its expected value. With partial information and targets that are not affine, the revealed share no longer summarizes the information relevant to the gap, and convexity alone does not order disagreement across priors and signals. Affine targets satisfy the condition, and for them the result holds with partial information: the gap grows linearly as the state moves away from the point where the targets cross. If the squared gap is instead largest near the mean of the state, as for targets that saturate at different thresholds, the variance of the state does not determine the effect. A uniform dispersion makes every period more turbulent and lowers disagreement. A local dispersion leaves most periods calm, makes rare crises more severe, and can raise disagreement (Proposition~\ref{prop:dispersion} and Figure~\ref{fig:targets}).

\paragraph{Several states.} Let the state be $\omega=(\omega_1,\dots,\omega_N)$ with independent components, such as financial distress, migration pressure, and technological disruption, and let policy $n$ respond to $\omega_n$. Player $i$'s type is $\theta_i=(\beta_{in},\alpha_{in})_{n=1}^N$, her target for policy $n$ is $\tau_n(\omega_n,\theta_i)=\beta_{in}+\alpha_{in}\omega_n$, and $S_n$ is the unpredictability of events about $\omega_n$. The cost of delegation is then
\[
\Lambda_i^j=\sum_{n=1}^N\Big[(\beta_{jn}-\beta_{in})^2+(\alpha_{jn}-\alpha_{in})^2S_n\Big],
\]
and the main results hold with these sums (Propositions~\ref{prop:multistate} and~\ref{prop:issues}). A rise in $S_n$ raises the weight on $\alpha_{in}$ in each voter's comparison and leaves the weights on $\beta_{in'}$ and on $\alpha_{in'}$, $n'\ne n$, unchanged, so the dividing line rotates only in the dimension of $\alpha_n$. Realignment after a shock to $\omega_n$ therefore predicts realignment after a shock to $\omega_{n'}$ only if $\alpha_{in}$ and $\alpha_{in'}$ are correlated across voters: a voter who wanted a drastic response to the financial crisis need not want one to artificial intelligence. The prediction also depends on the candidates and on $S_{n'}$ for the other states. The SI extends the results to a responsiveness matrix, in which one state bears on several policies, as migration pressure bears on border enforcement, welfare, and the labour market, and to correlated states. SI Section~\ref{app:tests} uses differences in exposure to these states across regions to derive regional predictions.

\paragraph{Distribution of voters.} Propositions~\ref{prop:auto-polar} and~\ref{prop:realign} compare fixed candidates and hold for any distribution of voters. The equilibrium analysis uses a median point to obtain closed forms. The median point keeps the decisive voter fixed while the dividing line rotates. Without a median point, the shock can also change the decisive voter, which adds a second source of realignment. The main comparative statics then hold qualitatively under moderate electoral sensitivity and local regularity (Proposition~\ref{prop:general} and Corollary~\ref{cor:leadgeneral}). Local regularity requires distinct, interior party ideals and an electoral margin that is smooth in the candidates' types near those ideals, with unpredictability bounded away from zero (Assumption~\ref{ass:local}).

\paragraph{Electoral noise.} The common electoral noise stands for events that shift every voter's evaluation of the candidates, such as a scandal or the state of the economy on election day. Uniform noise gives closed forms. The comparative statics also hold with log-concave noise (Proposition~\ref{prop:logconcave}) and with idiosyncratic noise under symmetry conditions (SI). Table~\ref{tab:general} in the SI lists the conditions behind each result.

\section{Policymaking}\label{sec:policy}

I solve the game backward and start with the officeholder's policy choice.

\paragraph{Policy choice.}
Candidates cannot commit to policy. After winning and forming posterior $\mu$, candidate $j\in\{l,r\}$ chooses
\[
x^*_j(\mu)=\arg\max_{x}\;\mathbb{E}_\mu\!\left[-\big(x-\tau(\omega,\theta_j)\big)^2\right]=\mathbb{E}_\mu\big[\tau(\omega,\theta_j)\big]=\beta_j+\alpha_j\,\mathbb{E}_\mu[\omega].
\]
The same formula gives $x^*_i(\mu)$, the policy that any player $i$ would choose with the officeholder's information. Before the state and the signal are realized, player $i$'s expected policy payoff from electing $j$ is
\[
U_i^j=\int \mathbb{E}_\mu\!\left[-\big(x^*_j(\mu)-\tau(\omega,\theta_i)\big)^2\right]\mathrm{d}F_\mu(\mu),
\]
and her payoff gap is $\Delta U_i\equiv U_i^l-U_i^r$. The parties' payoffs become $\Pi_L=\pi\big(U_L^l+B\big)+(1-\pi)U_L^r$ and $\Pi_R=\pi U_R^l+(1-\pi)\big(U_R^r+B\big)$.

\paragraph{Unpredictability of future events.}
The officeholder's policy moves with her forecast $\mathbb{E}_\mu[\omega]$, which summarizes what the events that she observes tell her about the state. I define the \emph{unpredictability of future events} as the variance of this forecast across possible futures,
\[
S(F_0,F_\mu)\equiv\operatorname{Var}_{F_\mu}\!\big(\mathbb{E}_\mu[\omega]\big)\in\big[0,\operatorname{Var}(\omega)\big],\qquad R^2(F_0,F_\mu)\equiv\frac{S(F_0,F_\mu)}{\operatorname{Var}(\omega)}\in[0,1].
\]
I drop the arguments unless they are needed, write $S$ and $R^2$, and refer to $S$ as unpredictability and to $R^2$ as the revealed share. Unpredictability measures how much the policy-relevant content of future events varies. When $S$ is low, we can anticipate what we will see; when it is high, the events could tell us very different things about the state, and the officeholder's policy will vary with them. Uncertainty measures how far the state can move from its mean: in the first example of Table~\ref{tab:examples}, how severe financial distress can become. The revealed share measures the quality of the officeholder's information, the fraction of the state's variance that her signal explains: how much of the distress that bank supervision and stress tests reveal in time. Unpredictability equals the product of the two, $S=R^2\operatorname{Var}(\omega)$, and it rises when the officeholder's information improves in the Blackwell order. A more informative signal leaves less uncertainty about the state once the officeholder has learned, but it makes what she will see, and hence her policy, harder to predict in advance.

Unpredictability, so defined, has three equivalent readings. Because the mean forecast is the prior mean, $S=\mathbb{E}\big[(\mathbb{E}_\mu[\omega]-\mathbb{E}[\omega])^2\big]$ is the expected squared movement of the forecast between the election and the time policy is made, so it measures how little we can anticipate today what we will believe tomorrow. The first reading is the suspense of \citet{ElyFrankelKamenica2015} applied to the forecast: they measure suspense by the variance of next period's beliefs and extend the measure to an agent who cares only about the expectation of the state.\footnote{In their baseline, \citet{ElyFrankelKamenica2015} take the square root of the variance of next period's beliefs, so $S$ is the square of their measure in the version for expectations, and the two measures rank information structures in the same way. Their extension assumes a bounded state, which their results need and the definition does not. With several states and policies (SI Section~\ref{app:multi}), unpredictability is a matrix, and the measure that corresponds to suspense in a comparison between two players is the suspense about the gap between their policies.} By the law of total variance, $S=\operatorname{Var}(\omega)-\mathbb{E}\big[\operatorname{Var}_\mu(\omega)\big]$ is also the expected amount of uncertainty that will be resolved by then, so the unpredictability of tomorrow's forecast equals how much we will learn. Finally, $\alpha_j^2S$ is the expected squared change in officeholder $j$'s policy from the policy that she would choose if tomorrow were like today (Lemma~\ref{lem:volatility}), so unpredictability also measures how much policy will change. A future about which we will learn nothing new is predictable in this sense, however uncertain it is, and a future that will teach us much is unpredictable. The uncertainty that remains unresolved, $\mathbb{E}[\operatorname{Var}_\mu(\omega)]$, cannot be acted on and costs every player the same whoever governs, so $S$ is the part of uncertainty that politics responds to.

Throughout the paper, a change in uncertainty holds the revealed share fixed. I reserve the word \emph{shock} for a change in uncertainty or in the revealed share, that is, a change in the distribution of the state or of the officeholder's signal. Realizations of the state, of the signal, and of the electoral noise introduced below are not shocks in this sense. Lemma~\ref{lem:payoff} shows that both kinds of shock affect every comparison between candidates only through unpredictability. Two information environments with the same $S$ are therefore politically equivalent. Figure~\ref{fig:same-effect} shows the two routes to greater unpredictability.

\begin{figure}[t]
\centering
\begin{tikzpicture}[x=1.45cm, y=2.5cm,
  ebox/.style={draw, rounded corners, fill=white, align=center, inner sep=4pt, font=\small},
  earr/.style={->, thick}]
    \draw[->] (-3.4,0) -- (3.5,0) node[right, font=\footnotesize, align=left] {officeholder's\\ forecast $\mathbb{E}_\mu[\omega]$};
    % before the shock: low unpredictability
    \draw[gray!65, thick, domain=-3.2:3.2, samples=80] plot (\x, {0.798*exp(-\x*\x/(2*0.25))});
    \node[font=\footnotesize, gray!70!black, anchor=west] at (0.45,0.6) {before: low $S$};
    \draw[<->, gray!65] (-0.5,-0.09) -- (0.5,-0.09);
    % after the shock: high unpredictability, reached by either route
    \draw[black, very thick, domain=-3.2:3.2, samples=80] plot (\x, {0.399*exp(-\x*\x/2)});
    \node[font=\footnotesize, anchor=south west] at (1.95,0.05) {after: high $S$};
    \draw[<->, thick] (-1,-0.19) -- (1,-0.19) node[right, font=\scriptsize] {unpredictability $S$};
    \node[ebox, text width=3.3cm] (u) at (-2.6,1.12) {\textbf{uncertainty rises}\\ \footnotesize $\operatorname{Var}(\omega)\uparrow$, revealed share $R^2$ fixed};
    \node[ebox, text width=3.3cm] (b) at (2.7,1.12) {\textbf{the signal improves}\\ \footnotesize $\operatorname{Var}(\omega)$ fixed, $R^2\uparrow$};
    \draw[earr] (u.south) to[out=-90, in=160] (-1.25,0.19);
    \draw[earr] (b.south) to[out=-90, in=20] (1.25,0.19);
\end{tikzpicture}
\caption{Why uncertainty and information have the same political effect. Greater uncertainty with the revealed share fixed, and a more informative signal with uncertainty fixed, both widen the distribution of the officeholder's forecast $\mathbb{E}_\mu[\omega]$. The variance of that distribution is $S=R^2\operatorname{Var}(\omega)$, and the officeholder's policy $\beta_j+\alpha_j\mathbb{E}_\mu[\omega]$ moves more with the state. Schematic.}
\label{fig:same-effect}
\end{figure}

\paragraph{Policy volatility and polarization.}
For any quantity $z$ that differs across players, write $\Delta z_{ij}\equiv z_i-z_j$, and drop the subscripts for the two candidates: $\Delta z\equiv\Delta z_{lr}$. In particular, $\Delta\alpha=\alpha_l-\alpha_r$, $\Delta\beta=\beta_l-\beta_r$, and $\Delta x^*(\mu)=x^*_l(\mu)-x^*_r(\mu)$.

\begin{lemma}[Policy volatility]\label{lem:volatility}
For any officeholder $j$, $\operatorname{Var}_{F_\mu}\big(x^*_j(\mu)\big)=\alpha_j^2 S$, and for any two officeholders $j$ and $k$, $\operatorname{Var}_{F_\mu}\big(x^*_j(\mu)-x^*_k(\mu)\big)=(\Delta\alpha_{jk})^2 S$. A rise in $S$ therefore increases how much the policy of any officeholder whose ideal policy responds to the state ($\alpha_j\neq0$) varies across events, and how much the gap between two officeholders' policies varies whenever their responses differ ($\Delta\alpha_{jk}\neq0$).
\end{lemma}

Both parts follow from $x^*_j(\mu)=\beta_j+\alpha_j\mathbb{E}_\mu[\omega]$. I measure polarization by the expected squared distance between the candidates' policies, which is the squared distance between their expected policies plus the variance of the gap.

\begin{proposition}[Automatic polarization]\label{prop:auto-polar}
Fix the candidates. Polarization,
\[
D(\theta_l,\theta_r)\equiv\int\big(\Delta x^*(\mu)\big)^2\,\mathrm{d}F_\mu(\mu)=(\Delta\beta)^2+(\Delta\alpha)^2\,S ,
\]
rises with $S$ whenever $\Delta\alpha\ne0$. The squared distance between the candidates' expected policies, $\big(\mathbb{E}[\Delta x^*(\mu)]\big)^2=(\Delta\beta)^2$, does not depend on $S$.
\end{proposition}

Greater unpredictability makes each candidate's policy vary more with events (Lemma~\ref{lem:volatility}). The candidates move at different rates, $\alpha_l$ and $\alpha_r$, so their policies are far apart in extreme states, and the expected squared distance rises. The distance between the expected policies keeps only the baseline term $(\Delta\beta)^2$. Measures of positions capture this distance, so polarization can rise while the measured distance between positions stays constant.

\paragraph{Value of office.}
A party's total value of office is the office benefit plus the policy value of winning: $V_L\equiv B+\Delta U_L$ for party $L$ and $V_R\equiv B-\Delta U_R$ for party $R$.

\begin{corollary}[Value of office]\label{cor:office}
Holding the office benefit fixed, $V_L$ rises with $S$ if and only if $|\Delta\alpha_{lL}|<|\Delta\alpha_{rL}|$, and symmetrically for party $R$.
\end{corollary}

When its own candidate's responsiveness is closer to its own than the rival candidate's is, a party loses more from the rival's victory as unpredictability rises, because the rival's policy then departs further from the party's target in extreme states.

\section{Induced preferences over officeholders: realignment and hardening}\label{sec:realign}

Players have preferences over policies, and these preferences induce preferences over officeholders, that is, over the types who would choose policy. I call the second kind induced preferences. In this section, the candidates $\theta_l$ and $\theta_r$ are fixed, and I examine how uncertainty and information change voters' induced preferences between them. Proofs appear in the SI.

\begin{lemma}[Induced preferences]\label{lem:payoff}
For any types $\theta_i$ and $\theta_j$ and any distribution of posteriors $F_\mu$,
\[
U_i^j=-\alpha_i^2\operatorname{Var}(\omega)-(\Delta\beta_{ij})^2+\alpha_j(2\alpha_i-\alpha_j)\,S,
\]
and
\[
\Delta U_i=2\Delta\beta\,(\beta_i-\beta_m)+2\Delta\alpha\,S\,(\alpha_i-\alpha_m).
\]
If $\beta_r>\beta_l$, voter $i$ votes for $l$ if and only if $\Delta U_i\ge\eta$, that is, if and only if $\beta_i\le\psi_\eta(\alpha_i)$, where
\[
\gamma\equiv\frac{\Delta\alpha\,S}{\Delta\beta},\qquad
\psi(\alpha)\equiv\beta_m-\gamma\,(\alpha-\alpha_m),\qquad
\psi_\eta(\alpha)\equiv\psi(\alpha)+\frac{\eta}{2\Delta\beta}.
\]
\end{lemma}

The first term of $U_i^j$, $-\alpha_i^2\operatorname{Var}(\omega)$, is the same whichever candidate governs, so it cancels from every comparison between candidates. Comparisons therefore depend on the information environment only through $S=R^2\operatorname{Var}(\omega)$. Greater uncertainty with the revealed share fixed and a higher revealed share with uncertainty fixed both raise $S$, so both move every comparison in the same direction, in proportion to $R^2$ and to $\operatorname{Var}(\omega)$ respectively. Uncertainty enters separately only through payoff levels, which matter for welfare (Section~\ref{sec:effects}).

Without noise, the voters who are indifferent between the candidates lie on the line $\psi$, which passes through the candidates' midpoint $\theta_m=(\alpha_m,\beta_m)$ and has slope $-\gamma$. The line passes through the midpoint because a voter at that point is equally distant from both candidates in each dimension. The noise shifts the line vertically without changing its slope, which depends only on $S$.

For any players $i$ and $j$, define the cost of delegation as the expected squared gap between the policies that they would choose,
\[
\Lambda_i^j\equiv\int\big(\Delta x^*_{ij}(\mu)\big)^2\,\mathrm{d}F_\mu(\mu).
\]
Polarization is the cost of delegation between the two candidates, $D=\Lambda_l^r$.

\begin{proposition}[Realignment and hardening]\label{prop:realign}\label{prop:office}
Fix candidates with $\beta_r>\beta_l$.
\begin{enumerate}
  \item[(a)] \emph{Realignment.} If $\Delta\alpha\ne0$, then as $S$ rises, $|\gamma|$ rises and $\psi$ rotates around $\theta_m$, from the horizontal line $\beta_i=\beta_m$ toward the vertical line $\alpha_i=\alpha_m$. If $\alpha_r>\alpha_l$, a voter with $\beta_i<\beta_m$ and $\alpha_i>\alpha_m$ switches from $l$ to $r$ once $S$ passes her threshold, provided that the threshold is attainable; a voter with $\beta_i>\beta_m$ and $\alpha_i<\alpha_m$ switches from $r$ to $l$ under the same condition; and each voter switches at most once.
  \item[(b)] \emph{Hardening.} For any players $i$ and $j$, $U_i^j=U_i^i-\Lambda_i^j$, where $U_i^i=-\alpha_i^2\big(\operatorname{Var}(\omega)-S\big)$ is $i$'s payoff if she governed herself and
  \[
  \Lambda_i^j=(\Delta\beta_{ij})^2+(\Delta\alpha_{ij})^2\,S .
  \]
  The cost of being represented by any $j$ with $\Delta\alpha_{ij}\ne0$ rises linearly in $S$, with slope $(\Delta\alpha_{ij})^2$.
\end{enumerate}
\end{proposition}

Figure~\ref{fig:realign} illustrates a reversal of the majority. Candidate $l$ is closer to the electorate's centre in baseline, and $r$ is closer in responsiveness. With no information and a uniform distribution on the square, 70 percent prefer $l$. As $S$ rises, voters in the lower-right region switch to $r$ because they prefer $r$'s response; these voters still prefer $l$'s baseline. Voters in the upper-left region switch in the opposite direction. The first group is larger, so only 45 percent prefer $l$ at $S=2$. If the two groups had equal mass, the coalitions would change without a change in vote shares.

\begin{figure}[t]
\centering
\begin{tikzpicture}[scale=3.4, every node/.style={font=\small}]
  % preferences when I_q = 2, with switchers (relative to I_q = 0) in lighter shades
  \fill[blue!22] (0,-1) -- (1.4,-1) -- (0.7,0.4) -- (0,0.4) -- cycle;
  \fill[red!22]  (0.4,1) -- (2,1) -- (2,0.4) -- (0.7,0.4) -- cycle;
  \fill[red!8]   (0.7,0.4) -- (2,0.4) -- (2,-1) -- (1.4,-1) -- cycle;
  \fill[blue!8]  (0,0.4) -- (0.7,0.4) -- (0.4,1) -- (0,1) -- cycle;
  % frame and axes
  \draw[gray!60] (0,-1) rectangle (2,1);
  \node[below] at (1,-1.08) {$\alpha$ (responsiveness)};
  \node[rotate=90,above] at (-0.1,0) {$\beta$ (baseline)};
  % indifferent voters for rising I_q
  \draw[thick,dotted] (0,0.4) -- (2,0.4) node[right] {$\psi(\alpha;0)$};
  \draw[thick,dashed] (0,0.75) -- (2,-0.25) node[right] {$\psi(\alpha;\tfrac12)$};
  \draw[thick] (0.4,1) -- (1.4,-1) node[below] {$\psi(\alpha;2)$};
  \draw[->,thick,gray] (1.05,0.4) arc (0:-60:0.35);
  % candidates and midpoint
  \fill (0.2,-0.1) circle (0.03) node[left] {$\theta_l$};
  \fill (1.2,0.9) circle (0.03) node[right] {$\theta_r$};
  \draw[fill=white] (0.7,0.4) circle (0.03);
  \node[below left] at (0.7,0.4) {$\theta_m$};
  % switching labels
  \node at (1.65,-0.45) {$l\to r$};
  \node at (0.18,0.72) {$r\to l$};
\end{tikzpicture}
\caption{Realignment with fixed candidates. Types are $\theta_l=(0.2,-0.1)$ and $\theta_r=(1.2,0.9)$. Candidate $l$ is closer to the electorate's centre in baseline, and $r$ in responsiveness. The indifference line passes through their midpoint and has slope $-S$. Voters below it prefer $l$ (blue); voters above it prefer $r$ (red). Shading shows preferences at $S=2$, with lighter regions identifying voters who have switched since $S=0$. For a uniform electorate on the square, $l$'s support falls from 70 percent at $S=0$ to 45 percent at $S=2$.}
\label{fig:realign}
\end{figure}

Since $\Delta U_i=2\Delta\beta\big(\beta_i-\psi(\alpha_i)\big)$, the vertical distance from a voter's type to the indifference line measures the intensity of her induced preference. Consider a voter with $\beta_i=\beta_m$. She is indifferent when $S=0$, but for positive $S$ her preference has intensity
\[
|\Delta U_i|=2|\Delta\alpha|\,|\alpha_i-\alpha_m|\,S.
\]
The voter's preference can become stronger than that of a voter with an extreme baseline but responsiveness near $\alpha_m$. A former moderate can thus become a committed partisan. Empirical studies link shocks of this kind to shifts in alignment: financial crises, bank failures, import competition, and the decline of middle-class occupations have moved voters toward radical and populist parties \citep{Funke2016,Doerr2022,ColantoneStanig2018,Kurer2020,Algan2017}.

Part (b) is a first sense in which induced preferences harden. As $S$ rises, every voter becomes less willing to delegate to anyone whose responsiveness differs from hers, because the officeholder acts more on her own response. Figure~\ref{fig:delegation} shows why: the gap between the two policies is linear in the officeholder's forecast, greater unpredictability puts more weight where the gap is large, and the payoff, minus the squared gap, is concave. Parties face the same cost of delegation as voters.

\begin{figure}[tbp]
\centering
\begin{subfigure}[t]{0.49\textwidth}
\centering
\begin{tikzpicture}
\begin{axis}[width=7.0cm, height=6.6cm,
  xmin=-1.9, xmax=1.9, ymin=-3.1, ymax=2.6,
  axis lines=left, xtick={-1.2,-0.6,0,0.6,1.2}, xticklabels={$-1.2$,$-0.6$,$0$,$0.6$,$1.2$}, ytick=\empty,
  xlabel={officeholder's forecast $\mathbb{E}_\mu[\omega]$}, ylabel={policy},
  label style={font=\footnotesize}, tick label style={font=\scriptsize}, clip=false]
  % distributions of the forecast
  \addplot[gray!70!black, thick, domain=-1.9:1.9, samples=120] {-3.1+1.2*exp(-x^2/(2*0.36))/(0.6*sqrt(2*pi))};
  \addplot[black, very thick, domain=-1.9:1.9, samples=120] {-3.1+1.2*exp(-x^2/(2*1.44))/(1.2*sqrt(2*pi))};
  \node[font=\scriptsize, anchor=west] at (axis cs:-1.85,2.35) {distribution of the forecast:};
  \draw[black, very thick] (axis cs:-1.82,2.0) -- (axis cs:-1.57,2.0);
  \node[font=\scriptsize, anchor=west] at (axis cs:-1.55,2.0) {high $S$};
  \draw[gray!70!black, thick] (axis cs:-1.82,1.68) -- (axis cs:-1.57,1.68);
  \node[font=\scriptsize, anchor=west, gray!70!black] at (axis cs:-1.55,1.68) {low $S$};
  % forecasts at plus and minus sqrt(I)
  \draw[gray!70!black, densely dotted] (axis cs:-0.6,-3.1) -- (axis cs:-0.6,-0.47);
  \draw[gray!70!black, densely dotted] (axis cs:0.6,-3.1) -- (axis cs:0.6,-0.13);
  \draw[black, densely dotted] (axis cs:-1.2,-3.1) -- (axis cs:-1.2,-1.19);
  \draw[black, densely dotted] (axis cs:1.2,-3.1) -- (axis cs:1.2,-0.01);
  % policies
  \addplot[gray!70!black, thick, domain=-1.8:1.8] {-0.25+0.2*x} node[pos=1, right, font=\footnotesize] {$x_i^*(\mu)$};
  \addplot[black, very thick, domain=-1.8:1.8] {0.25+1.2*x} node[pos=1, right, font=\footnotesize] {$x_j^*(\mu)$};
  % gaps at plus and minus sqrt(I)
  \draw[gray!70!black, line width=2.2pt] (axis cs:-0.6,-0.47) -- (axis cs:-0.6,-0.37);
  \draw[gray!70!black, line width=2.2pt] (axis cs:0.6,-0.13) -- (axis cs:0.6,0.97);
  \draw[black, line width=2.2pt] (axis cs:-1.2,-1.19) -- (axis cs:-1.2,-0.49);
  \draw[black, line width=2.2pt] (axis cs:1.2,-0.01) -- (axis cs:1.2,1.69);
  \node[font=\footnotesize, anchor=west] at (axis cs:1.24,0.85) {gap};
\end{axis}
\end{tikzpicture}
\caption{Gap between the policies}
\end{subfigure}\hfill
\begin{subfigure}[t]{0.49\textwidth}
\centering
\begin{tikzpicture}
\begin{axis}[width=7.0cm, height=6.6cm,
  xmin=-1.9, xmax=1.9, ymin=-4.6, ymax=0.2,
  axis lines=left, xtick={-1.2,-0.6,0,0.6,1.2}, xticklabels={$-1.2$,$-0.6$,$0$,$0.6$,$1.2$},
  ytick={-0.25,-0.61,-1.69}, yticklabels={$-(\Delta\beta_{ij})^2$,{$-\Lambda_i^j$, low $S$},{$-\Lambda_i^j$, high $S$}},
  xlabel={officeholder's forecast $\mathbb{E}_\mu[\omega]$}, ylabel={$-\big(\Delta x^*_{ij}(\mu)\big)^2$},
  label style={font=\footnotesize}, tick label style={font=\scriptsize}, clip=false]
  % payoff at the average forecast
  \addplot[gray, dashed, domain=-1.9:1.9] {-0.25};
  % payoff from j's policy relative to i's own policy
  \addplot[black, very thick, domain=-1.5:1.6, samples=120] {-(0.5+x)^2};
  % chords: the expected payoff is the midpoint of the chord between the forecasts at plus and minus sqrt(S)
  \addplot[gray!70!black, thick] coordinates {(-0.6,-0.01) (0.6,-1.21)};
  \addplot[black, thick] coordinates {(-1.2,-0.49) (1.2,-2.89)};
  \draw[gray!70!black, densely dotted] (axis cs:-1.9,-0.61) -- (axis cs:0,-0.61);
  \draw[black, densely dotted] (axis cs:-1.9,-1.69) -- (axis cs:0,-1.69);
  \addplot[only marks, mark=*, mark size=1.5pt, gray!70!black] coordinates {(-0.6,-0.01) (0.6,-1.21)};
  \addplot[only marks, mark=*, mark size=1.5pt, black] coordinates {(-1.2,-0.49) (1.2,-2.89)};
  \addplot[only marks, mark=*, mark size=2.3pt, gray!70!black] coordinates {(0,-0.61)};
  \addplot[only marks, mark=*, mark size=2.3pt, black] coordinates {(0,-1.69)};
  \addplot[only marks, mark=*, mark size=2.3pt, mark options={fill=white, draw=black, thick}] coordinates {(0,-0.25)};
\end{axis}
\end{tikzpicture}
\caption{Payoff from delegating}
\end{subfigure}
\caption{Why the cost of delegation rises with unpredictability. Panel (a): given the officeholder's forecast $\mathbb{E}_\mu[\omega]$, officeholder $j$ chooses $x_j^*(\mu)=\beta_j+\alpha_j\mathbb{E}_\mu[\omega]$, and player $i$ would choose $x_i^*(\mu)=\beta_i+\alpha_i\mathbb{E}_\mu[\omega]$. The gap between the two policies, $\Delta x^*_{ji}(\mu)=\Delta\beta_{ji}+\Delta\alpha_{ji}\mathbb{E}_\mu[\omega]$, is linear in the forecast. A higher $S$ spreads the forecast (bottom) and moves the forecasts at $\pm\sqrt{S}$ to where the gap is larger (bars), while the average gap stays at $\Delta\beta_{ji}$. Panel (b): relative to governing herself, player $i$'s payoff from delegating to $j$ is minus the expected squared gap, $U_i^j-U_i^i=-\Lambda_i^j$. Because this payoff is concave in the forecast, its expected value, the midpoint of the chord between the forecasts at $\pm\sqrt{S}$, lies below its value at the average forecast, $-(\Delta\beta_{ij})^2$ (open dot), by $(\Delta\alpha_{ij})^2S$. Parameters: $\Delta\beta_{ji}=0.5$ and $\Delta\alpha_{ji}=1$; $\sqrt{S}=0.6$ (gray) and $\sqrt{S}=1.2$ (black), so $-\Lambda_i^j$ falls from $-0.61$ to $-1.69$.}
\label{fig:delegation}
\end{figure}

\begin{remark}[Measurement]\label{rem:measurement}
A question about the policy that a voter wants today elicits $\beta_i+\alpha_i\mathbb{E}[\omega]=\beta_i$ and reveals nothing about $\alpha_i$. Measures of the distance between a voter and a candidate therefore compare expected policies, $\big(\mathbb{E}[x^*_j(\mu)]-\beta_i\big)^2=(\Delta\beta_{ij})^2$. Induced preferences depend on the expected distance $\Lambda_i^j=(\Delta\beta_{ij})^2+(\Delta\alpha_{ij})^2S$. When $S$ is large, a voter can prefer the candidate whose expected policy is farther from her own preferred policy, because that candidate's response is closer to hers.
\end{remark}

Because each voter's comparison is affine in $S$, a rise in uncertainty and an improvement in information cannot move the comparison in opposite directions, and a fall in $S$ reverses a switch. The restriction concerns comparisons across distributions of posteriors. Along a sequence of realized news, posterior means can move in either direction.

\paragraph{The electoral margin.}
For every slope, there is a line that divides the electorate into two equal groups. Such a line is called a median line \citep{DavisDeGrootHinich1972,McKelvey1986}. A median line is unique for a given slope because the density is bounded away from zero on a convex support.

\begin{lemma}\label{lem:median}
Fix candidates with $\beta_r>\beta_l$. There is a unique $M(\theta_l,\theta_r)$ such that
\[
G\big(\{\theta_i:\beta_i\le\psi_{M(\theta_l,\theta_r)}(\alpha_i)\}\big)=\tfrac12 ,
\]
and candidate $l$ wins if and only if $\eta<M(\theta_l,\theta_r)$.
\end{lemma}

Increasing $\eta$ lowers the indifference line and reduces $l$'s vote share, so $l$ wins precisely when the noise is below $M$. Since $\psi_M=\psi+M/(2\Delta\beta)$, the ratio $M/(2|\Delta\beta|)$ measures the vertical distance between the noiseless indifference line and the parallel median line. A positive $M$ means that more than half the electorate prefers $l$ without noise.

\begin{figure}[tbp]
\centering
\begin{minipage}[b]{0.49\linewidth}
\centering
\begin{tikzpicture}[scale=2.8, every node/.style={font=\small}]
  % (a) I_q = 1/2: l wins
  \fill[blue!22] (0,-1) -- (2,-1) -- (2,-0.25) -- (0,0.75) -- cycle;
  \fill[red!22]  (0,0.75) -- (2,-0.25) -- (2,1) -- (0,1) -- cycle;
  \fill[blue!60] (0,0.5) -- (0,0.75) -- (2,-0.25) -- (2,-0.5) -- cycle;
  \draw[gray!60] (0,-1) rectangle (2,1);
  \node[above] at (1,1.02) {(a) $S=\tfrac12$: $l$ wins};
  \node[below] at (1,-1.04) {$\alpha$ (responsiveness)};
  \node[rotate=90,above] at (-0.06,-0.5) {$\beta$ (baseline)};
  \draw[thick] (0,0.75) -- (2,-0.25);
  \node[above] at (1.7,-0.1) {$\psi(\alpha;\tfrac12)$};
  \draw[thick,dashed] (0,0.5) -- (2,-0.5);
  \node[below] at (1.6,-0.5) {$\psi_M(\alpha;\tfrac12)$};
  \draw[decorate,decoration={brace,amplitude=4pt}] (-0.02,0.5) -- (-0.02,0.75)
    node[midway,left=5pt] {$\dfrac{M}{2|\Delta\beta|}$};
  \fill (0.2,-0.1) circle (0.03) node[left] {$\theta_l$};
  \fill (1.2,0.9) circle (0.03) node[right] {$\theta_r$};
  \draw[fill=white] (0.7,0.4) circle (0.03);
  \node[below left] at (0.7,0.4) {$\theta_m$};
  \fill (1,0) ++(-0.025,-0.025) rectangle ++(0.05,0.05);
  \node[below left] at (0.97,-0.02) {$\theta_M$};
\end{tikzpicture}
\end{minipage}\hfill
\begin{minipage}[b]{0.49\linewidth}
\centering
\begin{tikzpicture}[scale=2.8, every node/.style={font=\small}]
  % (b) I_q = 2: r wins
  \fill[blue!22] (0,-1) -- (1.4,-1) -- (0.4,1) -- (0,1) -- cycle;
  \fill[red!22]  (0.4,1) -- (2,1) -- (2,-1) -- (1.4,-1) -- cycle;
  \fill[red!60] (0.4,1) -- (0.5,1) -- (1.5,-1) -- (1.4,-1) -- cycle;
  \draw[gray!60] (0,-1) rectangle (2,1);
  \node[above] at (1,1.02) {(b) $S=2$: $r$ wins};
  \node[below] at (1,-1.04) {$\alpha$ (responsiveness)};
  \node[rotate=90,above] at (-0.06,-0.5) {$\beta$ (baseline)};
  \draw[thick] (0.4,1) -- (1.4,-1);
  \node[left] at (1.33,-0.85) {$\psi(\alpha;2)$};
  \draw[thick,dashed] (0.5,1) -- (1.5,-1);
  \node[right] at (1.5,-0.85) {$\psi_M(\alpha;2)$};
  \draw[decorate,decoration={brace,amplitude=4pt}] (0.6,0.6) -- (0.6,0.8);
  \node[left] (Mlab) at (-0.02,0.7) {$\dfrac{|M|}{2|\Delta\beta|}$};
  \draw[->,thin] (Mlab.east) -- (0.555,0.7);
  \fill (0.2,-0.1) circle (0.03) node[left] {$\theta_l$};
  \fill (1.2,0.9) circle (0.03) node[right] {$\theta_r$};
  \draw[fill=white] (0.7,0.4) circle (0.03);
  \node[below left] at (0.7,0.4) {$\theta_m$};
  \fill (1,0) ++(-0.025,-0.025) rectangle ++(0.05,0.05);
  \node[right] at (1.04,0) {$\theta_M$};
\end{tikzpicture}
\end{minipage}
\caption{The electoral margin at two levels of unpredictability. Candidates are as in Figure~\ref{fig:realign}, and voters are uniform on the square. The solid line contains voters who are indifferent without noise. The dashed line is parallel to it, divides the electorate equally, and passes through $\theta_M$. Light shading marks each candidate's voters, blue for $l$ and red for $r$. The dark strip between the two lines is the winning margin. In panel (a), the dark blue strip gives $l$ a majority and $M>0$. In panel (b), the dark red strip gives $r$ a majority and $M<0$.}
\label{fig:median}
\end{figure}

\paragraph{A median point.}
Some results use the following restriction on the voter distribution.

\begin{assumption}[Median in all directions]\label{ass:median}
There is a point $\theta_M=(\alpha_M,\beta_M)$ such that every line through $\theta_M$ splits $G$ in half.
\end{assumption}

Under this assumption, all median lines pass through a common point. The coordinates of the common point are the marginal medians, and the point remains decisive as the direction of candidate disagreement changes. Central symmetry about $\theta_M$ is sufficient \citep{DavisDeGrootHinich1972}. Examples include uniform distributions on squares, parallelograms, ellipses, and rhombi, so the assumption allows linear correlation between baseline and responsiveness, as well as a centre that differs most in responsiveness (Figure~\ref{fig:median-types}(c)).

Some skewed distributions also have a median point, and many do not: for a uniform distribution on a triangle, median lines of different slopes do not meet at a common point (Figure~\ref{fig:median-types}(d)). The SI gives results without Assumption~\ref{ass:median}.

\begin{figure}[t]
\centering
\begin{tikzpicture}[scale=1.45, every node/.style={font=\footnotesize}]
  % (a) centrally symmetric: uniform on the square
  \begin{scope}
    \fill[gray!22] (0,-1) rectangle (2,1);
    \draw[gray!60] (0,-1) rectangle (2,1);
    \draw[semithick] (0,0) -- (2,0);
    \draw[semithick,dashed] (1,-1) -- (1,1);
    \draw[semithick,dotted] (0,1) -- (2,-1);
    \fill (1,0) ++(-0.045,-0.045) rectangle ++(0.09,0.09);
    \node[above right] at (1.02,0.02) {$\theta_M$};
    \node[above] at (1,1.04) {(a) Independent};
    \node[below] at (1,-1.04) {$\alpha$};
    \node[left] at (-0.04,0) {$\beta$};
  \end{scope}
  % (b) linear correlation: uniform on a tilted ellipse
  \begin{scope}[xshift=2.6cm]
    \fill[gray!22, rotate around={15:(1,0)}] (1,0) ellipse (0.95 and 0.28);
    \draw[gray!60] (0,-1) rectangle (2,1);
    \draw[semithick] (0,0) -- (2,0);
    \draw[semithick,dashed] (1,-1) -- (1,1);
    \draw[semithick,dotted] (0,1) -- (2,-1);
    \fill (1,0) ++(-0.045,-0.045) rectangle ++(0.09,0.09);
    \node[above right] at (1.02,0.02) {$\theta_M$};
    \node[above] at (1,1.04) {(b) Correlated};
    \node[below] at (1,-1.04) {$\alpha$};
  \end{scope}
  % (c) strong responsiveness concentrated at middle baselines: uniform on a rhombus
  \begin{scope}[xshift=5.2cm]
    \fill[gray!22] (0,0) -- (1,1) -- (2,0) -- (1,-1) -- cycle;
    \draw[gray!60] (0,-1) rectangle (2,1);
    \draw[semithick] (0,0) -- (2,0);
    \draw[semithick,dashed] (1,-1) -- (1,1);
    \draw[semithick,dotted] (0,1) -- (2,-1);
    \fill (1,0) ++(-0.045,-0.045) rectangle ++(0.09,0.09);
    \node[above right] at (1.02,0.02) {$\theta_M$};
    \node[above] at (1,1.04) {(c) Responsive centre};
    \node[below] at (1,-1.04) {$\alpha$};
  \end{scope}
  % (d) skewed: uniform on a triangle, no median point
  \begin{scope}[xshift=7.8cm]
    \fill[gray!22] (0,-1) -- (2,-1) -- (0,1) -- cycle;
    \draw[gray!60] (0,-1) rectangle (2,1);
    \fill[black!45] (0.586,-0.414) -- (0.828,-0.414) -- (0.586,-0.172) -- cycle;
    \draw[semithick] (0,-0.414) -- (2,-0.414);
    \draw[semithick,dashed] (0.586,-1) -- (0.586,1);
    \draw[semithick,dotted] (0,0.414) -- (1.414,-1);
    \draw[fill=white] (0.586,-0.414) ++(-0.045,-0.045) rectangle ++(0.09,0.09);
    \node[above] at (1,1.04) {(d) No median point};
    \node[below] at (1,-1.04) {$\alpha$};
  \end{scope}
\end{tikzpicture}
\caption{Voter distributions with and without a median point. In each panel, voters are uniform on the shaded region, and the solid, dashed, and dotted lines each divide it equally. Panels (a)--(c) are centrally symmetric, so every median line passes through $\theta_M$. They differ in how baseline and responsiveness are related: independent in (a); correlated in (b), where baselines are more concentrated than responsiveness, so the centre is thick when baselines divide voters and thin when responsiveness does; and uncorrelated in (c), where responsiveness is most dispersed among voters with baselines near the middle. In (d), the support is a triangle, and the three median lines enclose the dark region; the white square marks the marginal medians.}
\label{fig:median-types}
\end{figure}

\begin{corollary}[Median voter]\label{cor:median}
Under Assumption~\ref{ass:median},
\begin{align*}
M(\theta_l,\theta_r)
&=2\Delta\beta\big(\beta_M-\psi(\alpha_M)\big)=\Delta U_{\theta_M}\\
&=2\Delta\beta\,(\beta_M-\beta_m)+2\Delta\alpha\,S\,(\alpha_M-\alpha_m),
\end{align*}
so candidate $l$ wins if and only if the noise is below the payoff gap of a voter at the median point.
\end{corollary}

Figure~\ref{fig:median} uses a uniform distribution on the square. Candidate $l$ leads without noise exactly when the median point lies below the indifference line. At low unpredictability, $l$ leads because she is closer in baseline. As $S$ rises, the line rotates and crosses the median point at $S=4/3$. Candidate $r$ then leads because she is closer in responsiveness. Without a median point, the parallel halving line can also shift as the slope changes, which adds a second source of movement in $M$.

\paragraph{Winning probability.}
\begin{assumption}[Small $\phi$]\label{ass:small-h}
The noise is dispersed enough: $\phi\le\overline{\phi}$, where $\overline{\phi}\overline{M}<\tfrac12$ and $\overline{M}\equiv\sup|M(\theta_l,\theta_r)|$ over candidates in $\operatorname{supp}(G)$ and distributions of posteriors $F_\mu$. The value of the noise at which the vote is tied then lies strictly inside $\big[-\tfrac{1}{2\phi},\tfrac{1}{2\phi}\big]$ for every pair of candidates.
\end{assumption}

Under Assumption~\ref{ass:small-h}, $\pi=\tfrac12+\phi M$ (Lemma~\ref{lem:win}), so $\phi$ measures how strongly a party's winning probability responds to the electoral margin. I refer to $\phi$ as electoral sensitivity. A small $\phi$ means that electoral noise is diffuse, so winning probabilities respond only moderately to nominations. Results stated for moderate electoral sensitivity hold for $\phi$ below a further positive bound (Lemma~\ref{lem:smallh}); results stated for every $\phi$ do not use Assumption~\ref{ass:small-h}.

\begin{lemma}[Winning probability]\label{lem:win}
Under Assumption~\ref{ass:small-h}, for any candidates,
\[
\pi=\Pr(l\text{ wins})=\tfrac12+\phi\,M(\theta_l,\theta_r).
\]
The parties' expected payoffs are
\begin{align*}
\Pi_L&=\tfrac12\big(U_L^l+U_L^r+B\big)+\phi M\,V_L,\\
\Pi_R&=\tfrac12\big(U_R^l+U_R^r+B\big)-\phi M\,V_R.
\end{align*}
\end{lemma}

Once $S$ is large, the noise can decide only few votes, which is a second sense in which induced preferences harden.

\begin{remark}[Hardening]\label{rem:hardening}
Fix candidates with $\beta_r>\beta_l$ and $\Delta\alpha\ne0$.
\begin{enumerate}
  \item[(a)] If a voter's baseline and responsiveness favour the same candidate, that is, if $\Delta\beta(\beta_i-\beta_m)$ and $\Delta\alpha(\alpha_i-\alpha_m)$ have the same sign, her attachment $|\Delta U_i|$ rises with $S$. If they favour different candidates, her attachment first falls, she switches at her threshold, and her attachment then rises.
  \item[(b)] The share of voters whose vote depends on the realization of the noise, $G\big(\{\theta_i:|\Delta U_i|\le\tfrac{1}{2\phi}\}\big)$, and the share of voters whom a common shift $v$ in every payoff gap moves are each at most a constant times $1/S$, because $\|\nabla\Delta U\|=2\big((\Delta\beta)^2+(\Delta\alpha)^2S^2\big)^{1/2}$.
  \item[(c)] Under Assumption~\ref{ass:small-h}, a small common shift $v$ in every payoff gap changes $l$'s winning probability by $\phi v$, whatever the value of $S$.
  \item[(d)] The variance of payoff gaps across voters is quadratic in $S$, with leading coefficient $4(\Delta\alpha)^2\operatorname{Var}_G(\alpha_i)$, so it rises with $S$ once $S$ is large enough.
\end{enumerate}
\end{remark}

Parts (b) and (c) concern events that are orthogonal to future policy and common to all voters, such as a scandal about a candidate's conduct; I write $v$ for such a valence shift. An event that reveals how a candidate would respond to the state changes her responsiveness instead, and its weight in voters' comparisons rises with $S$. Voters whose baseline and responsiveness favour the same candidate become less willing to accept defeat. Voters for whom the two favour different candidates are the ones who switch, as former Obama voters and Labour voters in the Red Wall did. Once uncertainty is high, the noise and ordinary events, such as a scandal or a disappointing economy, can decide only a small share of votes. Part (b) is a bound: the share of voters that a scandal moves need not fall at every increase in $S$, and it can rise over an interval. Part (c) gives a further limit: while winning probabilities are interior, a shift that moves few voters still changes the winning probability by $\phi v$. The remark also holds the candidates fixed and does not describe hardening after nominations adjust. With these qualifications, the bound is consistent with the observation that radical-right candidates such as Trump have survived scandals that would have ended a candidacy in calmer times (Section~\ref{sec:evidence}), and a scandal can still decide a close election.

\begin{remark}[Responses versus conduct]\label{rem:responses}
Fix candidates with $\beta_r>\beta_l$ and $\Delta\alpha\ne0$ whose midpoint $\theta_m$ lies in the interior of $\operatorname{supp}(G)$. A small change $\delta$ in candidate $l$'s responsiveness moves voter $i$'s comparison by $2(\alpha_i-\alpha_l)S\delta$ to first order. The share of voters whom it moves is of order $\delta$, with a coefficient that converges to a positive limit as $S$ grows, whereas the share moved by a valence shift is at most a constant times $1/S$ (Remark~\ref{rem:hardening}(b)). Under Assumptions~\ref{ass:median} and~\ref{ass:small-h}, $\partial\pi/\partial\alpha_l=2\phi(\alpha_M-\alpha_l)S$, while a small valence shift $v$ changes $\pi$ by $\phi v$. The weight of news about a candidate's response, relative to news about her conduct, therefore rises linearly in $S$.
\end{remark}

A campaign event that reveals how a candidate would respond to the state, such as her reaction to an early crisis, therefore moves more votes as uncertainty rises, while a scandal about her conduct moves fewer.

\section{Nominations}\label{sec:competition}

In Sections~\ref{sec:policy} and~\ref{sec:realign}, the candidates were fixed. With the candidates fixed, greater uncertainty and better information widen the expected policy gap, raise the value of office, and realign voters. Parties now choose their candidates in response. When $S=0$, responsiveness affects no payoff, and each party is assumed to nominate a candidate with its own responsiveness.

\paragraph{Equilibrium nominations.}
Party $L$ chooses $\theta_l$ to maximize
\[
\Pi_L(\theta_l,\theta_r)=\Big(\tfrac12+\phi M\Big)\big(U_L^l+B\big)+\Big(\tfrac12-\phi M\Big)U_L^r.
\]
A new nominee changes the policy that party $L$ obtains if $l$ wins, and changes the probability that $l$ wins. The first effect is weighted by $\tfrac12+\phi M$. The second is worth $V_L$ per unit increase in the winning probability. Party $R$ faces the corresponding trade-off.

\begin{proposition}[Equilibrium nominations]\label{prop:nomination}
Suppose Assumption~\ref{ass:small-h} holds and $S>0$, and let $\overline{\theta}=(\overline{\alpha},\overline{\beta})$ be the average type of the voters on the median line $\psi_M$, weighted by $g$. In any equilibrium in which both candidates lie in the interior of $\operatorname{supp}(G)$, both total values of office are positive, and the tie value $M$ is differentiable in the candidates' types,
\[
\theta_l=w_L\,\theta_L+(1-w_L)\,\overline{\theta},
\qquad
\theta_r=w_R\,\theta_R+(1-w_R)\,\overline{\theta},
\]
where
\[
w_L=\frac{\tfrac12+\phi M}{\tfrac12+\phi M+\phi V_L},
\qquad
w_R=\frac{\tfrac12-\phi M}{\tfrac12-\phi M+\phi V_R}
\]
are evaluated at the equilibrium candidates and lie in $(0,1)$. Under Assumption~\ref{ass:median}, $\overline{\theta}=\theta_M$.
\end{proposition}

A party's own residual term, $U_J^J=-\alpha_J^2\big(\operatorname{Var}(\omega)-S\big)$, does not depend on the nominees, so nominations, like the comparisons in Section~\ref{sec:realign}, depend on the information environment only through $S$. Each candidate lies between her party's ideal type and $\overline{\theta}$, with the same share $w_J$ of the party's ideal type kept in both dimensions. The party chooses the moderation $1-w_J$ to balance the electoral benefit of moderation against its policy cost. The electoral benefit rises with the party's total value of office, and the policy cost arises only if the party's candidate wins. The average type $\overline{\theta}$ is therefore the decisive point for nominations. In general, $\overline{\theta}$ changes with the slope of the median line, and hence with the candidates and $S$. Under Assumption~\ref{ass:median}, $\overline{\theta}$ stays at $\theta_M$.


\section{The effects of unpredictability}\label{sec:effects}

In this section, I combine the two effects. With the candidates fixed, greater unpredictability widens the expected policy gap (Proposition~\ref{prop:auto-polar}); through nominations, greater unpredictability moves the candidates toward the median voter (Proposition~\ref{prop:nomination}). In equilibrium, the widening of the policy gap dominates the moderation. I also compare greater unpredictability with a higher office benefit $B$, which makes office more valuable in a different way. The symmetric benchmark gives exact comparative statics for both.

\paragraph{A symmetric benchmark.}

\begin{proposition}[Symmetric benchmark]\label{prop:symmetric}
Suppose Assumption~\ref{ass:median} holds, $S>0$, and the parties' distinct ideal types are symmetric about the median point, $\theta_L+\theta_R=2\theta_M$. For every $\phi>0$ and common office benefit $B\ge0$, the equilibrium is unique and symmetric. Each candidate wins with probability $\tfrac12$, and
\[
\theta_j=w\theta_J+(1-w)\theta_M,
\qquad
w=\frac{2}{1+2\phi B+\sqrt{(1+2\phi B)^2+8\phi\Lambda}},
\]
where $\Lambda\equiv\Lambda_L^R=(\beta_R-\beta_L)^2+(\alpha_R-\alpha_L)^2S$. Parts (a)--(c) assume $\alpha_R\ne\alpha_L$ and hold the benefit fixed.
\begin{enumerate}
  \item[(a)] Holding $B$ fixed, a higher $S$ brings the candidates closer together in both dimensions whenever $\alpha_R\ne\alpha_L$:
  \[
  \theta_r-\theta_l=w(\theta_R-\theta_L),
  \qquad
  \frac{\partial w}{\partial S}=-\frac{2\phi w^3}{1+2\phi\Lambda w^2}(\alpha_R-\alpha_L)^2<0.
  \]
  \item[(b)] The expected squared policy gap nevertheless rises:
  \[
  D^*=w^2\Lambda,
  \qquad
  \frac{\partial D^*}{\partial S}=\frac{(1+2\phi B)\,w^3}{1+2\phi\Lambda w^2}(\alpha_R-\alpha_L)^2>0.
  \]
  The gap $D^*$ is below $1/(2\phi)$ and approaches that ceiling as $\Lambda\to\infty$ at fixed $B$.
  \item[(c)] Each party's policy value of office is $w\Lambda$, and its total value is $B+w\Lambda$. Both rise with $S$ at rate $w(\alpha_R-\alpha_L)^2/(1+2\phi\Lambda w^2)$.
  \item[(d)] Holding the information environment fixed, a higher common office benefit brings candidates and policies closer together:
  \[
  \frac{\partial w}{\partial B}=-\frac{2\phi w^2}{1+2\phi\Lambda w^2}<0,
  \qquad
  \frac{\partial D^*}{\partial B}=-\frac{4\phi\Lambda w^3}{1+2\phi\Lambda w^2}<0.
  \]
  The policy value of office falls, while the total value rises, with $\partial(B+w\Lambda)/\partial B=1/(1+2\phi\Lambda w^2)>0$.
  \item[(e)] Suppose that the parties place weight $P_J>0$ on policy and the voters weight $P_V>0$. Then
  \[
  w=\frac{2}{1+2\phi P_VB/P_J+\sqrt{(1+2\phi P_VB/P_J)^2+8\phi P_V\Lambda}}.
  \]
  A higher $P_J$ raises $w$ and $D^*$ when $B>0$, and a higher $P_V$ lowers both. A common weight $P=P_J=P_V$ also lowers both:
  \[
  \frac{\partial w}{\partial P}=-\frac{2\phi\Lambda w^3}{1+2\phi P\Lambda w^2}<0,
  \qquad
  \frac{\partial D^*}{\partial P}=-\frac{4\phi\Lambda^2w^4}{1+2\phi P\Lambda w^2}<0.
  \]
\end{enumerate}
\end{proposition}

A higher $B$ makes victory more valuable at a given $\Lambda$, so the smaller difference in the nominees' types also means a smaller expected policy gap. Information raises the policy cost of losing and also increases $\Lambda$, the disagreement between the parties' ideal policy rules. The direct increase in $\Lambda$ outweighs the equilibrium moderation. A party moderates only until the electoral gain from further moderation equals the policy cost of moderation, so the moderation offsets only part of the increase in $\Lambda$.

The nomination condition is
\[
1-w=2\phi w(B+w\Lambda),
\qquad\text{so}\qquad
D^*=\frac{1-w}{2\phi}-Bw.
\]
When $B=0$, the equation reduces to $D^*=(1-w)/(2\phi)$: as parties moderate more in response to information, the policy gap rises in proportion to the moderation $1-w$. The implication of more valuable office for polarization therefore depends on why winning has become more valuable.

Part (e) checks the mechanism against a simpler source of higher stakes, and I state it only for the symmetric benchmark. If only the parties come to care more about policy, a higher $P_J$ works like a lower benefit, and policies diverge, as in \citet{Wittman1983} and \citet{Calvert1985}. If voters come to care more about policy, a higher $P_V$ makes votes respond more to policy, and parties gain more from moderation. Policies then converge, also when the two weights rise together. Greater uncertainty and better information also raise voters' stakes, and policies still diverge. A common weight revalues the same disagreement $\Lambda$. Uncertainty and information raise $\Lambda$ itself. Propositions~\ref{prop:auto-polar}--\ref{prop:symmetric} hold in general with several states and several policies, including correlated states, with unpredictability as a matrix (Proposition~\ref{prop:multistate}).

\begin{table}[H]
\centering
\small
\renewcommand{\arraystretch}{1.25}
\begin{tabular}{@{}lccc@{}}
\toprule
{\bfseries\boldmath Change} & {\bfseries\boldmath Type gap} & {\bfseries\boldmath Policy gap} & {\bfseries\boldmath Realignment} \\
\hline\hline
$B\uparrow$ & $\downarrow$ & $\downarrow$ & No \\ \midrule
$P_J\uparrow$ (parties) & $\uparrow$ & $\uparrow$ & No \\ \midrule
$P_V\uparrow$ (voters) & $\downarrow$ & $\downarrow$ & No \\ \midrule
$P_J=P_V\uparrow$ (all players) & $\downarrow$ & $\downarrow$ & No \\ \midrule
$S\uparrow$ & $\downarrow$ & $\uparrow$ & Yes \\
\bottomrule
\end{tabular}
\caption{Comparative statics in the symmetric benchmark. The weights $P_J$ and $P_V$ multiply the parties' and the voters' policy payoffs (Proposition~\ref{prop:symmetric}(e)); the row for $P_J$ requires $B>0$, and the last row requires $\alpha_R\ne\alpha_L$. Greater uncertainty and a higher revealed share have the effects of a higher $S$ (Section~\ref{sec:realign}). ``Policy gap'' is the expected squared difference between policies.}
\label{tab:benefit-comparison}
\end{table}

Figure~\ref{fig:symmetric} illustrates the information comparison at $B=0$, and Figure~\ref{fig:twoeffects} decomposes the change in polarization into the widening with fixed candidates and the moderation that follows.

\begin{figure}[tbp]
\centering
\begin{minipage}[b]{0.49\linewidth}
\centering
\begin{tikzpicture}[scale=2.8, every node/.style={font=\small}]
  \fill[blue!22] (0,-1) -- (1.1944,-1) -- (1,0) -- (0,0.4286) -- cycle;
  \fill[blue!8]  (0,0.4286) -- (1,0) -- (0.8056,1) -- (0,1) -- cycle;
  \fill[red!22]  (1,0) -- (2,-0.4286) -- (2,1) -- (0.8056,1) -- cycle;
  \fill[red!8]   (1,0) -- (1.1944,-1) -- (2,-1) -- (2,-0.4286) -- cycle;
  \draw[gray!60] (0,-1) rectangle (2,1);
  \node[above] at (1,1.02) {(a) Equilibrium candidates};
  \draw[thick,dotted] (0,0.4286) -- (2,-0.4286);
  \node[above] at (1.72,-0.24) {$\psi(\alpha;\tfrac12)$};
  \draw[thick,dashdotted] (0.4167,1) -- (1.5833,-1);
  \node[left] at (0.50,0.86) {$\psi(\alpha;2)$};
  \draw[thick] (0.8056,1) -- (1.1944,-1);
  \node[left] at (1.16,-0.85) {$\psi(\alpha;6)$};
  \node at (1.78,-0.62) {$l\to r$};
  \node at (0.35,0.40) {$r\to l$};
  \node[below] at (1,-1.04) {$\alpha$ (responsiveness)};
  \node[rotate=90,above] at (-0.06,0) {$\beta$ (baseline)};
  \draw[gray!50,dotted] (0.4,-0.7) -- (1.6,0.7);
  \draw[gray,thick,->] (0.6044,-0.4615) -- (0.7819,-0.2545);
  \draw[gray,thick,->] (1.3956,0.4615) -- (1.2181,0.2545);
  \fill[gray!35] (0.6363,-0.4243) circle (0.03);
  \fill[gray!65] (0.6967,-0.3538) circle (0.03);
  \fill (0.7709,-0.2673) circle (0.03) node[below right] {$\theta_l$};
  \fill[gray!35] (1.3637,0.4243) circle (0.03);
  \fill[gray!65] (1.3033,0.3538) circle (0.03);
  \fill (1.2291,0.2673) circle (0.03) node[above left] {$\theta_r$};
  \draw[thick] (0.4,-0.7) +(0,0.05) -- +(0.05,0) -- +(0,-0.05) -- +(-0.05,0) -- cycle;
  \node[left] at (0.36,-0.7) {$\theta_L$};
  \draw[thick] (1.6,0.7) +(0,0.05) -- +(0.05,0) -- +(0,-0.05) -- +(-0.05,0) -- cycle;
  \node[right] at (1.64,0.7) {$\theta_R$};
  \fill (1,0) ++(-0.025,-0.025) rectangle ++(0.05,0.05);
  \node[above right] at (1.02,0.02) {$\theta_M$};
\end{tikzpicture}
\end{minipage}\hfill
\begin{minipage}[b]{0.49\linewidth}
\centering
\begin{tikzpicture}
\begin{axis}[
  width=7.2cm, height=6.4cm,
  xmin=0, xmax=7, ymin=0, ymax=1,
  axis lines=left, every axis plot/.append style={thick},
  xlabel={$S$}, xtick={0,2,4,6}, ytick={0,0.5,1},
  title={(b) Type gap and policy gap},
  title style={font=\small}, label style={font=\small}, tick label style={font=\small}, clip=false]
  \addplot[black, domain=0:7, samples=80] {2/(1+sqrt(1+1.6*(1.96+1.44*x)))};
  \addplot[black, dashed, domain=0:7, samples=80] {1-2/(1+sqrt(1+1.6*(1.96+1.44*x)))};
  \node[font=\small, anchor=west] at (axis cs:3.2,0.30) {type gap, $w$};
  \node[font=\small, anchor=west] at (axis cs:3.2,0.76) {policy gap, $2\phi D^*=1-w$};
  \addplot[only marks, mark=*, mark size=1.8pt, gray!35] coordinates {(0.5,0.6061) (0.5,0.3939)};
  \addplot[only marks, mark=*, mark size=1.8pt, gray!65] coordinates {(2,0.5054) (2,0.4946)};
  \addplot[only marks, mark=*, mark size=1.8pt, black] coordinates {(6,0.3818) (6,0.6182)};
\end{axis}
\end{tikzpicture}
\end{minipage}
\caption{Equilibrium in the symmetric benchmark with $B=0$. Party ideals (diamonds) are $\theta_L=(0.4,-0.7)$ and $\theta_R=(1.6,0.7)$, symmetric about $\theta_M=(1,0)$, and $\phi=0.2$. In panel (a), light gray, gray, and black dots mark nominees at $S=\tfrac12$, $2$, and $6$. Arrows show their movement toward the median point. All indifference lines pass through that point, so the election remains tied. Shading shows preferences at $S=6$; lighter regions identify voters who switched after $S=\tfrac12$. Panel (b) shows the remaining fraction of the type gap, $w$ (solid), and the policy gap as a fraction of its ceiling, $2\phi D^*=1-w$ (dashed). Greater uncertainty at a fixed positive revealed share has the same effect.}
\label{fig:symmetric}
\end{figure}

\begin{figure}[tbp]
\centering
\begin{tikzpicture}[x=1.7cm, y=0.95cm, every node/.style={font=\footnotesize},
  lab/.style={below, align=center, text width=2.1cm}]
  \draw (-0.3,0) -- (5.3,0);
  \draw[->] (-0.3,0) -- (-0.3,4.5) node[above] {polarization $D$};
  % equilibrium at S = 1/2
  \fill[gray!45] (0,0) rectangle (0.8,0.98);
  \node[above] at (0.4,0.98) {$0.98$};
  % widening with the nominees fixed
  \draw[gray, dashed] (0.8,0.98) -- (1.4,0.98);
  \fill[black!70] (1.4,0.98) rectangle (2.2,3.89);
  \node[above] at (1.8,3.89) {$+2.91$};
  % moderation by the parties
  \draw[gray, dashed] (2.2,3.89) -- (2.8,3.89);
  \draw[thick] (2.8,1.55) rectangle (3.6,3.89);
  \node[below] at (3.2,1.5) {$-2.34$};
  % equilibrium at S = 6
  \draw[gray, dashed] (3.6,1.55) -- (4.2,1.55);
  \fill[black] (4.2,0) rectangle (5.0,1.55);
  \node[above] at (4.6,1.55) {$1.55$};
  % labels
  \node[lab] at (0.4,0) {equilibrium at $S=\tfrac12$};
  \node[lab] at (1.8,0) {widening, nominees fixed};
  \node[lab] at (3.2,0) {moderation by the parties};
  \node[lab] at (4.6,0) {equilibrium at $S=6$};
\end{tikzpicture}
\caption{Two effects of greater unpredictability on polarization in the symmetric benchmark. Starting from the equilibrium at $S=\tfrac12$, raising $S$ to $6$ with the nominees held fixed widens polarization from $0.98$ to $3.89$ (Proposition~\ref{prop:auto-polar}). The parties then move their nominees toward $\theta_M$ (Proposition~\ref{prop:symmetric}(a)), which lowers polarization to its new equilibrium value, $1.55$. The moderation offsets only part of the widening (Proposition~\ref{prop:symmetric}(b)). Parameters as in Figure~\ref{fig:symmetric}.}
\label{fig:twoeffects}
\end{figure}

\paragraph{Asymmetric parties.}
For each party, define its disagreement with the median voter as
\[
\Lambda_M^J\equiv(\beta_J-\beta_M)^2+(\alpha_J-\alpha_M)^2\,S.
\]

\begin{proposition}[Asymmetric parties]\label{prop:asymmetric}
Suppose Assumptions~\ref{ass:median} and~\ref{ass:small-h} hold, $S>0$, and the parties' ideal types are distinct, lie in the interior of $\operatorname{supp}(G)$, and satisfy $\alpha_R\ne\alpha_L$. Then for all sufficiently small $\phi>0$, in the equilibrium near the parties' ideal types, which is unique (Lemma~\ref{lem:smallh}):
\begin{enumerate}
  \item[(a)] candidate $l$ wins with probability above $\tfrac12$ if and only if $\Lambda_M^L<\Lambda_M^R$: the party whose ideal type is closer to the median point leads;
  \item[(b)] as $S$ rises, $\Lambda_M^R-\Lambda_M^L$ changes at rate $(\alpha_R-\alpha_M)^2-(\alpha_L-\alpha_M)^2$, so the lead moves toward the party whose ideal responsiveness is closer to $\alpha_M$; if $L$ is closer to the median point in baseline and $R$ in responsiveness, $L$ leads for $S<S^\dagger$ and $R$ for $S>S^\dagger$, so the lead switches if $S^\dagger$ lies in the attainable range of unpredictability, where
  \[
  S^\dagger\equiv\frac{(\beta_R-\beta_M)^2-(\beta_L-\beta_M)^2}{(\alpha_L-\alpha_M)^2-(\alpha_R-\alpha_M)^2};
  \]
  \item[(c)] as unpredictability rises, both candidates move toward $\theta_M$ along the segments from their parties' ideal types, and if the ideal types lie on opposite sides of the median point, $(\theta_L-\theta_M)\cdot(\theta_R-\theta_M)\le0$ in the Euclidean inner product, the distance between the candidates' types falls;
  \item[(d)] as unpredictability rises, policies diverge: the equilibrium expected squared policy gap rises, and so do both parties' values of office.
\end{enumerate}
\end{proposition}

Figure~\ref{fig:equilibrium} sets $B=0$ and uses the earlier candidates as parties' ideal types: $\theta_L=(0.2,-0.1)$ and $\theta_R=(1.2,0.9)$. Voters are uniform on the square, and $\phi=\tfrac14$. Because some feasible deviations would win or lose with certainty at large $S$, the figure illustrates the game with clipped winning probabilities. At low unpredictability, $L$ leads because it is closer in baseline. The trailing party $R$ moderates more. Most of the movement of $R$'s candidate occurs in baseline, where $R$ is farther from the centre.

As $S$ rises, $L$'s advantage shrinks, and most of $L$'s movement occurs in responsiveness. At $S=4/3$, the parties tie; beyond that level, $R$ leads and $L$ concedes faster. At the tie, both parties face the same electoral stakes and choose the same $w$. Nomination therefore leaves the threshold unchanged.

Each candidate also moves on the dimension where its party already has an advantage: $L$ moderates slightly in baseline, and $R$ in responsiveness. Proposition~\ref{prop:nomination} requires the same $w_J$ in both dimensions, so a party concedes in both dimensions at once, including the dimension where it is ahead. With several issues, Proposition~\ref{prop:asymmetric} holds issue by issue when each issue has its own state and the states are independent. Under Assumption~\ref{ass:median} in every direction of the larger type space, or locally under local regularity (Assumption~\ref{ass:local}), a rise in the unpredictability of events about one issue moves the lead toward the party whose responsiveness on that issue is closer to the decisive voter's (Proposition~\ref{prop:issues}).

\begin{figure}[tbp]
\centering
\begin{minipage}[b]{0.49\linewidth}
\centering
\begin{tikzpicture}[scale=2.8, every node/.style={font=\small}]
  % preferences at I_q = 4, with switchers (relative to I_q = 1/4) in lighter shades
  \fill[blue!22] (0,-1) -- (1.268,-1) -- (0.898,0.206) -- (0,0.464) -- cycle;
  \fill[red!22]  (0.898,0.206) -- (0.655,1) -- (2,1) -- (2,-0.110) -- cycle;
  \fill[red!8]   (0.898,0.206) -- (2,-0.110) -- (2,-1) -- (1.268,-1) -- cycle;
  \fill[blue!8]  (0.898,0.206) -- (0,0.464) -- (0,1) -- (0.655,1) -- cycle;
  \fill[red!60] (0.655,1) -- (0.693,1) -- (1.307,-1) -- (1.268,-1) -- cycle;
  \draw[gray!60] (0,-1) rectangle (2,1);
  \node[above] at (1,1.02) {(a) Equilibrium candidates};
  \node[below] at (1,-1.04) {$\alpha$ (responsiveness)};
  \node[rotate=90,above] at (-0.06,0) {$\beta$ (baseline)};
  % indifferent voters at I_q = 1/4 (dotted) and I_q = 4 (solid); halving line at I_q = 4 (dashed)
  \draw[thick,dotted] (0,0.464) -- (2,-0.110);
  \node[above] at (1.72,-0.02) {$\psi(\alpha;\tfrac14)$};
  \draw[thick] (0.655,1) -- (1.268,-1);
  \node[left] at (1.20,-0.85) {$\psi(\alpha;4)$};
  \draw[thick,dashed] (0.693,1) -- (1.307,-1);
  \node[right] at (1.29,-0.85) {$\psi_M(\alpha;4)$};
  % paths of the candidates as I_q rises from 0 to 4
  \draw[gray,thick,->] plot[smooth] coordinates {(0.337,-0.083) (0.382,-0.077) (0.424,-0.072) (0.491,-0.064) (0.527,-0.059) (0.597,-0.050) (0.648,-0.044) (0.676,-0.040)};
  \draw[gray,thick,->] plot[smooth] coordinates {(1.131,0.589) (1.127,0.572) (1.124,0.559) (1.120,0.541) (1.118,0.532) (1.114,0.514) (1.111,0.500) (1.109,0.490)};
  % party ideals (diamonds)
  \draw[thick] (0.2,-0.1) +(0,0.05) -- +(0.05,0) -- +(0,-0.05) -- +(-0.05,0) -- cycle;
  \node[below] at (0.2,-0.15) {$\theta_L$};
  \draw[thick] (1.2,0.9) +(0,0.05) -- +(0.05,0) -- +(0,-0.05) -- +(-0.05,0) -- cycle;
  \node[right] at (1.25,0.9) {$\theta_R$};
  % candidates at I_q = 1/4 (gray) and I_q = 4 (black)
  \fill[gray!60] (0.382,-0.077) circle (0.03);
  \fill[gray!60] (1.127,0.572) circle (0.03);
  \fill (0.676,-0.040) circle (0.03) node[below right] {$\theta_l$};
  \fill (1.109,0.490) circle (0.03) node[right] {$\theta_r$};
  % switching labels
  \node at (1.65,-0.62) {$l\to r$};
  \node at (0.22,0.78) {$r\to l$};
  % candidates' midpoint at I_q = 4 and median point
  \draw[fill=white] (0.8925,0.225) circle (0.03);
  \node[left] at (0.86,0.25) {$\theta_m$};
  \fill (1,0) ++(-0.025,-0.025) rectangle ++(0.05,0.05);
  \node[right] at (1.05,0.02) {$\theta_M$};
\end{tikzpicture}
\end{minipage}\hfill
\begin{minipage}[b]{0.49\linewidth}
\centering
\begin{tikzpicture}
\begin{axis}[
  width=7.2cm, height=6.4cm,
  xmin=0, xmax=4, ymin=0.4, ymax=1.1,
  xlabel={$S$}, title={(b) Evolution with unpredictability},
  axis lines=left, every axis plot/.append style={thick},
  tick label style={font=\small}, label style={font=\small}, title style={font=\small}]
  \addplot[black] table {
    0.020 0.5828
    0.219 0.5605
    0.418 0.5433
    0.617 0.5299
    0.816 0.5192
    1.015 0.5106
    1.214 0.5036
    1.413 0.4978
    1.612 0.4929
    1.811 0.4888
    2.010 0.4853
    2.209 0.4822
    2.408 0.4795
    2.607 0.4772
    2.806 0.4751
    3.005 0.4733
    3.204 0.4717
    3.403 0.4702
    3.602 0.4689
    3.801 0.4677
    4.000 0.4667
  };
  \addplot[black, dashed] table {
    0.020 0.4630
    0.219 0.5487
    0.418 0.6188
    0.617 0.6762
    0.816 0.7240
    1.015 0.7641
    1.214 0.7984
    1.413 0.8281
    1.612 0.8542
    1.811 0.8772
    2.010 0.8979
    2.209 0.9166
    2.408 0.9335
    2.607 0.9491
    2.806 0.9634
    3.005 0.9767
    3.204 0.9891
    3.403 1.0007
    3.602 1.0116
    3.801 1.0218
    4.000 1.0315
  };
  \addplot[black, dashdotted] table {
    0.020 1.0379
    0.219 0.9945
    0.418 0.9569
    0.617 0.9242
    0.816 0.8957
    1.015 0.8708
    1.214 0.8487
    1.413 0.8291
    1.612 0.8115
    1.811 0.7957
    2.010 0.7813
    2.209 0.7681
    2.408 0.7561
    2.607 0.7449
    2.806 0.7346
    3.005 0.7249
    3.204 0.7160
    3.403 0.7075
    3.602 0.6996
    3.801 0.6921
    4.000 0.6850
  };
  \addplot[gray, dotted, domain=0:4] {0.5};
  \draw[gray, dotted] (axis cs:1.333,0.4) -- (axis cs:1.333,1.1);
  \node[font=\small, anchor=west] at (axis cs:2.6,0.52) {$\Pr(l\text{ wins})$};
  \node[font=\small, anchor=west] at (axis cs:1.55,0.94) {$D^*$};
  \node[font=\small, anchor=west] at (axis cs:2.4,0.65) {$\|\theta_r-\theta_l\|$};
\end{axis}
\end{tikzpicture}
\end{minipage}
\caption{Equilibrium nominations and policies. Party ideals are $\theta_L=(0.2,-0.1)$ and $\theta_R=(1.2,0.9)$, voters are uniform on the square, and $\phi=\tfrac14$, with $B=0$. In panel (a), diamonds mark party ideals, gray dots mark nominees at $S=\tfrac14$, and black dots mark nominees at $S=4$. Arrows trace the nominations as $S$ rises. Shading shows voter preferences at $S=4$; lighter regions contain voters who have switched since $S=\tfrac14$. The dotted and solid lines are the indifference lines at those two levels. The dashed median line shows that $r$ leads at $S=4$; the dark red strip is the winning margin. In panel (b), $l$'s winning probability crosses $\tfrac12$ at $S=4/3$. The distance between candidate types falls while their expected squared policy gap $D^*$ rises.}
\label{fig:equilibrium}
\end{figure}

\paragraph{Robustness of polarization.}\label{sec:general-G}
Policy divergence from greater uncertainty or information holds exactly in the symmetric benchmark at every $\phi$. Policy divergence also holds qualitatively for general voter distributions and party ideals under moderate electoral sensitivity and local regularity (Proposition~\ref{prop:general}). Under the same conditions, higher common office benefits reduce both the type gap and the policy gap (Corollary~\ref{cor:benefits}). Under high electoral sensitivity, policy divergence can fail (SI).

\paragraph{Welfare.}
Uncertainty and information affect electoral comparisons through the same quantity, $S$, and still differ in their effects on welfare. Write $U_i^*=\pi U_i^l+(1-\pi)U_i^r$ for a player's expected policy payoff. Party payoffs additionally include $\pi B$ for $L$ and $(1-\pi)B$ for $R$. Recall that the revealed share $R^2=S/\operatorname{Var}(\omega)$ is held fixed when uncertainty changes.

The uncertainty that remains unresolved, $\operatorname{Var}(\omega)-S=\mathbb{E}[\operatorname{Var}_\mu(\omega)]$, and the uncertainty that is resolved, $S$, affect welfare differently. Lemma~\ref{lem:payoff} gives
\[
U_i^j=-\alpha_i^2\big(\operatorname{Var}(\omega)-S\big)-(\Delta\beta_{ij})^2-(\Delta\alpha_{ij})^2S:
\]
unresolved uncertainty costs a player in proportion to her own responsiveness, and resolved uncertainty costs her in proportion to her disagreement with the officeholder's response. Greater uncertainty at a fixed revealed share raises both, so it hurts every player. Better information converts unresolved into resolved uncertainty one for one, so its value to player $i$, $\alpha_i^2-(\Delta\alpha_{ij})^2=\alpha_j(2\alpha_i-\alpha_j)$ per unit, depends on her position. A change that raises uncertainty while lowering $S$, as when events cannot distinguish among extreme states, raises the unresolved uncertainty by more than the rise in uncertainty and lowers the cost of delegation.

\begin{proposition}[Uncertainty hurts everyone]\label{prop:welfare-unc}
Suppose $S>0$ and $\alpha_R\ne\alpha_L$.
\begin{enumerate}
  \item[(a)] For any candidates and any player, $\mathrm{d}U_i^j/\mathrm{d}\operatorname{Var}(\omega)=-(1-R^2)\alpha_i^2-R^2(\alpha_j-\alpha_i)^2\le0$.
  \item[(b)] For any distribution of voters and parties' ideal types satisfying the local regularity condition of Section~\ref{app:small-h} (Assumption~\ref{ass:local}), every player's policy payoff in equilibrium falls with uncertainty for all sufficiently small $\phi>0$.
  \item[(c)] Under Assumption~\ref{ass:median} with symmetric parties, every player's policy payoff in equilibrium falls with uncertainty for every $\phi$.
\end{enumerate}
\end{proposition}

With fixed finite office benefits, the strict welfare conclusions in parts (b) and (c) also hold for the parties including their benefits. At fixed candidates, greater uncertainty increases residual risk and the cost of disagreement over responses. Both effects hurt every player. In equilibrium, nominations and winning probabilities also change. Under moderate electoral sensitivity, the changes in nominations and winning probabilities are too small to offset the direct loss. In the symmetric benchmark, the policy lottery between candidates widens as well, so every player loses at any level of electoral sensitivity.

For better information, the symmetric benchmark gives the clearest characterization. Write $\alpha_L=\alpha_M-\delta$ and $\alpha_R=\alpha_M+\delta$, with $\delta\ne0$.

\begin{proposition}[Information is positional]\label{prop:welfare}
Suppose Assumption~\ref{ass:median} holds, the parties' ideal types are symmetric about the median point, and $\alpha_R\ne\alpha_L$. In the equilibrium of Proposition~\ref{prop:symmetric}, player $i$'s expected policy payoff is
\[
U_i^*=-\alpha_i^2\operatorname{Var}(\omega)-(\beta_M-\beta_i)^2+\alpha_M(2\alpha_i-\alpha_M)\,S-\tfrac14D^*,
\]
and
\[
\frac{\mathrm{d}U_i^*}{\mathrm{d}S}=\alpha_M(2\alpha_i-\alpha_M)-\tfrac14\,\frac{\mathrm{d}D^*}{\mathrm{d}S}.
\]
If $\alpha_M=0$, every player is worse off when unpredictability rises, and by continuity this holds for all $\alpha_M$ close enough to zero. If $\alpha_M\ne0$, the gain is increasing in $\alpha_i$ when $\alpha_M>0$ and decreasing when $\alpha_M<0$: the party whose ideal responsiveness has the same sign as $\alpha_M$ and the larger magnitude gains most, the median voter gains less, and the other party gains least. Player $i$ gains if and only if $\alpha_i\operatorname{sign}(\alpha_M)>|\alpha_M|/2+\big(\mathrm{d}D^*/\mathrm{d}S\big)/\big(8|\alpha_M|\big)$.
\end{proposition}

Because each party's expected benefit is $B/2$, the result also holds for total party payoffs. Information affects policy welfare through two terms. The term $\alpha_M(2\alpha_i-\alpha_M)S$ captures the effect of the response at the candidates' midpoint, which is the median point. The term $-D^*/4$ captures the cost of the lottery between candidates. Better information increases that lottery cost.

Figure~\ref{fig:welfare} illustrates the ordering of gains in Proposition~\ref{prop:welfare} for the parties and the median voter.

\begin{figure}[t]
\centering
\begin{tikzpicture}
\begin{axis}[
  width=11cm, height=7cm,
  xmin=-1.5, xmax=1.5, ymin=-0.6, ymax=2,
  xlabel={$\alpha_M$}, ylabel={$\mathrm{d}U_i^*/\mathrm{d}S$},
  axis lines=left, samples=200, domain=-1.5:1.5, clip=true,
  every axis plot/.append style={thick},
  tick label style={font=\small}, label style={font=\small}]
  \fill[gray!15] (axis cs:-0.1357,-0.6) rectangle (axis cs:0.1357,2);
  \addplot[gray, thin] {0};
  \addplot[red!70!black] {x*(x+1)-0.1541};
  \addplot[black, dashed] {x^2-0.1541};
  \addplot[blue!70!black, dashdotted] {x*(x-1)-0.1541};
  \node[font=\small, anchor=south east, red!70!black] at (axis cs:0.92,1.56) {party $R$};
  \node[font=\small, anchor=south east] at (axis cs:0.92,0.66) {median voter};
  \node[font=\small, anchor=south east, blue!70!black] at (axis cs:1.47,0.42) {party $L$};
  \node[font=\small, anchor=south] at (axis cs:0,1.75) {all lose};
\end{axis}
\end{tikzpicture}
\caption{Welfare effects of information in the symmetric benchmark. The parties have $\alpha_L=\alpha_M-\tfrac12$ and $\alpha_R=\alpha_M+\tfrac12$, with $\beta_R-\beta_L=1$, $B=0$, $\phi=0.05$, and $S=\tfrac12$. Both parties and the median voter lose in the shaded band; at $\alpha_M=0$, every player loses. For $\alpha_M>0$, party $R$ gains most, the median voter next, and party $L$ least. The ordering reverses for $\alpha_M<0$.}
\label{fig:welfare}
\end{figure}

The SI discusses the welfare effect of information without symmetry.

\paragraph{Endogenous learning in office.}
The benchmark model fixes the officeholder's information. Suppose instead that the winner chooses the precision of her signal after the election and bears the cost herself. In a normal example, $\omega\sim\operatorname{Norm}(0,\sigma_\omega^2)$ and the winner observes $s=\omega+\varepsilon$ with independent noise $\varepsilon\sim\operatorname{Norm}(0,\sigma_s^2)$, so that $S=\sigma_\omega^4/(\sigma_\omega^2+\sigma_s^2)$. The winner chooses precision $\rho\ge0$, so $\sigma_s^2=1/\rho$, at cost $\zeta \rho^2/2$ with $\zeta>0$, and the unpredictability of the events that she observes is $S(\rho)=\sigma_\omega^4\rho/(\sigma_\omega^2\rho+1)$. Because she cannot commit to a level of learning before the election, voters and parties anticipate the precision $\rho_j^*$ that a candidate of type $\theta_j$ would choose. Responsiveness then determines both a candidate's use of information and the amount of information that she acquires.

\begin{corollary}[Endogenous learning]\label{cor:endolearn}
\begin{enumerate}
  \item[(a)] For $\sigma_\omega^2>0$, a candidate learns, $\rho_j^*>0$, if and only if $\alpha_j\ne0$; for $\alpha_j\ne0$, her precision rises with $|\alpha_j|$, with uncertainty, and as learning becomes cheaper, that is, as $\zeta$ falls.
  \item[(b)] For $\alpha_j\ne0$, the unpredictability of the events that she observes rises with uncertainty faster than at fixed precision.
  \item[(c)] Let $\alpha_R=-\alpha_L\ne0$. Under the regularity conditions stated in the SI, for sufficiently low electoral sensitivity each party nominates a candidate who is more responsive than itself, and the expected squared policy gap rises with uncertainty and as learning becomes cheaper, faster with uncertainty than at fixed precision.
\end{enumerate}
\end{corollary}

The force behind part (c) is the lack of commitment. The officeholder bears the cost of learning, and her party also benefits from the officeholder's learning. A party cannot require its nominee to learn more once the nominee is in office, so the party chooses a nominee whose own responsiveness makes more learning worthwhile to the nominee: a candidate more responsive than the party itself. This selection motive works against the moderation of the benchmark model and adds to polarization. Every officeholder with nonzero responsiveness learns, even when every player would prefer less information, as when the median voter wants no response (Proposition~\ref{prop:welfare}). A party can commit to no learning only by nominating a candidate with $\alpha_j=0$.

\paragraph{Learning before the election.}
The benchmark model also assumes that all learning occurs after the election. Suppose instead that public news $s^1$ arrives before nomination and that the winner observes $s^2$ after the election. Conditional on $s^1$, the game is the benchmark model with each player's baseline shifted to $\beta_i+\alpha_i\mathbb{E}[\omega\mid s^1]$ and with unpredictability equal to that of the winner's later signal. With early news, part of the disagreement over future responses becomes disagreement over known positions. With a median point and symmetric parties, the symmetric benchmark applies after each realization of the news, and I select its symmetric equilibrium.

\begin{corollary}[Timing of information]\label{cor:timing}
Suppose Assumption~\ref{ass:median} holds and the parties' ideal types are symmetric about the median point.
\begin{enumerate}
  \item[(a)] With candidates fixed, the expected squared policy gap depends only on the officeholder's final information, whatever part of it arrives before the election.
  \item[(b)] In the symmetric continuation equilibrium after every realization of the news, holding the officeholder's final information fixed, disclosing more of it before nomination weakly lowers expected polarization, strictly when the additional disclosure changes the parties' disagreement.
  \item[(c)] As electoral sensitivity vanishes, expected polarization approaches its value with candidates at the parties' ideal types, so better final information raises it when $\alpha_R\ne\alpha_L$, as in the benchmark model.
\end{enumerate}
\end{corollary}

With early news, parties can adapt their nominees to the disagreement that they expect. Parties moderate more when the news places their positions far apart and less when the news brings their positions closer. Because the policy gap in equilibrium is concave in the parties' disagreement, this adaptation lowers the average gap. The effect vanishes as electoral sensitivity tends to zero.


\section{Strategic uncertainty}\label{sec:pre-election}

The previous sections take uncertainty as given. Once a party holds office, it can also change uncertainty, and this section studies, within the same framework, one way in which the changed equilibrium can evolve. Suppose now that an incumbent can raise uncertainty before the challenger party chooses its nominee. The incumbent's type $\theta_{\mathcal I}$ is fixed because she runs for reelection. Uncertainty rises with the revealed share $R^2>0$ held fixed, as elsewhere in the paper, so unpredictability $S=R^2 \sigma_\omega^2$ rises with uncertainty. As in the main model, winning brings an office benefit in addition to policy payoffs. Here the benefits may differ: the incumbent receives $B_{\mathcal I}$, and the challenger party receives $B_{\mathcal C}$. Let $\pi=\tfrac12+\phi M$ be the incumbent's probability of winning. The incumbent's payoff is $\Pi_{\mathcal I}=\pi\,(U_{\mathcal I}^{\mathcal I}+B_{\mathcal I})+(1-\pi)\,U_{\mathcal I}^c$, and the challenger party chooses $\theta_c$ to maximize $(1-\pi)(U_{\mathcal C}^c+B_{\mathcal C})+\pi\,U_{\mathcal C}^{\mathcal I}$. To compare benefits with electoral sensitivity, hold $\phi B_{\mathcal I}$ and $\phi B_{\mathcal C}$ fixed as $\phi$ varies. As $\phi\to0$, the challenger's nominee moves a fixed fraction $2\phi B_{\mathcal C}/(1+2\phi B_{\mathcal C})$ of the way from its party's ideal type to the median point. The nominee's responsiveness tends to
\[
\overline{\alpha}_c\equiv\alpha_M+\frac{\alpha_{\mathcal C}-\alpha_M}{1+2\phi B_{\mathcal C}}.
\]

\begin{proposition}[Strategic uncertainty]\label{prop:strategic}
Suppose Assumption~\ref{ass:median} holds, the incumbent's type and the challenger party's ideal type are distinct and interior, uncertainty lies in a compact interval of positive values, and the incumbent is closer to the median voter in responsiveness than the challenger's nominee, $|\overline{\alpha}_c-\alpha_M|>|\alpha_{\mathcal I}-\alpha_M|$. If
\[
\phi B_{\mathcal I}>\overline{B}\equiv\frac{(1-R^2)\,\alpha_{\mathcal I}^2+\tfrac12R^2\,(\overline{\alpha}_c-\alpha_{\mathcal I})^2}{R^2\big[(\overline{\alpha}_c-\alpha_M)^2-(\alpha_{\mathcal I}-\alpha_M)^2\big]},
\]
then there is $\overline{\phi}>0$ such that for all $\phi\le\overline{\phi}$:
\begin{enumerate}
  \item[(a)] the incumbent's probability of winning and her equilibrium payoff strictly increase with uncertainty;
  \item[(b)] if $\overline{\alpha}_c\ne\alpha_{\mathcal I}$, the expected squared gap between the policies of the incumbent and the challenger's nominee strictly increases with uncertainty;
  \item[(c)] if $\overline{\alpha}_c\ne\alpha_{\mathcal I}$, every player's policy payoff, excluding office benefits, strictly decreases with uncertainty.
\end{enumerate}
For a higher revealed share at a fixed level of uncertainty, the winning-probability and polarization conclusions in (a) and (b) also hold, and the incumbent gains with the lower threshold $\underline B\equiv\big[\tfrac12(\overline{\alpha}_c-\alpha_{\mathcal I})^2-\alpha_{\mathcal I}^2\big]\big/\big[(\overline{\alpha}_c-\alpha_M)^2-(\alpha_{\mathcal I}-\alpha_M)^2\big]$, which is negative for a sufficiently responsive incumbent.
\end{proposition}

The incumbent raises uncertainty, and the higher uncertainty lowers every player's policy payoff, including the incumbent's own and the median voter's (part (c)). Only the incumbent's payoff including the office benefit rises. With greater uncertainty, responsiveness receives more weight in the electoral comparison. The incumbent has the advantage in responsiveness, so her probability of keeping office rises. With a large enough benefit, the electoral gain outweighs the policy loss. Policy polarizes in the environment that the incumbent creates, because uncertainty amplifies the remaining difference in responsiveness between the incumbent and the challenger's nominee. The condition on the challenger matters. A challenger party that values office highly moves its nominee toward the median point. Once the nominee's responsiveness is closer to the median voter's than the incumbent's, the incentive reverses: the incumbent would rather reduce uncertainty. Raising uncertainty therefore requires an opposition whose nominee remains farther from the median voter in responsiveness than the incumbent.

Voters become more likely to reelect such an incumbent, and the uncertainty that the incumbent creates lowers the voters' payoffs as well. Voters choose between candidates within the environment that the incumbent has created, and electing the challenger would not remove the uncertainty. Electing the challenger would replace the incumbent's response with one that the decisive voter likes less, because the incumbent is closer to the decisive voter in responsiveness. The incumbent can thus become the lesser evil in the conditions that she created. To an outside observer, the incumbent's choice may look like a failure of accountability. In the model, the incumbent chooses the uncertainty as a strategy.

As in models of strategic challenger entry \citep{GordonHuberLanda2007}, the challenger responds to the incumbent, here through its nomination. The incumbent can expand information for the same reason, with a lower threshold, because better information adds nothing to her own residual risk. With several independent issues, the result holds issue by issue by the same argument as Proposition~\ref{prop:asymmetric}: an incumbent can gain by raising uncertainty about an issue on which she is closer to the median voter in responsiveness than the challenger's nominee, whatever her positions on the other issues (Proposition~\ref{prop:issues}).

\section{Reinterpreting realignment}\label{sec:evidence}

A voter's preferences combine a baseline with a responsiveness to each state. In ordinary times, alignment is governed mostly by baselines and by responsiveness to the states that are already uncertain. Responsiveness to other states is latent: responsiveness to those states matters little for votes, and observers rarely measure it. A shock that raises uncertainty about one state shifts weight from baselines toward responsiveness to that state and rotates the line that divides the electorate (Propositions~\ref{prop:realign} and~\ref{prop:multistate}). The response to the shock becomes consequential, and voters must choose a side. When that responsiveness is only weakly correlated with the old alignment, the new alignment runs nearly orthogonal to the old alignment.

\begin{remark}[Cross-cutting realignment]\label{rem:crosscut}
Suppose that a shock raises the unpredictability of events about a state, and that responsiveness to that state is weakly correlated with baselines. As the dividing line rotates (Proposition~\ref{prop:realign}), the voters who switch come from both sides of the old dividing line, and the new coalitions cut across the old alignment.
\end{remark}

Partisan alignments can reverse in this way: the alignment of US states in the presidential election of 2000 was close to the reverse of that in 1896 \citep{MillerSchofield2003}. The weak correlation also makes the extremism that follows a shock hard to read in pre-shock left--right terms, and observers struggle to identify the source of the new cleavage. A large and stable centre in ordinary times reveals little about how the centre will react to a shock.

\paragraph{Realignment after 2008.}
Remark~\ref{rem:crosscut} and Proposition~\ref{prop:asymmetric} apply to the realignment after 2008. After the financial crisis of 2008, mainstream parties lost to challengers on both sides \citep{Funke2016,Hernandez2016}. In the model's terms, an ordinary recession is a realization of the state, a bad draw from a distribution that voters already know. A financial crisis is a structural shock: the crisis shows that financial conditions can vary far more than voters had believed, and future choices over bailouts, austerity, and regulation become more consequential. In each country, the electorate that the new right assembled cut across the old alignment. In the United States, districts exposed to Chinese import competition replaced moderate legislators with more extreme ones at both ends \citep{Autor2020}, and between 11 and 13 percent of Trump's 2016 voters had voted for Obama in 2012 \citep{Skelley2017}. In Britain, Leave identities cut across party lines \citep{HoboltLeeperTilley2021}, and Reform UK drew mostly on 2019 Conservatives, a coalition that had itself absorbed Labour Leavers \citep{HeathEtAl2024}. In Germany, the AfD's largest gain in 2025 came from former non-voters, and much of the rest from the parties of the centre \citep{NeuPokorny2025}. In France, the RN's working-class vote has been traced both to former Communist voters and to conservative or previously apolitical workers \citep{Perrineau1997,Mayer2002}. Table~\ref{tab:episodes} summarizes these and other episodes. The model offers one interpretation of this pattern, which the evidence alone does not establish. Voters who were far apart on the old axis may share a desired response to immigration, trade, or sovereignty. A shock that makes that response consequential then brings them together. Once a party's responsiveness is closer to that of decisive voters, the party also acquires a policy-based electoral advantage (Proposition~\ref{prop:asymmetric}).

\begin{table}[t]
\centering
\small
\renewcommand{\arraystretch}{1.3}
\begin{tabular}{@{}L{2.6cm}L{3.0cm}L{3.6cm}L{2.6cm}L{2.8cm}@{}}
\toprule
{\bfseries\boldmath Shock} & {\bfseries\boldmath Response at stake} & {\bfseries\boldmath Political change} & {\bfseries\boldmath Gainer} & {\bfseries\boldmath Sources} \\
\hline\hline
Depression, Germany (1929--33) & Recovery, order, redistribution & Working poor to Nazis, unemployed to Communists & Both extremes, right more & \citet{King2008}; \citet{Hamilton1982} \\ \midrule
Financial crisis (2008) & Bailouts, austerity, regulation & Mainstream losses, fragmentation & Radical right; radical left in the South & \citet{Funke2016}; \citet{Hernandez2016} \\ \midrule
China shock (2000s) & Protection, transfers & Moderates replaced at both ends & Both; Republicans more & \citet{Autor2020} \\ \midrule
Brexit (2016--) & Trade, immigration & Radical right rose despite costs & Radical right & \citet{BornEtAl2019}; \citet{Curtice2025} \\ \midrule
COVID-19 (2020--) & Lockdowns, vaccination & Partisan divergence & --- & \citet{Allcott2020} \\ \midrule
Climate change & Mitigation, adaptation & Predicted realignment & Side closer in responsiveness & \citet{IPCC2023} \\ \midrule
Artificial intelligence & Labour adjustment, state use of AI & Predicted realignment, wider policy gap & Side closer in responsiveness & \citet{Eloundou2024}; \citet{Engstrom2020} \\
\bottomrule
\end{tabular}
\caption{Episodes interpreted through the model. The last two rows are predictions.}
\label{tab:episodes}
\end{table}

\paragraph{The interwar parallel.}
In 1928 the Nazis won 2.6 percent of the vote \citep{Evans2003}; by 1932, the Nazis and Communists together held a majority of seats. The Nazi electorate came from across classes and parties \citep{Hamilton1982,Childers1983,Falter1991}, democracy collapsed as conservative politicians tried to co-opt the movement \citep{Hett2018}, and across interwar Europe right-wing anti-system parties grew where hardship persisted \citep{deBromhead2013}. Classic accounts describe particular configurations of baseline and responsiveness. \citet{Lipset1960} located fascism in the liberal middle class: in the model's terms, a middle baseline with a strong desired response to economic collapse (Figure~\ref{fig:median-types}(c)), a correlation that \citet{Hamilton1982} and \citet{Childers1983} dispute. The reactionary modernism of the Weimar right \citep{Herf1984} attached strong responsiveness to a radical-right baseline.

\begin{remark}[A responsive centre]\label{rem:centre}
Suppose that voters with middle baselines differ most in responsiveness, as in Figure~\ref{fig:median-types}(c). A rise in unpredictability then reclassifies mostly middle-baseline voters, and these voters split toward both candidates according to their responsiveness.
\end{remark}

The leading economic account of the Nazi vote shows what is lost without this distinction.

\paragraph{What a retrospective reading misses.}
\citet{King2008} find that, among groups severely affected by the Depression, the working poor moved toward the Nazis and the unemployed toward the Communists, and they interpret these movements as economic voting. Retrospective punishment explains why voters abandoned the governing parties. Retrospective punishment does not explain why voters turned to the Nazis and the Communists over the established opposition, or why groups hit by the same crisis moved to opposite extremes. In the model, voters choose the side whose response they prefer, given the state that the Depression made consequential, and groups that want different responses move in opposite directions. Remark~\ref{rem:centre} describes the split of the centre toward both extremes. Recent reactions to costly outcomes look irrational on a retrospective reading for the same reason. Estimates put Brexit's cost at several percent of UK output \citep{BornEtAl2019,AlabreseEtAl2026b}, and radical-right support rose afterward, most where the local costs were largest \citep{AlabreseEtAl2026}. In the United States, Trump's first-term trade war imposed import tariffs in 2018 and 2019, aimed mainly at China. The tariffs did not raise employment in protected industries, and foreign retaliation cost jobs. Residents of exposed regions nevertheless became more supportive of tariffs and more likely to vote for Trump in 2020 \citep{AutorBeckDornHanson2024}. The authors attribute this support to voters who doubted the economic benefits of tariffs but valued the intention to confront foreign competition. In the model, valuing the intention is a preference over responses: a voter can judge a past decision costly and still prefer the same side's response to the circumstances ahead.

\begin{remark}[Costly outcomes]\label{rem:costly}
In the model, a vote compares the candidates' responses to the state ahead. If a costly outcome raises uncertainty, support shifts toward the side that is closer to decisive voters in responsiveness (Proposition~\ref{prop:asymmetric}), even if that side is responsible for the outcome.
\end{remark}

The prospective channel can coexist with retrospective punishment and with the explanations reviewed in the introduction.

\paragraph{Accountability.}
Remark~\ref{rem:hardening}(b) also bears on scandals. A scandal shifts every voter's comparison against the candidate, and once uncertainty is high, the share of voters whom it can move is small. Trump has survived allegations that would have ended a candidacy in calmer times, and radical-right leaders elsewhere have survived theirs. These politicians still try to hide their scandals, so a scandal still costs votes, but it costs fewer. \citet{Rottinghaus2025} documents that fewer voters have been swayed by scandals over time. In a survey experiment in Germany, supporters of radical parties stayed with corrupt politicians mainly when they lacked a viable alternative, and they switched as often as other voters when they found an alternative acceptable \citep{JankowskiJuenLewandowsky2022}. In the model, greater uncertainty raises the cost of the other side's response, so fewer alternatives remain viable. Close elections can still turn on small shifts, so the hardening appears as stable approval combined with narrow margins.

\paragraph{Shocks under way.}
Proposition~\ref{prop:welfare}, Corollary~\ref{cor:timing}, and Proposition~\ref{prop:multistate} bear on shocks now under way. COVID-19 was a shock on both margins: it raised uncertainty about health and economic conditions, and the real-time data that soon became available improved governments' policy information about those conditions \citep{Chetty2024}. The partisan differences in policy and behaviour that followed common news \citep{Allcott2020,Adolph2022} are consistent with differences in desired responses. Climate change raises uncertainty about damages while monitoring improves \citep{IPCC2023}. Artificial intelligence exposes a large share of tasks to change \citep{Eloundou2024}. Estimated aggregate effects are modest so far \citep{Acemoglu2025}, and earlier automation raised radical-right support among exposed workers \citep{Frey2018,Anelli2021}. AI also expands the state's policy information \citep{Engstrom2020}. Most empirical work on information shocks studies what voters learn \citep{Prior2007,Falck2014,Campante2018,Guriev2021,Zhuravskaya2020}; the channel here concerns the officeholder's policy information, so the growth of the state's capacity to gather information is directly relevant \citep{Brambor2020}. Because AI affects several states at once, the realignment that AI produces depends on agreement over responses across domains (Proposition~\ref{prop:multistate}). Realignment after 2008 turned on responses to financial distress and need not predict the realignment that AI produces.

\paragraph{Guiding measurement.}
Before the EU referendum, most forecasts expected Remain to win. Combined forecasts gave Remain a chance of about three in four \citep{FisherShorrocks2016}, and every final poll overstated its vote \citep{BPC2016}. Forecasts of the US presidential election gave Clinton about a 90 percent chance of winning. State polls underestimated Trump in the Upper Midwest, mainly because of a late swing toward him and because many state polls did not weight for education \citep{Kennedy2018}. Sampling and weighting explain part of these misses. The model points to a further problem.

The instruments used to read the electorate were built for an alignment governed by baselines. Left--right self-placement is the clearest case, and its content had been shifting. In the Netherlands between 1980 and 2006, left--right self-placement came to depend more on immigration attitudes and less on redistribution \citep{deVries2013}. Respondents also attach systematically different meanings to ``left'' and ``right'' \citep{Bauer2017}. These studies document the shift without identifying its cause. Remark~\ref{rem:measurement} supplies a cause. A survey question about policy today elicits a citizen's preferred policy at the expected state, $\beta_i+\alpha_i\mathbb{E}[\omega]$, which equals her baseline and reveals nothing about her responsiveness. Her comparison of two candidates, $U_i^r-U_i^l$, has slope $2(\beta_r-\beta_l)$ in $\beta_i$ and slope $2(\alpha_r-\alpha_l)S$ in $\alpha_i$, so as $S$ rises, the comparison that decides a vote rotates toward responsiveness. The same distinction reconciles the finding that politicians are polarized while voters are not \citep{FiorinaAbrams2009}: surveys measure positions, $\beta_i$, while votes depend on attachments, $|U_i^r-U_i^l|$, which load on $\alpha_i$ through $(\alpha_r-\alpha_l)S$ (Remark~\ref{rem:hardening}). Measures of the distance between a voter and a party compare expected positions, $\big(\mathbb{E}[x_j]-\mathbb{E}[\tau(\omega,\theta_i)]\big)^2=(\beta_j-\beta_i)^2$, while her payoff depends on the cost of delegation, $\Lambda_i^j=(\beta_j-\beta_i)^2+(\alpha_j-\alpha_i)^2S$. By Proposition~\ref{prop:auto-polar}, the expected signed gap between the candidates stays $\beta_r-\beta_l$ while $\Lambda_l^r$ rises with $S$, so measured polarization can stay flat while policies diverge, and a party with the larger $(\beta_j-\beta_i)^2$ can have the smaller $\Lambda_i^j$. Surveys of working-class voters in the United States and the Netherlands show a related pattern: many of these voters favour economically progressive policies and rank economic issues among their priorities, and many of them vote for parties of the right \citep{Kamphorst2025}. \citet{Kamphorst2025} attributes the pattern to misperceptions of the left's priorities and shows that information about Democrats' roll-call votes moves vote intentions. The model offers a complementary reading that does not require misperception. The survey items measure positions at the current state and the importance of a domain, and the vote compares the candidates' responses to the state that has become uncertain. With several states (Proposition~\ref{prop:issues}), let $n=1$ be economic policy and $n=2$ immigration policy. A voter with $\beta_{i1}$ close to $\beta_{L1}$ and $\alpha_{i2}$ close to $\alpha_{R2}$ can prefer $L$ when $S_2$ is small and $R$ when $S_2$ is large. Information about a party's record can also change voters' estimates of the party's responsiveness. Weights calibrated on past elections, in which education mattered less, face the same problem. A large and stable centre on the old scale could therefore coexist with a majority for the side whose response much of that centre preferred. The polling failures are a political version of the critique of \citet{Lucas1976}: predicting the response to a shock requires knowing how citizens want policy to respond to the state that the shock makes consequential.

Because the outcomes of the model are jointly determined, their correlation across countries or over time does not identify an effect of one on another. A regression of radical-right support on polarization, for example, captures their common dependence on uncertainty, and with the change in uncertainty held fixed, the model predicts that the correlation weakens. SI Section~\ref{app:tests} proposes ways to measure responsiveness and to separate the effect of uncertainty from the effect of a change in levels.

\section{Discussion}\label{sec:conclusion}

For decades, a thick centre and a median voter who usually got what she wanted were read as signs that democracy was structurally stable. In the model, the same signs can come from a calm environment. When uncertainty is low and policy information is poor, elections turn on baseline. Voters with very different responsiveness can then sit together in the centre, and measures of left--right placement do not reveal the difference (Remark~\ref{rem:measurement}). A shock reveals the distribution of responsiveness. The dividing line rotates, the centre can split toward both extremes, and the median voter remains decisive without restoring representation. Explanations specific to a country's culture or institutions do not by themselves account for this pattern, which recurs across countries (Section~\ref{sec:evidence}).

\paragraph{The cost of losing.}
The same shock raises the cost of being governed by the other side. The higher cost applies to politicians and, at fixed candidates, to voters whose preferred candidate is closer to them in responsiveness. Democracy survives only if losers prefer to comply \citep{Przeworski1991}, and a higher cost of losing weakens that preference. Backsliding today proceeds mostly through elected incumbents who weaken institutions step by step \citep{Bermeo2016,HaggardKaufman2021,Przeworski2019}. A higher cost of losing gives incumbents a reason to entrench themselves and gives their supporters a reason to tolerate entrenchment. Manipulation belongs to the evolution of the new equilibrium rather than its initial shift: once the Nazis ruled, radio propaganda helped them recruit party members and incited denunciations of Jews \citep{Adena2015}. Formal models of backsliding study how incumbents act on the motive to entrench. The model here explains what strengthens the motive, and Proposition~\ref{prop:strategic} describes one way in which an incumbent can act on it: an incumbent who is closer to decisive voters in responsiveness can gain by raising uncertainty, and electing the challenger would not remove the uncertainty. The tariff changes of 2025, which pushed an index of uncertainty about trade policy to a record high \citep{Caldara2020}, may be an example. Refusals to accept a loss have come from voters as well as politicians, in the attack on the US Capitol, the storming in Bras\'{i}lia, and the rallies against Yoon's impeachment. Supporters justify such refusals by the danger of the other side. Young protesters in Seoul said that only Yoon could save the country from ``the crisis of communization'' \citep{SBS2025}. One young protester could not stay home because ``the country seemed about to collapse'' \citep{FinanceToday2024}, and another believed that ``the country would be ruined'' if Yoon were impeached \citep{MBC2025}. Surveys show the same fear. After Germany's 2025 election, 91 percent of AfD voters said that they feared their lives would change too much \citep{AA2025}, and 79 percent of Reform UK voters were pessimistic about the future \citep{Byrne2025}. In France, pollsters described a ``negative politicisation'' driven by fear of the opposing camp \citep{AFP2024}.

\paragraph{Participation.}
Political scientists long treated low turnout and engagement as a problem and proposed remedies up to compulsory voting \citep{Wattenberg2002,Patterson2002,Mair2013,Lijphart1997,Macedo2005,Fiorina1999,Kornhauser1959,LiftEveryVoice2020}. Turnout in the 2020 US presidential election was the highest since 1900 \citep{McDonald2020}, and far more voters say that it matters who wins \citep{Pew2020}. In the calculus of voting, a larger stake raises the benefit of voting \citep{RikerOrdeshook1968}. In the model, the stake rises with the cost of losing, and under the conditions of Proposition~\ref{prop:welfare-unc}, every voter's expected policy payoff falls at the same time. Participation driven by the fear of losing is then a channel of democratic instability. The Nazi movement grew in this way before conservative politicians brought the movement into government \citep{Hett2018}. \citet{Lipset1960} drew the same lesson from 1930s Germany: a rise in participation can reflect a breakdown of democratic cohesion, and a stable democracy may rest on the belief that ``the outcome of an election will not make too great a difference.'' The model gives this observation a mechanism, because uncertainty and information raise the difference that an election makes.

\paragraph{Implications.}
If uncertainty and policy information keep rising, as with climate change, pandemics, and artificial intelligence, the model implies further realignment, more polarized policy, and more voters and politicians unwilling to accept defeat. Policymakers therefore need to manage uncertainty, policy information, and the timing of policy information. Investments in the state's policy information, including artificial intelligence, are also distributive choices, because better information benefits some voters at the expense of others. Institutional design should give officeholders incentives to reduce uncertainty, beyond limiting an incumbent's ability to manufacture uncertainty. A higher office benefit narrows the policy gap at a given level of uncertainty and also strengthens an incumbent's motive to raise uncertainty. The motive disappears when the opposition's nominee is closer than the incumbent to decisive voters in responsiveness. Policy information that arrives before parties choose candidates lowers expected polarization whoever governs.

\paragraph{Further research.}
Apart from the incumbent's choice of uncertainty and the officeholder's choice of precision, the paper treats uncertainty and the revealed share as exogenous. In a dynamic setting, both would be determined in equilibrium: a policy response can change later states and their variance, policies lose effectiveness as the world changes \citep{CallanderMartin2017}, and officeholders and voters learn from policy outcomes \citep{Callander2011AER,Callander2011APSR,CallanderHummel2014}. Who governs then shapes the information environment of the next election. Candidates could also communicate their responsiveness strategically, and an unconventional or offensive style may serve as a costly signal that a candidate would respond differently from established politicians.


\clearpage
\begin{thebibliography}{99}
\setlength{\itemsep}{2pt}

\bibitem[{Abramowitz and Webster(2016)}]{AbramowitzWebster2016}
Abramowitz, A. I., and S. Webster. 2016. ``The Rise of Negative Partisanship and the Nationalization of U.S. Elections in the 21st Century.'' \textit{Electoral Studies} 41: 12--22.

\bibitem[{Acemoglu(2025)}]{Acemoglu2025}
Acemoglu, D. 2025. ``The Simple Macroeconomics of AI.'' \textit{Economic Policy} 40 (121): 13--58.

\bibitem[{Acemoglu, Kong, and Ozdaglar(2026)}]{AcemogluKongOzdaglar2026}
Acemoglu, D., D. Kong, and A. Ozdaglar. 2026. ``AI, Human Cognition and Knowledge Collapse.'' NBER Working Paper 34910.

\bibitem[{Adena et al.(2015)}]{Adena2015}
Adena, M., R. Enikolopov, M. Petrova, V. Santarosa, and E. Zhuravskaya. 2015. ``Radio and the Rise of the Nazis in Prewar Germany.'' \textit{Quarterly Journal of Economics} 130 (4): 1885--1939.

\bibitem[{Adolph et~al.(2022)}]{Adolph2022}
Adolph, C., K. Amano, B. Bang-Jensen, N. Fullman, and J. Wilkerson. 2022. ``Pandemic Politics: Timing State-Level Social Distancing Responses to COVID-19.'' \textit{Journal of Health Politics, Policy and Law} 47 (2).

\bibitem[{Agence France-Presse(2024)}]{AFP2024}
Agence France-Presse. 2024. ``Far Right Scents Power as Tense France Braces for Snap Vote.'' 29 June.

\bibitem[{Ahir, Bloom, and Furceri(2022)}]{Ahir2022}
Ahir, H., N. Bloom, and D. Furceri. 2022. ``The World Uncertainty Index.'' NBER Working Paper 29763.

\bibitem[{Alabrese et~al.(2026a)}]{AlabreseEtAl2026}
Alabrese, E., J. Edenhofer, T. Fetzer, and S. Wang. 2026a. ``Levelling Up by Levelling Down: The Economic and Political Costs of Brexit.'' NICEP Working Paper 2026-06, University of Nottingham.

\bibitem[{Alabrese et~al.(2026b)}]{AlabreseEtAl2026b}
Alabrese, E., J. Edenhofer, T. Fetzer, and S. Wang. 2026b. ``Measuring the Regional Economic Cost of Brexit: Evidence as of 2026.'' Warwick Economics Research Paper 1617, University of Warwick.

\bibitem[{Aldrich(1983)}]{Aldrich1983}
Aldrich, J. H. 1983. ``A Downsian Spatial Model with Party Activism.'' \textit{American Political Science Review} 77 (4): 974--990.

\bibitem[{Alesina and Cukierman(1990)}]{AlesinaCukierman1990}
Alesina, A., and A. Cukierman. 1990. ``The Politics of Ambiguity.'' \textit{Quarterly Journal of Economics} 105 (4): 829--850.

\bibitem[{Algan et~al.(2017)}]{Algan2017}
Algan, Y., S. Guriev, E. Papaioannou, and E. Passari. 2017. ``The European Trust Crisis and the Rise of Populism.'' \textit{Brookings Papers on Economic Activity} 2017 (2): 309--400.

\bibitem[{Allcott and Gentzkow(2017)}]{AllcottGentzkow2017}
Allcott, H., and M. Gentzkow. 2017. ``Social Media and Fake News in the 2016 Election.'' \textit{Journal of Economic Perspectives} 31 (2): 211--236.

\bibitem[{Allcott et~al.(2020)}]{Allcott2020}
Allcott, H., L. Boxell, J. Conway, M. Gentzkow, M. Thaler, and D. Yang. 2020. ``Polarization and Public Health: Partisan Differences in Social Distancing during the Coronavirus Pandemic.'' \textit{Journal of Public Economics} 191: 104254.

\bibitem[{Alonso and C\^{a}mara(2016)}]{AlonsoCamara2016}
Alonso, R., and O. C\^{a}mara. 2016. ``Persuading Voters.'' \textit{American Economic Review} 106 (11): 3590--3605.

\bibitem[{Alonso and Matouschek(2008)}]{AlonsoMatouschek2008}
Alonso, R., and N. Matouschek. 2008. ``Optimal Delegation.'' \textit{Review of Economic Studies} 75 (1): 259--293.

\bibitem[{Amador and Bagwell(2013)}]{AmadorBagwell2013}
Amador, M., and K. Bagwell. 2013. ``The Theory of Optimal Delegation with an Application to Tariff Caps.'' \textit{Econometrica} 81 (4): 1541--1599.

\bibitem[{Anadolu Agency(2025)}]{AA2025}
Anadolu Agency. 2025. ``Germany Elections: `Far-Right AfD Rising by Stoking Fears, but Offers No Answers'.'' Interview with Gerd Bohmann, February. Reporting Infratest dimap's election-day survey for ARD.

\bibitem[{Anderson et al.(2005)}]{AndersonEtAl2005}
Anderson, C. J., A. Blais, S. Bowler, T. Donovan, and O. Listhaug. 2005. \textit{Losers' Consent: Elections and Democratic Legitimacy}. Oxford: Oxford University Press.

\bibitem[{Anelli, Colantone, and Stanig(2021)}]{Anelli2021}
Anelli, M., I. Colantone, and P. Stanig. 2021. ``Individual Vulnerability to Industrial Robot Adoption Increases Support for the Radical Right.'' \textit{Proceedings of the National Academy of Sciences} 118 (47): e2111611118.

\bibitem[{Arendt(1951)}]{Arendt1951}
Arendt, H. 1951. \textit{The Origins of Totalitarianism}. New York: Harcourt, Brace.

\bibitem[{Ashworth(2012)}]{Ashworth2012}
Ashworth, S. 2012. ``Electoral Accountability: Recent Theoretical and Empirical Work.'' \textit{Annual Review of Political Science} 15: 183--201.

\bibitem[{Autor et~al.(2020)}]{Autor2020}
Autor, D., D. Dorn, G. Hanson, and K. Majlesi. 2020. ``Importing Political Polarization? The Electoral Consequences of Rising Trade Exposure.'' \textit{American Economic Review} 110 (10): 3139--3183.

\bibitem[{Autor et~al.(2024)}]{AutorBeckDornHanson2024}
Autor, D., A. Beck, D. Dorn, and G. H. Hanson. 2024. ``Help for the Heartland? The Employment and Electoral Effects of the Trump Tariffs in the United States.'' NBER Working Paper 32082.

\bibitem[{Barber(2016)}]{Barber2016}
Barber, M. J. 2016. ``Ideological Donors, Contribution Limits, and the Polarization of American Legislatures.'' \textit{Journal of Politics} 78 (1): 296--310.

\bibitem[{Barker(2025)}]{Barker2025}
Barker, E. 2025. ``Why I Don't Regret Leaving the Democratic Party.'' \textit{The Free Press}, 18 November.

\bibitem[{Bauer et~al.(2017)}]{Bauer2017}
Bauer, P. C., P. Barber\'{a}, K. Ackermann, and A. Venetz. 2017. ``Is the Left-Right Scale a Valid Measure of Ideology? Individual-Level Variation in Associations with `Left' and `Right' and Left-Right Self-Placement.'' \textit{Political Behavior} 39 (3): 553--583.

\bibitem[{Bermeo(2016)}]{Bermeo2016}
Bermeo, N. 2016. ``On Democratic Backsliding.'' \textit{Journal of Democracy} 27 (1): 5--19.

\bibitem[{Bernhardt, C\^{a}mara, and Squintani(2011)}]{BernhardtCamaraSquintani2011}
Bernhardt, D., O. C\^{a}mara, and F. Squintani. 2011. ``Competence and Ideology.'' \textit{Review of Economic Studies} 78 (2): 487--522.

\bibitem[{Bertsou and Caramani(2022)}]{BertsouCaramani2022}
Bertsou, E., and D. Caramani. 2022. ``People Haven't Had Enough of Experts: Technocratic Attitudes among Citizens in Nine European Democracies.'' \textit{American Journal of Political Science} 66 (1): 5--23.

\bibitem[{Besley and Burgess(2002)}]{BesleyBurgess2002}
Besley, T., and R. Burgess. 2002. ``The Political Economy of Government Responsiveness: Theory and Evidence from India.'' \textit{Quarterly Journal of Economics} 117 (4): 1415--1451.

\bibitem[{Besley and Coate(1997)}]{BesleyCoate1997}
Besley, T., and S. Coate. 1997. ``An Economic Model of Representative Democracy.'' \textit{Quarterly Journal of Economics} 112 (1): 85--114.

\bibitem[{Black(1948)}]{Black1948}
Black, D. 1948. ``On the Rationale of Group Decision-Making.'' \textit{Journal of Political Economy} 56 (1): 23--34.

\bibitem[{Blum, Cowburn, and Masket(2024)}]{BlumCowburnMasket2024}
Blum, R. M., M. Cowburn, and S. Masket. 2024. ``Who Decides? Media, MAGA, Money, and Mentions in the 2022 Republican Primaries.'' \textit{Political Research Quarterly} 77 (4): 1314--1332.

\bibitem[{Boix(2022)}]{Boix2022}
Boix, C. 2022. ``AI and the Economic and Informational Foundations of Democracy.'' In J. B. Bullock et~al. (eds.), \textit{The Oxford Handbook of AI Governance}, 707--725. Oxford: Oxford University Press.

\bibitem[{Bonica(2014)}]{Bonica2014}
Bonica, A. 2014. ``Mapping the Ideological Marketplace.'' \textit{American Journal of Political Science} 58 (2): 367--386.

\bibitem[{Bonomi, Gennaioli, and Tabellini(2021)}]{BonomiGennaioliTabellini2021}
Bonomi, G., N. Gennaioli, and G. Tabellini. 2021. ``Identity, Beliefs, and Political Conflict.'' \textit{Quarterly Journal of Economics} 136 (4): 2371--2411.

\bibitem[{Born et~al.(2019)}]{BornEtAl2019}
Born, B., G. J. M\"{u}ller, M. Schularick, and P. Sedl\'{a}\v{c}ek. 2019. ``The Costs of Economic Nationalism: Evidence from the Brexit Experiment.'' \textit{Economic Journal} 129 (623): 2722--2744.

\bibitem[{Brambor et~al.(2020)}]{Brambor2020}
Brambor, T., A. Goenaga, J. Lindvall, and J. Teorell. 2020. ``The Lay of the Land: Information Capacity and the Modern State.'' \textit{Comparative Political Studies} 53 (2): 175--213.

\bibitem[{British Polling Council(2016)}]{BPC2016}
British Polling Council. 2016. Statement on the performance of the final polls in the EU referendum, June 24.

\bibitem[{Buchanan(2025)}]{Buchanan2025}
Buchanan, I. 2025. ``Local Elections 2025: Results and Analysis.'' House of Commons Library Research Briefing 10272.

\bibitem[{Buisseret and Van Weelden(2022)}]{BuisseretVanWeelden2022}
Buisseret, P., and R. Van Weelden. 2022. ``Polarization, Valence, and Policy Competition.'' \textit{American Economic Review: Insights} 4 (3): 341--352.

\bibitem[{Burnham(1970)}]{Burnham1970}
Burnham, W. D. 1970. \textit{Critical Elections and the Mainsprings of American Politics}. New York: W. W. Norton.

\bibitem[{Bursztyn, Egorov, and Fiorin(2020)}]{BursztynEgorovFiorin2020}
Bursztyn, L., G. Egorov, and S. Fiorin. 2020. ``From Extreme to Mainstream: The Erosion of Social Norms.'' \textit{American Economic Review} 110 (11): 3522--3548.

\bibitem[{Byrne(2025)}]{Byrne2025}
Byrne, L. 2025. ``How to Defeat Reform.'' Reporting Hope Not Hate's survey of Reform UK voters. June.

\bibitem[{Caldara et~al.(2020)}]{Caldara2020}
Caldara, D., M. Iacoviello, P. Molligo, A. Prestipino, and A. Raffo. 2020. ``The Economic Effects of Trade Policy Uncertainty.'' \textit{Journal of Monetary Economics} 109: 38--59.

\bibitem[{Callander(2011a)}]{Callander2011AER}
Callander, S. 2011a. ``Searching and Learning by Trial and Error.'' \textit{American Economic Review} 101 (6): 2277--2308.

\bibitem[{Callander(2011b)}]{Callander2011APSR}
Callander, S. 2011b. ``Searching for Good Policies.'' \textit{American Political Science Review} 105 (4): 643--662.

\bibitem[{Callander and Carbajal(2022)}]{CallanderCarbajal2022}
Callander, S., and J. C. Carbajal. 2022. ``Cause and Effect in Political Polarization: A Dynamic Analysis.'' \textit{Journal of Political Economy} 130 (4): 825--880.

\bibitem[{Callander and Hummel(2014)}]{CallanderHummel2014}
Callander, S., and P. Hummel. 2014. ``Preemptive Policy Experimentation.'' \textit{Econometrica} 82 (4): 1509--1528.

\bibitem[{Callander and Martin(2017)}]{CallanderMartin2017}
Callander, S., and G. J. Martin. 2017. ``Dynamic Policymaking with Decay.'' \textit{American Journal of Political Science} 61 (1): 50--67.

\bibitem[{Calvert(1985)}]{Calvert1985}
Calvert, R. L. 1985. ``Robustness of the Multidimensional Voting Model: Candidate Motivations, Uncertainty, and Convergence.'' \textit{American Journal of Political Science} 29 (1): 69--95.

\bibitem[{Campante, Durante, and Sobbrio(2018)}]{Campante2018}
Campante, F., R. Durante, and F. Sobbrio. 2018. ``Politics 2.0: The Multifaceted Effect of Broadband Internet on Political Participation.'' \textit{Journal of the European Economic Association} 16 (4): 1094--1136.

\bibitem[{Canen, Kendall, and Trebbi(2020)}]{Canen2020}
Canen, N., C. Kendall, and F. Trebbi. 2020. ``Unbundling Polarization.'' \textit{Econometrica} 88 (3): 1197--1233.

\bibitem[{Canes-Wrone and Shotts(2007)}]{CanesWroneShotts2007}
Canes-Wrone, B., and K. W. Shotts. 2007. ``When Do Elections Encourage Ideological Rigidity?'' \textit{American Political Science Review} 101 (2): 273--288.

\bibitem[{Carmines and Stimson(1989)}]{CarminesStimson1989}
Carmines, E. G., and J. A. Stimson. 1989. \textit{Issue Evolution: Race and the Transformation of American Politics}. Princeton: Princeton University Press.

\bibitem[{Chetty et~al.(2024)}]{Chetty2024}
Chetty, R., J. N. Friedman, M. Stepner, and the Opportunity Insights Team. 2024. ``The Economic Impacts of COVID-19: Evidence from a New Public Database Built Using Private Sector Data.'' \textit{Quarterly Journal of Economics} 139 (2): 829--889.

\bibitem[{Chicago Tribune Editorial Board(2026)}]{ChicagoTribune2026}
Chicago Tribune Editorial Board. 2026. ``The Number of Americans Who Feel Politically Homeless Is Rising. Understandably So.'' \textit{Chicago Tribune}, January.

\bibitem[{Childers(1983)}]{Childers1983}
Childers, T. 1983. \textit{The Nazi Voter: The Social Foundations of Fascism in Germany, 1919--1933}. Chapel Hill: University of North Carolina Press.

\bibitem[{Chiopris, Nalepa, and Vanberg(2025)}]{ChioprisNalepaVanberg2025}
Chiopris, C., M. Nalepa, and G. Vanberg. 2025. ``A Wolf in Sheep's Clothing: Citizen Uncertainty and Democratic Backsliding.'' \textit{Journal of Politics} 87 (4): 1272--1287.

\bibitem[{Colantone and Stanig(2018)}]{ColantoneStanig2018}
Colantone, I., and P. Stanig. 2018. ``The Trade Origins of Economic Nationalism: Import Competition and Voting Behavior in Western Europe.'' \textit{American Journal of Political Science} 62 (4): 936--953.

\bibitem[{Crawford and Sobel(1982)}]{CrawfordSobel1982}
Crawford, V. P., and J. Sobel. 1982. ``Strategic Information Transmission.'' \textit{Econometrica} 50 (6): 1431--1451.

\bibitem[{Curtice(2025)}]{Curtice2025}
Curtice, J. 2025. ``Reform's Rise Has Been Dizzying, but the Polls Hide a Landmine.'' \textit{GB News}, 29 October.

\bibitem[{Davis, DeGroot, and Hinich(1972)}]{DavisDeGrootHinich1972}
Davis, O. A., M. H. DeGroot, and M. J. Hinich. 1972. ``Social Preference Orderings and Majority Rule.'' \textit{Econometrica} 40 (1): 147--157.

\bibitem[{de~Bromhead, Eichengreen, and O'Rourke(2013)}]{deBromhead2013}
de Bromhead, A., B. Eichengreen, and K. H. O'Rourke. 2013. ``Political Extremism in the 1920s and 1930s: Do German Lessons Generalize?'' \textit{Journal of Economic History} 73 (2): 371--406.

\bibitem[{DellaVigna and Kaplan(2007)}]{DellaVignaKaplan2007}
DellaVigna, S., and E. Kaplan. 2007. ``The Fox News Effect: Media Bias and Voting.'' \textit{Quarterly Journal of Economics} 122 (3): 1187--1234.

\bibitem[{de~Vries, Hakhverdian, and Lancee(2013)}]{deVries2013}
de Vries, C. E., A. Hakhverdian, and B. Lancee. 2013. ``The Dynamics of Voters' Left/Right Identification: The Role of Economic and Cultural Attitudes.'' \textit{Political Science Research and Methods} 1 (2): 223--238.

\bibitem[{Difford(2025)}]{Difford2025}
Difford, D. 2025. ``How Do Britons Feel about Brexit Five Years On?'' YouGov, January 29.

\bibitem[{Dixit and Weibull(2007)}]{DixitWeibull2007}
Dixit, A. K., and J. W. Weibull. 2007. ``Political Polarization.'' \textit{Proceedings of the National Academy of Sciences} 104 (18): 7351--7356.

\bibitem[{Doerr et~al.(2022)}]{Doerr2022}
Doerr, S., S. Gissler, J.-L. Peydr\'{o}, and H.-J. Voth. 2022. ``Financial Crises and Political Radicalization: How Failing Banks Paved Hitler's Path to Power.'' \textit{Journal of Finance} 77 (6): 3339--3372.

\bibitem[{Downs(1957)}]{Downs1957}
Downs, A. 1957. \textit{An Economic Theory of Democracy}. New York: Harper.

\bibitem[{Downs and Rocke(1994)}]{DownsRocke1994}
Downs, G. W., and D. M. Rocke. 1994. ``Conflict, Agency, and Gambling for Resurrection: The Principal-Agent Problem Goes to War.'' \textit{American Journal of Political Science} 38 (2): 362--380.

\bibitem[{Eloundou et~al.(2024)}]{Eloundou2024}
Eloundou, T., S. Manning, P. Mishkin, and D. Rock. 2024. ``GPTs Are GPTs: Labor Market Impact Potential of LLMs.'' \textit{Science} 384 (6702): 1306--1308.

\bibitem[{Ely, Frankel, and Kamenica(2015)}]{ElyFrankelKamenica2015}
Ely, J., A. Frankel, and E. Kamenica. 2015. ``Suspense and Surprise.'' \textit{Journal of Political Economy} 123 (1): 215--260.

\bibitem[{Engstrom et~al.(2020)}]{Engstrom2020}
Engstrom, D. F., D. E. Ho, C. M. Sharkey, and M.-F. Cu\'{e}llar. 2020. \textit{Government by Algorithm: Artificial Intelligence in Federal Administrative Agencies}. Report submitted to the Administrative Conference of the United States.

\bibitem[{Evans(2003)}]{Evans2003}
Evans, R. J. 2003. \textit{The Coming of the Third Reich}. London: Allen Lane.

\bibitem[{Falck, Gold, and Heblich(2014)}]{Falck2014}
Falck, O., R. Gold, and S. Heblich. 2014. ``E-lections: Voting Behavior and the Internet.'' \textit{American Economic Review} 104 (7): 2238--2265.

\bibitem[{Falter(1991)}]{Falter1991}
Falter, J. W. 1991. \textit{Hitlers W\"{a}hler}. Munich: C.~H. Beck.

\bibitem[{Finance Today(2024)}]{FinanceToday2024}
Finance Today. 2024. ``Young People in Their 20s and 30s at the Gwanghwamun Rally: `We Oppose President Yoon's Impeachment Too''' [in Korean]. 28 December. Quotation translated by the author.

\bibitem[{Fiorina(1999)}]{Fiorina1999}
Fiorina, M. P. 1999. ``Extreme Voices: A Dark Side of Civic Engagement.'' In T. Skocpol and M. P. Fiorina (eds.), \textit{Civic Engagement in American Democracy}, 395--425. Washington, DC: Brookings Institution Press.

\bibitem[{Fiorina and Abrams(2009)}]{FiorinaAbrams2009}
Fiorina, M. P., and S. J. Abrams. 2009. \textit{Disconnect: The Breakdown of Representation in American Politics}. Norman: University of Oklahoma Press.

\bibitem[{Fisher and Shorrocks(2016)}]{FisherShorrocks2016}
Fisher, S., and R. Shorrocks. 2016. ``Final Combined EU Referendum Forecast.'' \textit{Elections Etc.}, June 23.

\bibitem[{Folkenflik(2016)}]{Folkenflik2016}
Folkenflik, D. 2016. ``Infowars' Alex Jones Claims Ear of Trump, Pushes `Pizzagate' Fictions.'' NPR, December.

\bibitem[{Frey, Berger, and Chen(2018)}]{Frey2018}
Frey, C. B., T. Berger, and C. Chen. 2018. ``Political Machinery: Did Robots Swing the 2016 US Presidential Election?'' \textit{Oxford Review of Economic Policy} 34 (3): 418--442.

\bibitem[{Fromm(1941)}]{Fromm1941}
Fromm, E. 1941. \textit{Escape from Freedom}. New York: Farrar and Rinehart.

\bibitem[{Fujiwara, M\"{u}ller, and Schwarz(2024)}]{FujiwaraMullerSchwarz2024}
Fujiwara, T., K. M\"{u}ller, and C. Schwarz. 2024. ``The Effect of Social Media on Elections: Evidence from the United States.'' \textit{Journal of the European Economic Association} 22 (3): 1495--1539.

\bibitem[{Funda\c{c}\~{a}o FHC(2021)}]{FHC2021}
Funda\c{c}\~{a}o FHC. 2021. ``The Evolution of Bolsonarism: Qualitative Analysis of the Perceptions of Bolsonaro's Electorate in 2019 and 2020.'' S\~{a}o Paulo: Funda\c{c}\~{a}o FHC.

\bibitem[{Funke, Schularick, and Trebesch(2016)}]{Funke2016}
Funke, M., M. Schularick, and C. Trebesch. 2016. ``Going to Extremes: Politics after Financial Crises, 1870--2014.'' \textit{European Economic Review} 88: 227--260.

\bibitem[{Garicano and Rayo(2025)}]{GaricanoRayo2025}
Garicano, L., and L. Rayo. 2025. ``Training in the Age of AI: A Theory of Career Viability.'' CEPR Discussion Paper 20634.

\bibitem[{Gennaioli and Tabellini(2025)}]{GennaioliTabellini2025}
Gennaioli, N., and G. Tabellini. 2025. ``Presidential Address: Identity Politics.'' \textit{Econometrica} 93 (6): 1937--1967.

\bibitem[{George(2026)}]{George2026}
George, M. 2026. ``Local Elections 2026: Sweeping Reform Gains, Deep Labour Losses.'' \textit{Local Government Chronicle}, May 8.

\bibitem[{Gethin, Mart\'{i}nez-Toledano and Piketty(2022)}]{Gethin2022}
Gethin, A., C. Mart\'{i}nez-Toledano, and T. Piketty. 2022. ``Brahmin Left versus Merchant Right: Changing Political Cleavages in 21 Western Democracies, 1948--2020.'' \textit{Quarterly Journal of Economics} 137 (1): 1--48.

\bibitem[{Glaeser, Ponzetto, and Shapiro(2005)}]{GlaeserPonzettoShapiro2005}
Glaeser, E. L., G. A. M. Ponzetto, and J. M. Shapiro. 2005. ``Strategic Extremism: Why Republicans and Democrats Divide on Religious Values.'' \textit{Quarterly Journal of Economics} 120 (4): 1283--1330.

\bibitem[{Gordon, Huber, and Landa(2007)}]{GordonHuberLanda2007}
Gordon, S. C., G. A. Huber, and D. Landa. 2007. ``Challenger Entry and Voter Learning.'' \textit{American Political Science Review} 101 (2): 303--320.

\bibitem[{Graham and Svolik(2020)}]{GrahamSvolik2020}
Graham, M. H., and M. W. Svolik. 2020. ``Democracy in America? Partisanship, Polarization, and the Robustness of Support for Democracy in the United States.'' \textit{American Political Science Review} 114 (2): 392--409.

\bibitem[{Gratton and Lee(2024)}]{GrattonLee2024}
Gratton, G., and B. E. Lee. 2024. ``Liberty, Security, and Accountability: The Rise and Fall of Illiberal Democracies.'' \textit{Review of Economic Studies} 91 (1): 340--371.

\bibitem[{Grillo and Prato(2023)}]{GrilloPrato2023}
Grillo, E., and C. Prato. 2023. ``Reference Points and Democratic Backsliding.'' \textit{American Journal of Political Science} 67 (1): 71--88.

\bibitem[{Grillo et al.(2024)}]{GrilloLuoNalepaPrato2024}
Grillo, E., Z. Luo, M. Nalepa, and C. Prato. 2024. ``Theories of Democratic Backsliding.'' \textit{Annual Review of Political Science} 27: 381--400.

\bibitem[{Groseclose(2001)}]{Groseclose2001}
Groseclose, T. 2001. ``A Model of Candidate Location When One Candidate Has a Valence Advantage.'' \textit{American Journal of Political Science} 45 (4): 862--886.

\bibitem[{Grossmann and Hopkins(2024)}]{GrossmannHopkins2024}
Grossmann, M., and D. A. Hopkins. 2024. \textit{Polarized by Degrees: How the Diploma Divide and the Culture War Transformed American Politics}. Cambridge: Cambridge University Press.

\bibitem[{Guriev and Papaioannou(2022)}]{GurievPapaioannou2022}
Guriev, S., and E. Papaioannou. 2022. ``The Political Economy of Populism.'' \textit{Journal of Economic Literature} 60 (3): 753--832.

\bibitem[{Guriev, Melnikov, and Zhuravskaya(2021)}]{Guriev2021}
Guriev, S., N. Melnikov, and E. Zhuravskaya. 2021. ``3G Internet and Confidence in Government.'' \textit{Quarterly Journal of Economics} 136 (4): 2533--2613.

\bibitem[{Haggard and Kaufman(2021)}]{HaggardKaufman2021}
Haggard, S., and R. Kaufman. 2021. \textit{Backsliding: Democratic Regress in the Contemporary World}. Cambridge: Cambridge University Press.

\bibitem[{Hall(2019)}]{Hall2019}
Hall, A. B. 2019. \textit{Who Wants to Run? How the Devaluing of Political Office Drives Polarization}. Chicago: University of Chicago Press.

\bibitem[{Hamilton(1982)}]{Hamilton1982}
Hamilton, R. F. 1982. \textit{Who Voted for Hitler?} Princeton: Princeton University Press.

\bibitem[{Heath et~al.(2024)}]{HeathEtAl2024}
Heath, O., C. Prosser, H. Southall, and P. Aucott. 2024. ``The 2024 General Election and the Rise of Reform UK.'' \textit{The Political Quarterly}.

\bibitem[{Helmke, Kroeger, and Paine(2022)}]{HelmkeKroegerPaine2022}
Helmke, G., M. Kroeger, and J. Paine. 2022. ``Democracy by Deterrence: Norms, Constitutions, and Electoral Tilting.'' \textit{American Journal of Political Science} 66 (2): 434--450.

\bibitem[{Herf(1984)}]{Herf1984}
Herf, J. 1984. \textit{Reactionary Modernism: Technology, Culture, and Politics in Weimar and the Third Reich}. Cambridge: Cambridge University Press.

\bibitem[{Hern\'{a}ndez and Kriesi(2016)}]{Hernandez2016}
Hern\'{a}ndez, E., and H. Kriesi. 2016. ``The Electoral Consequences of the Financial and Economic Crisis in Europe.'' \textit{European Journal of Political Research} 55 (2): 203--224.

\bibitem[{Hett(2018)}]{Hett2018}
Hett, B. C. 2018. \textit{The Death of Democracy: Hitler's Rise to Power and the Downfall of the Weimar Republic}. New York: Henry Holt.

\bibitem[{Hobolt, Leeper, and Tilley(2021)}]{HoboltLeeperTilley2021}
Hobolt, S. B., T. J. Leeper, and J. Tilley. 2021. ``Divided by the Vote: Affective Polarization in the Wake of the Brexit Referendum.'' \textit{British Journal of Political Science} 51 (4): 1476--1493.

\bibitem[{Hofstadter(1964)}]{Hofstadter1964}
Hofstadter, R. 1964. ``The Paranoid Style in American Politics.'' \textit{Harper's Magazine}, November, 77--86.

\bibitem[{Hooghe and Marks(2018)}]{HoogheMarks2018}
Hooghe, L., and G. Marks. 2018. ``Cleavage Theory Meets Europe's Crises: Lipset, Rokkan, and the Transnational Cleavage.'' \textit{Journal of European Public Policy} 25 (1): 109--135.

\bibitem[{House of Commons Library(2024)}]{HoCL2024}
House of Commons Library. 2024. ``General Election 2024 Results.'' Research Briefing CBP-10009.

\bibitem[{Howell and Wolton(2018)}]{HowellWolton2018}
Howell, W. G., and S. Wolton. 2018. ``The Politician's Province.'' \textit{Quarterly Journal of Political Science} 13 (2): 119--146.

\bibitem[{Howell, Shepsle, and Wolton(2023)}]{HowellShepsleWolton2023}
Howell, W. G., K. A. Shepsle, and S. Wolton. 2023. ``Executive Absolutism: The Dynamics of Authority Acquisition in a System of Separated Powers.'' \textit{Quarterly Journal of Political Science} 18 (2): 243--275.

\bibitem[{Ide and Talam\`{a}s(2025)}]{IdeTalamas2025}
Ide, E., and E. Talam\`{a}s. 2025. ``Artificial Intelligence in the Knowledge Economy.'' \textit{Journal of Political Economy} 133 (12): 3762--3800.

\bibitem[{IPCC(2023)}]{IPCC2023}
Intergovernmental Panel on Climate Change. 2023. \textit{Climate Change 2023: Synthesis Report. Contribution of Working Groups I, II and III to the Sixth Assessment Report}. Geneva: IPCC.

\bibitem[{Ipsos(2025)}]{Ipsos2025}
Ipsos bva. 2025. ``Fractures fran\c{c}aises 2025.'' Survey for \textit{Le Monde}, Cevipof, Fondation Jean-Jaur\`{e}s and Institut Montaigne, October.

\bibitem[{Iyengar et~al.(2019)}]{Iyengar2019}
Iyengar, S., Y. Lelkes, M. Levendusky, N. Malhotra, and S. J. Westwood. 2019. ``The Origins and Consequences of Affective Polarization in the United States.'' \textit{Annual Review of Political Science} 22: 129--146.

\bibitem[{Izzo(2023)}]{Izzo2023}
Izzo, F. 2023. ``Ideology for the Future.'' \textit{American Political Science Review} 117 (3): 1089--1104.

\bibitem[{Izzo, Invernizzi, and Crosson(2025)}]{IzzoInvernizziCrosson2025}
Izzo, F., G. M. Invernizzi, and J. M. Crosson. 2025. ``Learning in a Complex World: How Multidimensionality Affects Policymaking.'' Working paper.

\bibitem[{Izzo, Martin, and Callander(2023)}]{IzzoMartinCallander2023}
Izzo, F., G. J. Martin, and S. Callander. 2023. ``Ideological Competition.'' \textit{American Journal of Political Science} 67 (3): 687--700.

\bibitem[{Jang(2025)}]{Jang2025}
Jang, S. 2025. ``Who Falls for the Election Fraud Conspiracies? The Interaction Effects of Affective Polarization and Populist Attitudes.'' \textit{Korean Political Science Review} 59 (4): 99--119. In Korean.

\bibitem[{Jankowski, Juen, and Lewandowsky(2022)}]{JankowskiJuenLewandowsky2022}
Jankowski, M., C.-M. Juen, and M. Lewandowsky. 2022. ``Turning a Blind Eye on the Black Sheep: When Are Voters Loyal to Corrupt Politicians?'' \textit{Acta Politica} 57 (4): 765--791.

\bibitem[{Jeong(2025)}]{Jeong2025}
Jeong, Y. 2025. ``Who Supports Martial Law? Korean Voters' Political Attitudes toward Violations of Democratic Norms.'' \textit{Korean Party Studies Review} 24 (3). In Korean.

\bibitem[{Kamphorst(forthcoming)}]{Kamphorst2025}
Kamphorst, J. Forthcoming. ``The Left Misunderstood: How Voters' Perceptions of What Matters to the Left Benefit the Right.'' \textit{Journal of Politics}.

\bibitem[{Kartik, Lee, and Suen(2021)}]{KartikLeeSuen2021}
Kartik, N., F. X. Lee, and W. Suen. 2021. ``Information Validates the Prior: A Theorem on Bayesian Updating and Applications.'' \textit{American Economic Review: Insights} 3 (2): 165--182.

\bibitem[{Kartik, Van~Weelden, and Wolton(2017)}]{KVW2017}
Kartik, N., R. Van Weelden, and S. Wolton. 2017. ``Electoral Ambiguity and Political Representation.'' \textit{American Journal of Political Science} 61 (4): 958--970.

\bibitem[{Kater(1989)}]{Kater1989}
Kater, M. H. 1989. \textit{Doctors under Hitler}. Chapel Hill: University of North Carolina Press.

\bibitem[{Kennedy et~al.(2018)}]{Kennedy2018}
Kennedy, C., M. Blumenthal, S. Clement, J. D. Clinton, C. Durand, C. Franklin, K. McGeeney, L. Miringoff, K. Olson, D. Rivers, L. Saad, G. E. Witt, and C. Wlezien. 2018. ``An Evaluation of the 2016 Election Polls in the United States.'' \textit{Public Opinion Quarterly} 82 (1): 1--33.

\bibitem[{Key(1955)}]{Key1955}
Key, V. O., Jr. 1955. ``A Theory of Critical Elections.'' \textit{Journal of Politics} 17 (1): 3--18.

\bibitem[{Kim(2022)}]{KimHanna2022}
Kim, H. 2022. ``Gender Cleavage and Vote Choice among Voters in Their Twenties: The 20th Presidential Election.'' \textit{Trends and Prospects} 115: 296--338. In Korean.

\bibitem[{King et~al.(2008)}]{King2008}
King, G., O. Rosen, M. Tanner, and A. F. Wagner. 2008. ``Ordinary Economic Voting Behavior in the Extraordinary Election of Adolf Hitler.'' \textit{Journal of Economic History} 68 (4): 951--996.

\bibitem[{Kornhauser(1959)}]{Kornhauser1959}
Kornhauser, W. 1959. \textit{The Politics of Mass Society}. Glencoe, IL: Free Press.

\bibitem[{Krasa and Polborn(2010)}]{KrasaPolborn2010}
Krasa, S., and M. Polborn. 2010. ``Competition between Specialized Candidates.'' \textit{American Political Science Review} 104 (4): 745--765.

\bibitem[{Kriesi et~al.(2008)}]{Kriesi2008}
Kriesi, H., E. Grande, R. Lachat, M. Dolezal, S. Bornschier, and T. Frey. 2008. \textit{West European Politics in the Age of Globalization}. Cambridge: Cambridge University Press.

\bibitem[{Kurer(2020)}]{Kurer2020}
Kurer, T. 2020. ``The Declining Middle: Occupational Change, Social Status, and the Populist Right.'' \textit{Comparative Political Studies} 53 (10--11): 1798--1835.

\bibitem[{Ladd and Lenz(2009)}]{LaddLenz2009}
Ladd, J. M., and G. S. Lenz. 2009. ``Exploiting a Rare Communication Shift to Document the Persuasive Power of the News Media.'' \textit{American Journal of Political Science} 53 (2): 394--410.

\bibitem[{La~Raja and Schaffner(2015)}]{LaRaja2015}
La Raja, R. J., and B. F. Schaffner. 2015. \textit{Campaign Finance and Political Polarization: When Purists Prevail}. Ann Arbor: University of Michigan Press.

\bibitem[{Lee(2016)}]{Lee2016}
Lee, F. E. 2016. \textit{Insecure Majorities: Congress and the Perpetual Campaign}. Chicago: University of Chicago Press.

\bibitem[{Levitsky and Ziblatt(2018)}]{LevitskyZiblatt2018}
Levitsky, S., and D. Ziblatt. 2018. \textit{How Democracies Die}. New York: Crown.

\bibitem[{Lijphart(1997)}]{Lijphart1997}
Lijphart, A. 1997. ``Unequal Participation: Democracy's Unresolved Dilemma.'' \textit{American Political Science Review} 91 (1): 1--14.

\bibitem[{Lindbeck and Weibull(1987)}]{Lindbeck1987}
Lindbeck, A., and J. W. Weibull. 1987. ``Balanced-Budget Redistribution as the Outcome of Political Competition.'' \textit{Public Choice} 52 (3): 273--297.

\bibitem[{Lipset(1960)}]{Lipset1960}
Lipset, S. M. 1960. \textit{Political Man: The Social Bases of Politics}. Garden City, NY: Doubleday.

\bibitem[{Loeper and Dziuda(2024)}]{LoeperDziuda2024}
Loeper, A., and W. Dziuda. 2024. ``Voters and the Trade-Off between Policy Stability and Responsiveness.'' \textit{Journal of Public Economics} 232: 105093.

\bibitem[{Lucas(1976)}]{Lucas1976}
Lucas, R. E., Jr. 1976. ``Econometric Policy Evaluation: A Critique.'' \textit{Carnegie-Rochester Conference Series on Public Policy} 1: 19--46.

\bibitem[{Luo and Przeworski(2023)}]{LuoPrzeworski2023}
Luo, Z., and A. Przeworski. 2023. ``Democracy and Its Vulnerabilities: Dynamics of Democratic Backsliding.'' \textit{Quarterly Journal of Political Science} 18 (1): 105--130.

\bibitem[{Macedo et~al.(2005)}]{Macedo2005}
Macedo, S., et al. 2005. \textit{Democracy at Risk: How Political Choices Undermine Citizen Participation, and What We Can Do About It}. Washington, DC: Brookings Institution Press.

\bibitem[{Mair(2013)}]{Mair2013}
Mair, P. 2013. \textit{Ruling the Void: The Hollowing of Western Democracy}. London: Verso.

\bibitem[{Margalit(2019)}]{Margalit2019}
Margalit, Y. 2019. ``Economic Insecurity and the Causes of Populism, Reconsidered.'' \textit{Journal of Economic Perspectives} 33 (4): 152--170.

\bibitem[{Martin and Yurukoglu(2017)}]{MartinYurukoglu2017}
Martin, G. J., and A. Yurukoglu. 2017. ``Bias in Cable News: Persuasion and Polarization.'' \textit{American Economic Review} 107 (9): 2565--2599.

\bibitem[{Marx(1852)}]{Marx1852}
Marx, K. 1852. ``Der achtzehnte Brumaire des Louis Napoleon.'' \textit{Die Revolution} 1. New York.

\bibitem[{Mason(2018)}]{Mason2018}
Mason, L. 2018. \textit{Uncivil Agreement: How Politics Became Our Identity}. Chicago: University of Chicago Press.

\bibitem[{Mayer(2002)}]{Mayer2002}
Mayer, N. 2002. \textit{Ces Fran\c{c}ais qui votent Le Pen}. Paris: Flammarion.

\bibitem[{MBC(2025)}]{MBC2025}
MBC. 2025. \textit{PD Su-cheop}, ``Conspiracy Theories in the Square, Part 1: Young People with Flags'' [in Korean]. 18 February. As described by producer Park Hae-in to \textit{OhmyStar}; quotation translated by the author.

\bibitem[{McCarty, Poole, and Rosenthal(2016)}]{McCarty2016}
McCarty, N., K. T. Poole, and H. Rosenthal. 2016. \textit{Polarized America: The Dance of Ideology and Unequal Riches}, 2nd ed. Cambridge, MA: MIT Press.

\bibitem[{McCoy(1999)}]{McCoy1999}
McCoy, J. 1999. ``Chávez and the End of `Partyarchy' in Venezuela.'' \textit{Journal of Democracy} 10 (3): 64--77.

\bibitem[{McCoy, Rahman, and Somer(2018)}]{McCoyRahmanSomer2018}
McCoy, J., T. Rahman, and M. Somer. 2018. ``Polarization and the Global Crisis of Democracy: Common Patterns, Dynamics, and Pernicious Consequences for Democratic Polities.'' \textit{American Behavioral Scientist} 62 (1): 16--42.

\bibitem[{McDonald(2020)}]{McDonald2020}
McDonald, M. P. 2020. ``2020 November General Election Turnout Rates.'' United States Elections Project.

\bibitem[{McKelvey(1976)}]{McKelvey1976}
McKelvey, R. D. 1976. ``Intransitivities in Multidimensional Voting Models and Some Implications for Agenda Control.'' \textit{Journal of Economic Theory} 12 (3): 472--482.

\bibitem[{McKelvey(1986)}]{McKelvey1986}
McKelvey, R. D. 1986. ``Covering, Dominance, and Institution-Free Properties of Social Choice.'' \textit{American Journal of Political Science} 30 (2): 283--314.

\bibitem[{Melumad and Shibano(1991)}]{MelumadShibano1991}
Melumad, N. D., and T. Shibano. 1991. ``Communication in Settings with No Transfers.'' \textit{RAND Journal of Economics} 22 (2): 173--198.

\bibitem[{Mian, Sufi, and Trebbi(2014)}]{Mian2014}
Mian, A., A. Sufi, and F. Trebbi. 2014. ``Resolving Debt Overhang: Political Constraints in the Aftermath of Financial Crises.'' \textit{American Economic Journal: Macroeconomics} 6 (2): 1--28.

\bibitem[{Miller and Schofield(2003)}]{MillerSchofield2003}
Miller, G., and N. Schofield. 2003. ``Activists and Partisan Realignment in the United States.'' \textit{American Political Science Review} 97 (2): 245--260.

\bibitem[{Mudde(2004)}]{Mudde2004}
Mudde, C. 2004. ``The Populist Zeitgeist.'' \textit{Government and Opposition} 39 (4): 541--563.

\bibitem[{Mutz(2018)}]{Mutz2018}
Mutz, D. C. 2018. ``Status Threat, Not Economic Hardship, Explains the 2016 Presidential Vote.'' \textit{Proceedings of the National Academy of Sciences} 115 (19): E4330--E4339.

\bibitem[{National Review(2016)}]{NationalReview2016}
National Review. 2016. ``Against Trump.'' Editorial and symposium, January 21.

\bibitem[{Neu and Pokorny(2025)}]{NeuPokorny2025}
Neu, V., and S. Pokorny. 2025. ``Bundestag Election in Germany on 23 February 2025: Final Results.'' Monitor Electoral and Social Research, Konrad-Adenauer-Stiftung.

\bibitem[{Nohlen and St\"{o}ver(2010)}]{NohlenStover2010}
Nohlen, D., and P. St\"{o}ver, eds. 2010. \textit{Elections in Europe: A Data Handbook}. Baden-Baden: Nomos.

\bibitem[{Norris and Inglehart(2019)}]{NorrisInglehart2019}
Norris, P., and R. Inglehart. 2019. \textit{Cultural Backlash: Trump, Brexit, and Authoritarian Populism}. Cambridge: Cambridge University Press.

\bibitem[{NSDAP(1920)}]{NSDAP1920}
Nationalsozialistische Deutsche Arbeiterpartei. 1920. \textit{Programm der Nationalsozialistischen Deutschen Arbeiterpartei}, 24 February 1920.

\bibitem[{OECD(2024)}]{OECDTrust2024}
OECD. 2024. \textit{OECD Survey on Drivers of Trust in Public Institutions---2024 Results: Building Trust in a Complex Policy Environment}. Paris: OECD Publishing. \url{https://doi.org/10.1787/9a20554b-en}.

\bibitem[{Osborne and Slivinski(1996)}]{OsborneSlivinski1996}
Osborne, M. J., and A. Slivinski. 1996. ``A Model of Political Competition with Citizen-Candidates.'' \textit{Quarterly Journal of Economics} 111 (1): 65--96.

\bibitem[{Patterson(2002)}]{Patterson2002}
Patterson, T. E. 2002. \textit{The Vanishing Voter: Public Involvement in an Age of Uncertainty}. New York: Knopf.

\bibitem[{PBS NewsHour(2024)}]{PBS2024}
PBS NewsHour. 2024. ``Americans Voters Explain Their Election Choices and Discuss Views on the Future.'' 8 November.

\bibitem[{Perrineau(1997)}]{Perrineau1997}
Perrineau, P. 1997. \textit{Le Sympt\^{o}me Le Pen: Radiographie des \'{e}lecteurs du Front national}. Paris: Fayard.

\bibitem[{Pew Research Center(2020)}]{Pew2020}
Pew Research Center. 2020. ``Election 2020: Voters Are Highly Engaged, but Nearly Half Expect to Have Difficulties Voting.'' Report, August 13.

\bibitem[{Pew Research Center(2025a)}]{Pew2025}
Pew Research Center. 2025a. ``Trump's Tariffs and `One Big Beautiful Bill' Face More Opposition Than Support as His Job Rating Slips.'' Report, August 14.

\bibitem[{Pew Research Center(2025b)}]{PewDemocracy2025}
Pew Research Center. 2025b. ``Dissatisfaction with Democracy Remains Widespread in Many Nations.'' June 30. \url{https://www.pewresearch.org/short-reads/2025/06/30/dissatisfaction-with-democracy-remains-widespread-in-many-nations/}.

\bibitem[{Pew Research Center(2025c)}]{PewTrust2025}
Pew Research Center. 2025c. ``Public Trust in Government: 1958--2025.'' December 4. \url{https://www.pewresearch.org/politics/2025/12/04/public-trust-in-government-1958-2025/}.

\bibitem[{Pew Research Center(2026)}]{Pew2026}
Pew Research Center. 2026. ``Americans Largely Disapprove of Trump's Tariff Increases.'' Report, February 4.

\bibitem[{Plott(1967)}]{Plott1967}
Plott, C. R. 1967. ``A Notion of Equilibrium and Its Possibility Under Majority Rule.'' \textit{American Economic Review} 57 (4): 787--806.

\bibitem[{Price(2004)}]{Price2004}
Price, R. 2004. \textit{People and Politics in France, 1848--1870}. Cambridge: Cambridge University Press.

\bibitem[{Prior(2007)}]{Prior2007}
Prior, M. 2007. \textit{Post-Broadcast Democracy: How Media Choice Increases Inequality in Political Involvement and Polarizes Elections}. Cambridge: Cambridge University Press.

\bibitem[{Przeworski(1991)}]{Przeworski1991}
Przeworski, A. 1991. \textit{Democracy and the Market: Political and Economic Reforms in Eastern Europe and Latin America}. Cambridge: Cambridge University Press.

\bibitem[{Przeworski(2019)}]{Przeworski2019}
Przeworski, A. 2019. \textit{Crises of Democracy}. Cambridge: Cambridge University Press.

\bibitem[{Przeworski(2023)}]{Przeworski2023}
Przeworski, A. 2023. ``Formal Models of Authoritarian Regimes: A Critique.'' \textit{Perspectives on Politics} 21 (3): 979--988.

\bibitem[{Przeworski and Sprague(1986)}]{PrzeworskiSprague1986}
Przeworski, A., and J. Sprague. 1986. \textit{Paper Stones: A History of Electoral Socialism}. Chicago: University of Chicago Press.

\bibitem[{Pulzer(1988)}]{Pulzer1988}
Pulzer, P. 1988. \textit{The Rise of Political Anti-Semitism in Germany and Austria}. Revised edition. Cambridge, MA: Harvard University Press.

\bibitem[{Putnam(2000)}]{Putnam2000}
Putnam, R. D. 2000. \textit{Bowling Alone: The Collapse and Revival of American Community}. New York: Simon \& Schuster.

\bibitem[{Quinnipiac University Poll(2018)}]{Quinnipiac2018}
Quinnipiac University Poll. 2018. ``U.S. Voters Oppose Family Separation.'' National poll released June 18, 2018.

\bibitem[{Riker and Ordeshook(1968)}]{RikerOrdeshook1968}
Riker, W. H., and P. C. Ordeshook. 1968. ``A Theory of the Calculus of Voting.'' \textit{American Political Science Review} 62 (1): 25--42.

\bibitem[{Rock(1987)}]{Rock1987}
Rock, D. 1987. \textit{Argentina, 1516--1987: From Spanish Colonization to Alfons\'{i}n}. Berkeley: University of California Press.

\bibitem[{Roemer(1994)}]{Roemer1994}
Roemer, J. E. 1994. ``The Strategic Role of Party Ideology When Voters Are Uncertain about How the Economy Works.'' \textit{American Political Science Review} 88 (2): 327--335.

\bibitem[{Roemer(2001)}]{Roemer2001}
Roemer, J. E. 2001. \textit{Political Competition: Theory and Applications}. Cambridge, MA: Harvard University Press.

\bibitem[{Rogoff(1985)}]{Rogoff1985}
Rogoff, K. 1985. ``The Optimal Degree of Commitment to an Intermediate Monetary Target.'' \textit{Quarterly Journal of Economics} 100 (4): 1169--1189.

\bibitem[{Rooduijn et~al.(2024)}]{Rooduijn2024}
Rooduijn, M., A. L. P. Pirro, D. Halikiopoulou, et al. 2024. ``The PopuList: A Database of Populist, Far-Left, and Far-Right Parties Using Expert-Informed Qualitative Comparative Classification (EiQCC).'' \textit{British Journal of Political Science}.

\bibitem[{Rottinghaus(2025)}]{Rottinghaus2025}
Rottinghaus, B. 2025. \textit{Scandal: Why Politicians Survive Controversy in a Partisan Era}. New York: Columbia University Press.

\bibitem[{SBS(2025)}]{SBS2025}
SBS. 2025. \textit{Newstory}, ``The 2030 Generation Shouting `Against Impeachment' at Rallies'' [in Korean]. March. Quotation translated by the author.

\bibitem[{Schmidt(2024)}]{Schmidt2024}
Schmidt, G. 2024. ``Climate Models Can't Explain 2023's Huge Heat Anomaly---We Could Be in Uncharted Territory.'' \textit{Nature} 627: 467.

\bibitem[{Schofield(2007)}]{Schofield2007}
Schofield, N. 2007. ``The Mean Voter Theorem: Necessary and Sufficient Conditions for Convergent Equilibrium.'' \textit{Review of Economic Studies} 74 (3): 965--980.

\bibitem[{Schultz(1996)}]{Schultz1996}
Schultz, C. 1996. ``Polarization and Inefficient Policies.'' \textit{Review of Economic Studies} 63 (2): 331--344.

\bibitem[{Scott(1998)}]{Scott1998}
Scott, J. C. 1998. \textit{Seeing Like a State: How Certain Schemes to Improve the Human Condition Have Failed}. New Haven: Yale University Press.

\bibitem[{Seawright(2012)}]{Seawright2012}
Seawright, J. 2012. \textit{Party-System Collapse: The Roots of Crisis in Peru and Venezuela}. Stanford: Stanford University Press.

\bibitem[{Skelley(2017)}]{Skelley2017}
Skelley, G. 2017. ``Just How Many Obama 2012--Trump 2016 Voters Were There?'' \textit{Sabato's Crystal Ball}, University of Virginia Center for Politics, June 1.

\bibitem[{Sniderman, Hagendoorn, and Prior(2004)}]{SnidermanHagendoornPrior2004}
Sniderman, P. M., L. Hagendoorn, and M. Prior. 2004. ``Predisposing Factors and Situational Triggers: Exclusionary Reactions to Immigrant Minorities.'' \textit{American Political Science Review} 98 (1): 35--49.

\bibitem[{Snyder and Str\"{o}mberg(2010)}]{SnyderStromberg2010}
Snyder, J. M., Jr., and D. Str\"{o}mberg. 2010. ``Press Coverage and Political Accountability.'' \textit{Journal of Political Economy} 118 (2): 355--408.

\bibitem[{Spenkuch and Tillmann(2018)}]{SpenkuchTillmann2018}
Spenkuch, J. L., and P. Tillmann. 2018. ``Elite Influence? Religion and the Electoral Success of the Nazis.'' \textit{American Journal of Political Science} 62 (1): 19--36.

\bibitem[{Svolik(2019)}]{Svolik2019}
Svolik, M. W. 2019. ``Polarization versus Democracy.'' \textit{Journal of Democracy} 30 (3): 20--32.

\bibitem[{Teperoglou and Tsatsanis(2014)}]{TeperoglouTsatsanis2014}
Teperoglou, E., and E. Tsatsanis. 2014. ``Dealignment, De-legitimation and the Implosion of the Two-Party System in Greece: The Earthquake Election of 6 May 2012.'' \textit{Journal of Elections, Public Opinion and Parties} 24 (2): 222--242.

\bibitem[{The Independent(2026)}]{Independent2026}
The Independent. 2026. ``Local Election in Numbers: Inside Reform's Back-to-Back Record Results and Labour's Collapse.'' May 10.

\bibitem[{Thomsen(2017)}]{Thomsen2017}
Thomsen, D. M. 2017. \textit{Opting Out of Congress: Partisan Polarization and the Decline of Moderate Candidates}. Cambridge: Cambridge University Press.

\bibitem[{Tocqueville(2016)}]{Tocqueville2016}
Tocqueville, A. de. 2016. \textit{Recollections: The French Revolution of 1848 and Its Aftermath}. Edited by O. Zunz. Translated by A. Goldhammer. Charlottesville: University of Virginia Press.

\bibitem[{Truex and Arceneaux(2026)}]{TruexArceneaux2026}
Truex, R., and K. Arceneaux. 2026. ``Election Legitimacy Tracking Survey.'' Data dashboard, accessed November 2026.

\bibitem[{UNDP(2022)}]{UNDP2022}
United Nations Development Programme. 2022. \textit{Human Development Report 2021/2022: Uncertain Times, Unsettled Lives}. New York: UNDP.

\bibitem[{Vaeth(2026)}]{Vaeth2024}
Vaeth, M. 2026. ``Rational Voter Learning, Issue Alignment, and Polarization.'' Working paper, version of August 31, 2026.

\bibitem[{Wattenberg(2002)}]{Wattenberg2002}
Wattenberg, M. P. 2002. \textit{Where Have All the Voters Gone?} Cambridge, MA: Harvard University Press.

\bibitem[{Wei{\ss}(1932)}]{Weiss1932}
Wei{\ss}, B. 1932. ``Mehr Selbstbewu{\ss}tsein.'' \textit{C.V.-Zeitung} 11 (23), June 3. Translated as ``More Self-Confidence'' in \textit{German History in Documents and Images}, German Historical Institute, Washington, DC.

\bibitem[{Weyland(2003)}]{Weyland2003}
Weyland, K. 2003. ``Economic Voting Reconsidered: Crisis and Charisma in the Election of Hugo Ch\'{a}vez.'' \textit{Comparative Political Studies} 36 (7): 822--848.

\bibitem[{Wittman(1983)}]{Wittman1983}
Wittman, D. 1983. ``Candidate Motivation: A Synthesis of Alternative Theories.'' \textit{American Political Science Review} 77 (1): 142--157.

\bibitem[{Working Group on Universal Voting(2020)}]{LiftEveryVoice2020}
Working Group on Universal Voting. 2020. \textit{Lift Every Voice: The Urgency of Universal Civic Duty Voting}. Washington, DC: Brookings Institution and Ash Center for Democratic Governance and Innovation.

\bibitem[{World Meteorological Organization(2025)}]{WMO2025}
World Meteorological Organization. 2025. \textit{State of the Global Climate 2024}. Geneva: WMO.

\bibitem[{Zhuravskaya, Petrova, and Enikolopov(2020)}]{Zhuravskaya2020}
Zhuravskaya, E., M. Petrova, and R. Enikolopov. 2020. ``Political Effects of the Internet and Social Media.'' \textit{Annual Review of Economics} 12: 415--438.

\bibitem[{Zhuravskaya, Petrova, and Enikolopov(2020)}]{ZhuravskayaPetrovaEnikolopov2020}
Zhuravskaya, E., M. Petrova, and R. Enikolopov. 2020. ``Political Effects of the Internet and Social Media.'' \textit{Annual Review of Economics} 12: 415--438.

\end{thebibliography}

\clearpage
\appendix
\setcounter{page}{1}\renewcommand{\thepage}{A\arabic{page}}
\begin{center}
{\Large Supplementary Information for}\\[0.8ex]
{\Large ``Realignment and Polarization in an Unpredictable World''}
\end{center}
\vspace{1em}
\phantomsection\label{app:proofs}
\setcounter{section}{0}\renewcommand{\thesection}{A\arabic{section}}
\providecommand{\theHsection}{}\renewcommand{\theHsection}{A\arabic{section}}
\addtocontents{toc}{\protect\setcounter{tocdepth}{1}}
\makeatletter\renewcommand*\l@section{\@dottedtocline{1}{0em}{2.6em}}\makeatother
\begingroup\renewcommand{\contentsname}{Contents}\tableofcontents\endgroup
\clearpage
\setcounter{equation}{0}\renewcommand{\theequation}{A.\arabic{equation}}
\setcounter{lemma}{0}\renewcommand{\thelemma}{A.\arabic{lemma}}
\setcounter{proposition}{0}\renewcommand{\theproposition}{A.\arabic{proposition}}
\setcounter{corollary}{0}\renewcommand{\thecorollary}{A.\arabic{corollary}}
\setcounter{assumption}{0}\renewcommand{\theassumption}{A.\arabic{assumption}}
\setcounter{table}{0}\renewcommand{\thetable}{A\arabic{table}}\providecommand{\theHtable}{}\renewcommand{\theHtable}{A\arabic{table}}
\setcounter{figure}{0}\renewcommand{\thefigure}{A\arabic{figure}}\providecommand{\theHfigure}{}\renewcommand{\theHfigure}{A\arabic{figure}}
% Keep appendix hyperlinks distinct from main-text theorem counters.
\providecommand{\theHlemma}{}
\renewcommand{\theHlemma}{A.\arabic{lemma}}
\providecommand{\theHproposition}{}
\renewcommand{\theHproposition}{A.\arabic{proposition}}
\providecommand{\theHcorollary}{}
\renewcommand{\theHcorollary}{A.\arabic{corollary}}
\providecommand{\theHassumption}{}
\renewcommand{\theHassumption}{A.\arabic{assumption}}


The supplement collects the proofs of the results in the text and the arguments for their generality. Its sections are numbered A1, A2, and so on, and formal statements and equations introduced here are numbered A.1, A.2, and so on. Unless stated otherwise, the parties receive the same fixed office benefit $B\ge0$. Bounds under moderate electoral sensitivity are uniform over any fixed compact interval of benefits; they do not cover benefits that grow like $1/\phi$, as in the strategic extension.

\section{Notation}\label{app:notation}
Table~\ref{tab:notation-si} lists all notation that does not appear in Table~\ref{tab:notation}, including equilibrium quantities and notation used only in this supplement.

\begin{table}[ht]
\centering
\scriptsize
\renewcommand{\arraystretch}{1.15}
\setlength{\aboverulesep}{0pt}\setlength{\belowrulesep}{0pt}\setlength{\extrarowheight}{2pt}
\begin{adjustbox}{max width=\textwidth,center}
\begin{tabular}{@{}L{1.9cm}L{3.2cm}|L{1.9cm}L{3.2cm}|L{2.2cm}L{3.2cm}@{}}
\toprule
{\bfseries\boldmath Symbol} & {\bfseries\boldmath Meaning} & {\bfseries\boldmath Symbol} & {\bfseries\boldmath Meaning} & {\bfseries\boldmath Symbol} & {\bfseries\boldmath Meaning} \\
\hhline{==|==|==}
\multicolumn{2}{l|}{\textbf{Voting and realignment}} & \multicolumn{2}{l|}{\textbf{Information}} & \multicolumn{2}{l}{\textbf{Several states}}  \\
\hhline{==|==|==}
$x$; $x_j^*(\mu)$ & Policy; $j$'s policy & $s$; $\varepsilon$ & Signal; its noise & $N$; $\omega_n$ & Number of states; state $n$ \\
\hhline{--|--|--}
$m$ & Posterior mean & $\sigma_\omega^2$; $\sigma_s^2$ & Variances in the normal example & $\alpha_{in}$; $S_n$ & Responsiveness and $S$ for state $n$ \\
\hhline{--|--|--}
$\psi$; $\psi_\eta$; $\psi_M$ & Indifference line; shifted; median line & $\Lambda_i^j(\omega)$ & Cost of delegation given the state & $\Lambda_{i,-n}^j$ & Cost of delegation without state $n$ \\
\hhline{--|--|--}
$\gamma$ & Minus the slope of $\psi$ & $\rho$; $\zeta$ & Precision; cost scale & & \\
\hhline{--|--|--}
$\theta_M$ & Median point & $S(\rho)$; $\rho_j^*$; $S_j^*$ & Learning at $\rho$; $j$'s choices & \multicolumn{2}{l}{\textbf{Only in the SI}}  \\
\hhline{--|--|==}
$\overline{M}$; $\overline{\phi}$ & Bounds in Assumption~\ref{ass:small-h} & $s_j$ & Signal that candidate $j$ would observe & $O(\phi)$, $o(\phi)$; $\tilde o$ & Orders of $\phi$; constant \\
\hhline{--|--|--}
\multicolumn{2}{l|}{\textbf{Equilibrium}} & $s^1$; $s^2$ & Early news; later signal & $\hat \phi$; $\underline S$, $\overline S$ & Bounds on $\phi$ and $S$  \\
\hhline{==|--|--}
$\overline{\theta}$ & Decisive average type & $m^1$; $S^2(s^1)$ & Mean after news; learning left & $\nu\cdot\nu'$; $H$ & Policy inner product; its matrix \\
\hhline{--|--|--}
$\theta_l^*$, $\theta_r^*$ & Equilibrium nominees & $\Lambda(s^1)$; $D^*(\Lambda)$ & Disagreement given news; equilibrium policy gap & $\overline{\theta}_0$ & $\overline{\theta}$ at ideal types \\
\hhline{--|--|--}
$w_J$; $w$ & Share of $J$'s ideal kept; symmetric value & $\mathcal H$, $\mathcal H'$; $\mathcal F$ & Early; final information & $\mathcal N$; $\tilde\theta_l$ & Neighbourhood; maximizer \\
\hhline{--|--|--}
$\Lambda$; $\Lambda_M^J$ & $\Lambda_L^R$; $J$'s gap to $\theta_M$ & \multicolumn{2}{l|}{\textbf{Strategic uncertainty}} & $\iota$; $\Xi$; $\xi$; $y$ & Noise; cdf; density; argument  \\
\hhline{--|==|--}
$D$; $D^*$ & Policy gap; equilibrium value & $\theta_{\mathcal I}$; $\theta_{\mathcal C}$; $\theta_c$ & Incumbent; challenger ideal; nominee & & \\
\hhline{--|--|--}
$\Delta U_J^*$ & Equilibrium payoff gap & $\Pi_{\mathcal I}$ & Incumbent's payoff & $\widehat R^2$; $\widehat\alpha_r$, $\widehat \rho$, $\widehat D$, $\widehat\theta_J$ & Revealed share at the chosen precision; limits in Corollary~\ref{cor:endolearn} \\
\hhline{--|--|--}
$S^\dagger$ & $S$ at which the lead switches & $B_{\mathcal I}$; $B_{\mathcal C}$ & Incumbent's and challenger's benefits & $\Lambda(\mathcal F)$ & Disagreement given final information \\
\hhline{--|--|--}
$U_i^*$ & Equilibrium expected policy payoff & $\overline{\alpha}_c$; $w_c$ & Nominee's limiting $\alpha$; share kept & $i$; $j$; $J$; $n$ & Player; candidate; party; state \\
\hhline{--|--|--}
$\delta$ & Half the parties' gap in $\alpha$ & $\Delta U_{\mathcal C}$ & Challenger's payoff gap & $\varepsilon_i$; $F_\varepsilon$ & Idiosyncratic noise; its cdf \\
\hhline{--|--|--}
$\mathcal S_l(\eta)$ & Conditional vote share & $\overline{B}$; $\underline B$ & Thresholds for $\phi B_{\mathcal I}$ & $F_0^1$; $F_0^2$ & Prior distributions compared in Prop.~\ref{prop:dispersion} \\
\bottomrule
\end{tabular}
\end{adjustbox}
\caption{Notation not listed in Table~\ref{tab:notation}.}
\label{tab:notation-si}
\end{table}
For a quantity $a$ that depends on $\phi$, write $a=O(\phi)$ if there exist constants $\tilde o>0$ and $\hat \phi>0$ such that $|a|\le\tilde o \phi$ for all $\phi\le\hat \phi$, uniformly over candidates in the neighbourhood of the relevant profile, office benefits in a fixed compact interval, and unpredictability in a compact interval $[\underline S,\overline S]$ with $\underline S>0$. The constants depend on the primitives, on that neighbourhood and on that interval; they are not uniform as unpredictability approaches zero. A statement of the form ``$a=b+O(\phi)$, so $a>0$ for $\phi$ small enough'' with $b>0$ therefore means that there are such $\tilde o$ and $\hat \phi$ with $|a-b|\le\tilde o \phi$ for all $\phi\le\hat \phi$, so that $a>0$ whenever $\phi<\min\{\hat \phi,\,b/\tilde o\}$.

\section{Requirements of each result}\label{app:requirements}
Table~\ref{tab:general} lists, for each result, what it requires of the distribution of voters, the parties' ideal types, the targets and electoral sensitivity, and the results it builds on.

\paragraph{How the results connect.}
Every result starts from Lemma~\ref{lem:payoff}, which writes a player's payoff from a candidate in terms of their gaps in baseline and responsiveness and of $S$. With candidates fixed, it yields the rotating line of indifferent voters (Proposition~\ref{prop:realign}) and the cost of delegation (Proposition~\ref{prop:office}); the cost of delegation in turn gives the value of office (Corollary~\ref{cor:office}) and the polarization of a fixed pair of candidates (Proposition~\ref{prop:auto-polar}). Section~\ref{sec:competition} adds elections: Lemmas~\ref{lem:median} and~\ref{lem:win} turn voters' comparisons into a winning probability, Corollary~\ref{cor:median} places the decisive voter at the median point, and Proposition~\ref{prop:nomination} weighs that probability against the value of office to characterize nominations. Section~\ref{sec:effects} applies Proposition~\ref{prop:nomination}: Propositions~\ref{prop:symmetric} and~\ref{prop:asymmetric} show that nominations offset Proposition~\ref{prop:auto-polar} without undoing it, Propositions~\ref{prop:welfare-unc} and~\ref{prop:welfare} evaluate the resulting lottery between candidates with Proposition~\ref{prop:office}, and Corollaries~\ref{cor:endolearn} and~\ref{cor:timing} let information be chosen or arrive early. Section~\ref{sec:pre-election} lets the incumbent choose uncertainty, using Corollary~\ref{cor:median} for the margin and the nomination condition of Proposition~\ref{prop:nomination} for the challenger. The last column of Table~\ref{tab:general} lists each result's direct inputs.

\begin{table}[p]
\begin{spacing}{1}
\centering
\footnotesize
\renewcommand{\arraystretch}{1.1}
\begin{adjustbox}{max totalsize={\textwidth}{0.83\textheight},center}
\begin{tabular}{@{}L{3.6cm}L{2.6cm}L{2.2cm}L{2.5cm}L{1.9cm}L{3.2cm}@{}}
\toprule
{\bfseries\boldmath Result} & {\bfseries\boldmath Voters $G$} & {\bfseries\boldmath Party ideals} & {\bfseries\boldmath Targets} & {\bfseries\boldmath $\phi$} & {\bfseries\boldmath Builds on} \\
\hline\hline
Prop.~\ref{prop:auto-polar} & --- & Any & Any (information); convex gap under full information (uncertainty) & --- & Lemma~\ref{lem:payoff} \\ \midrule
Cor.~\ref{cor:office} & --- & Any & Affine & --- & Lemma~\ref{lem:payoff} \\ \midrule
Prop.~\ref{prop:realign} & --- & Any & (a) affine; (b) any (information), convex gap under full information (uncertainty) & --- & Lemma~\ref{lem:payoff} \\ \midrule
Lemma~\ref{lem:median} & Any & Any & Affine & --- & --- \\ \midrule
Lemma~\ref{lem:win} & Any & Any & Affine & Ass.~\ref{ass:small-h} & Lemma~\ref{lem:median} \\ \midrule
Cor.~\ref{cor:median} & MP & Any & Affine & --- & Lemmas~\ref{lem:payoff}, \ref{lem:median} \\ \midrule
Prop.~\ref{prop:nomination} & Any; $M$ differentiable & Any & Affine & Ass.~\ref{ass:small-h} & Lemmas~\ref{lem:payoff}--\ref{lem:win}, Cor.~\ref{cor:median} \\ \midrule
Prop.~\ref{prop:symmetric} & MP & Symmetric & Affine & Any & Prop.~\ref{prop:nomination}, Lemma~\ref{lem:global} \\ \midrule
Prop.~\ref{prop:asymmetric}(a) & MP; LR (Cor.~\ref{cor:leadgeneral}) & Distinct & Affine & Small & Prop.~\ref{prop:nomination}, Cor.~\ref{cor:median}, Lemma~\ref{lem:smallh} \\ \midrule
Prop.~\ref{prop:asymmetric}(b)--(d) & MP & Distinct & Affine & Small & Prop.~\ref{prop:asymmetric}(a), Prop.~\ref{prop:general} \\ \midrule
Cor.~\ref{cor:benefits}, Prop.~\ref{prop:general} & LR & Distinct & Affine & Small & Prop.~\ref{prop:nomination}, Lemma~\ref{lem:smallh} \\ \midrule
Prop.~\ref{prop:welfare-unc} & (a) ---; (b) LR; (c) MP & (b) distinct; (c) symmetric & Affine & (b) small; (c) any & Lemma~\ref{lem:payoff}; Prop.~\ref{prop:general}; Prop.~\ref{prop:welfare} \\ \midrule
Prop.~\ref{prop:welfare} & MP & Symmetric & Affine & Any & Props.~\ref{prop:office}, \ref{prop:symmetric} \\ \midrule
Cor.~\ref{cor:endolearn} & (c) smooth margins & (c) $\alpha_R=-\alpha_L$ & Normal signals, quadratic cost & (c) small & Lemma~\ref{lem:payoff}, Lemma~\ref{lem:smallh} \\ \midrule
Cor.~\ref{cor:timing} & MP & Symmetric & Affine & Any & Prop.~\ref{prop:auto-polar}, Prop.~\ref{prop:symmetric}, Lemma~\ref{lem:global} \\ \midrule
Prop.~\ref{prop:strategic} & MP & Distinct & Affine & Small, $B_{\mathcal I},B_{\mathcal C}\propto1/\phi$ & Cor.~\ref{cor:median}, Prop.~\ref{prop:nomination} \\ \midrule
Prop.~\ref{prop:multistate} & (c) MP for hyperplanes & As originally & Affine; correlated states allowed; (c) $S\succ0$ & As originally & Lemma~\ref{lem:payoff} and the original results \\ \midrule
Prop.~\ref{prop:issues} & (a), (c), (d) MP; (b) LR & Distinct & Affine, independent issues & (a)--(c) small & Props.~\ref{prop:asymmetric}, \ref{prop:strategic}, \ref{prop:welfare-unc}, \ref{prop:welfare} \\ \midrule
Prop.~\ref{prop:dispersion} & --- & --- & (a), (b) squared gap single-peaked at the prior; full information & --- & Lemma~\ref{lem:targets} \\
\bottomrule
\end{tabular}
\end{adjustbox}
\caption{What each result requires and the results it builds on. ``---'': not used; MP: median point (Assumption~\ref{ass:median}); LR: local regularity (Assumption~\ref{ass:local}); ``Distinct'': distinct, interior ideals; ``Small'': all sufficiently small $\phi>0$, with unpredictability bounded away from zero; ``Any'': every $\phi>0$. Unlisted parts hold without the condition. Section~\ref{sec:choices} discusses log-concave noise (Proposition~\ref{prop:logconcave}).}
\label{tab:general}
\end{spacing}
\end{table}

\section{Proofs for Sections~\ref{sec:policy} and~\ref{sec:realign}}

\begin{proof}[Proof of Lemma~\ref{lem:payoff}]
Write $m=\mathbb{E}_\mu[\omega]$. Then $x^*_j(\mu)=\beta_j+\alpha_j m$, $\mathbb{E}[m]=\mathbb{E}[\omega]=0$, and, by the law of iterated expectations,
\begin{equation}\label{eq:moments}
\mathbb{E}[m\,\omega]=\mathbb{E}[m^2]=S,\qquad \mathbb{E}[\omega^2]=\operatorname{Var}(\omega).
\end{equation}
Hence
\[
U_i^j=-\mathbb{E}\big[(\beta_j-\beta_i+\alpha_j m-\alpha_i\omega)^2\big]
=-(\beta_j-\beta_i)^2-\alpha_j^2S+2\alpha_i\alpha_jS-\alpha_i^2\operatorname{Var}(\omega),
\]
which is the first expression. Subtracting the expressions for $j=r$ and $j=l$,
\[
\Delta U_i=(\beta_r-\beta_i)^2-(\beta_l-\beta_i)^2+\big[\alpha_l(2\alpha_i-\alpha_l)-\alpha_r(2\alpha_i-\alpha_r)\big]S ,
\]
and $(\beta_r-\beta_i)^2-(\beta_l-\beta_i)^2=2\Delta\beta(\beta_i-\beta_m)$ and $\alpha_l(2\alpha_i-\alpha_l)-\alpha_r(2\alpha_i-\alpha_r)=2\Delta\alpha(\alpha_i-\alpha_m)$. If $\beta_r>\beta_l$, then $\Delta\beta<0$, and dividing $\Delta U_i\ge\eta$ by $2\Delta\beta$ gives $\beta_i-\beta_m+\gamma(\alpha_i-\alpha_m)\le\eta/(2\Delta\beta)$, that is, $\beta_i\le\psi_\eta(\alpha_i)$.
\end{proof}

\begin{proof}[Proof of Proposition~\ref{prop:auto-polar}]
With $m=\mathbb{E}_\mu[\omega]$, $x^*_l-x^*_r=\Delta\beta+\Delta\alpha\,m$, so by \eqref{eq:moments}, $D=(\Delta\beta)^2+(\Delta\alpha)^2S$, which increases in $S$ when $\Delta\alpha\ne0$. The difference between the expected policies is $\mathbb{E}[x^*_l-x^*_r]=\Delta\beta+\Delta\alpha\,\mathbb{E}[m]=\Delta\beta$.
\end{proof}

\begin{proof}[Proof of Corollary~\ref{cor:office}]
By Lemma~\ref{lem:payoff}, $U_L^j=-\alpha_L^2\operatorname{Var}(\omega)-(\beta_j-\beta_L)^2+\big[\alpha_L^2-(\alpha_j-\alpha_L)^2\big]S$, so $V_L=B+\Delta U_L=B+(\beta_r-\beta_L)^2-(\beta_l-\beta_L)^2+\big[(\alpha_r-\alpha_L)^2-(\alpha_l-\alpha_L)^2\big]S$. Its derivative with respect to $S$ is positive if and only if $|\alpha_l-\alpha_L|<|\alpha_r-\alpha_L|$.
\end{proof}

\begin{proof}[Proof of Proposition~\ref{prop:realign}]
(a)
By Lemma~\ref{lem:payoff}, voter $i$ prefers $l$ if and only if $\Delta U_i>0$. Dividing by $2\Delta\beta<0$, this is $\beta_i-\beta_m+\gamma(\alpha_i-\alpha_m)<0$, that is, $\beta_i<\psi(\alpha_i)$. The line $\psi$ passes through $\theta_m$ for every $S$, and its slope $-\gamma$ is proportional to $S$, so it rotates around $\theta_m$ as $S$ rises; a rise in uncertainty or in the revealed share raises $S$. Let $\alpha_r>\alpha_l$, so $\gamma>0$. A voter with $\beta_i<\beta_m$ and $\alpha_i>\alpha_m$ prefers $l$ when $S=0$; the term $\gamma(\alpha_i-\alpha_m)$ is positive and increasing in $S$, so she prefers $r$ once it exceeds $\beta_m-\beta_i$, and never switches back. The argument for voters with $\beta_i>\beta_m$ and $\alpha_i<\alpha_m$ is symmetric, and the preference of any other voter does not change sign. Since $\Delta U_i$ is linear in $S$, no voter switches more than once.

(b) With $m=\mathbb{E}_\mu[\omega]$, $x^*_j-x^*_i=\beta_j-\beta_i+(\alpha_j-\alpha_i)m$, so by \eqref{eq:moments}, $\Lambda_i^j=(\beta_j-\beta_i)^2+(\alpha_j-\alpha_i)^2S$. Lemma~\ref{lem:payoff} with $j=i$ gives $U_i^i=-\alpha_i^2\big(\operatorname{Var}(\omega)-S\big)$, and
\[
U_i^j-U_i^i=-(\beta_j-\beta_i)^2+\big[\alpha_j(2\alpha_i-\alpha_j)-\alpha_i^2\big]S=-\Lambda_i^j.
\]
\end{proof}

\begin{proof}[Proof of Lemma~\ref{lem:median}]
The line $\psi_\eta$ has slope $-\gamma$, which does not depend on $\eta$, and a height that is strictly decreasing and affine in $\eta$. Exactly one median line has this slope, so exactly one $\eta$, which is $M(\theta_l,\theta_r)$, makes $\psi_\eta$ that line. The share of voters voting for $l$, $G(\{\theta_i:\beta_i\le\psi_\eta(\alpha_i)\})$, is decreasing in $\eta$, so it exceeds $\tfrac12$ if and only if $\eta<M(\theta_l,\theta_r)$. Relabelling the candidates covers $\beta_r<\beta_l$. If $\beta_r=\beta_l$ but $(\alpha_r-\alpha_l)S\ne0$, the lines are vertical and the same argument applies with $\alpha$ and $\beta$ exchanged. If both differences relevant to policy are zero, all voters are indifferent and $M\equiv0$.
\end{proof}

\begin{proof}[Proof of Lemma~\ref{lem:win}]
By Lemma~\ref{lem:median}, $l$ wins if and only if $\eta<M(\theta_l,\theta_r)$. By Assumption~\ref{ass:small-h}, $M(\theta_l,\theta_r)$ lies in the support of $\eta$, which is uniform with density $\phi$, so $\Pr(\eta<M)=\tfrac12+\phi M$. Substituting into $\Pi_J$ gives the expression for the payoff.
\end{proof}

\begin{proof}[Proof of Corollary~\ref{cor:median}]
The halving line $\psi_M$ has slope $-\gamma$ and, by Assumption~\ref{ass:median}, passes through $\theta_M$, so $\psi_M(\alpha_M)=\beta_M$. Since $\psi_M=\psi-M/\big(2(\beta_r-\beta_l)\big)$, this gives $M=2(\beta_r-\beta_l)\big(\psi(\alpha_M)-\beta_M\big)$, which equals $\Delta U_{\theta_M}$ by Lemma~\ref{lem:payoff}.
\end{proof}

\begin{proof}[Proof of Remark~\ref{rem:hardening}]
By Lemma~\ref{lem:payoff}, $\Delta U_i=\Delta U_i^0+\Delta U_i^{\prime}S$ with $\Delta U_i^0=2\Delta\beta(\beta_i-\beta_m)$ and $\Delta U_i^{\prime}=2\Delta\alpha(\alpha_i-\alpha_m)$. If $\Delta U_i^0$ and $\Delta U_i^{\prime}$ have the same sign, $|\Delta U_i^0+\Delta U_i^{\prime}S|$ increases in $S$. If they have opposite signs, $|\Delta U_i^0+\Delta U_i^{\prime}S|$ decreases until $S=-\Delta U_i^0/\Delta U_i^{\prime}$, where the preference switches, and increases afterward, which proves (a). For (b), because $\Delta U_i$ is affine in $\theta_i$, the coarea formula \citep{EvansGariepy2015} gives the mass of voters with $\Delta U_i\in[t_1,t_2]$ as the integral over $t\in[t_1,t_2]$ of the $g$-weighted length of the line $\{\Delta U_i=t\}$, divided by $\|\nabla\Delta U\|$. The weighted length is at most $\sup g$ times the diameter of $\operatorname{supp}(G)$, and $\|\nabla\Delta U\|\ge2|\Delta\alpha|S$. Taking $[t_1,t_2]=[-\tfrac{1}{2\phi},\tfrac{1}{2\phi}]$ bounds the share of voters whose vote depends on the noise. A common shift $v$ acts like a shift of the noise by $-v$, and taking an interval of length $|v|$ bounds the share of voters whom it moves. For (c), the shift moves the tie value $M$ by $v$, so by Lemma~\ref{lem:win} the winning probability $\tfrac12+\phi M$ changes by $\phi v$, provided that the new tie value stays inside the support of the noise. For (d), $\operatorname{Var}_G(\Delta U_i^0+\Delta U_i^{\prime}S)=\operatorname{Var}_G(\Delta U_i^0)+2S\operatorname{Cov}_G(\Delta U_i^0,\Delta U_i^{\prime})+S^2\operatorname{Var}_G(\Delta U_i^{\prime})$, and $\operatorname{Var}_G(\Delta U_i^{\prime})=4(\Delta\alpha)^2\operatorname{Var}_G(\alpha_i)>0$.
\end{proof}

\begin{proof}[Proof of Remark~\ref{rem:responses}]
By Lemma~\ref{lem:payoff}, $\Delta U_i$ depends on $\alpha_l$ only through $\alpha_l(2\alpha_i-\alpha_l)S$, so a change $\delta$ moves voter $i$'s comparison by $\big(2(\alpha_i-\alpha_l)\delta-\delta^2\big)S=2(\alpha_i-\alpha_l)S\delta+O(\delta^2)$. For a given realization of the noise, the voters whose vote changes are those whose comparison crosses the tie value. Because $\Delta U_i$ is affine in $\theta_i$, the coarea formula, as in the proof of Remark~\ref{rem:hardening}(b), gives their mass as
\[
\delta\int_{\{\Delta U_i=\eta\}}g(\theta_i)\,\frac{2|\alpha_i-\alpha_l|\,S}{\|\nabla\Delta U\|}\,\mathrm{d}\ell+o(\delta),
\qquad
\|\nabla\Delta U\|=2\big((\Delta\beta)^2+(\Delta\alpha)^2S^2\big)^{1/2}.
\]
As $S\to\infty$, $S/\|\nabla\Delta U\|\to1/(2|\Delta\alpha|)$, and the line $\{\Delta U_i=\eta\}$, whose slope is $-\gamma$ and which passes within $|\eta|/(2|\Delta\beta|)$ of $\theta_m$ vertically, converges to the vertical line $\alpha_i=\alpha_m$, on which $|\alpha_i-\alpha_l|=|\Delta\alpha|/2$. The coefficient of $\delta$ therefore converges to $\tfrac12\int_{\{\alpha_i=\alpha_m\}}g\,\mathrm{d}\ell$, which is positive because $\theta_m$ lies in the interior of $\operatorname{supp}(G)$ and $g$ is bounded away from zero there. The bound for a valence shift is Remark~\ref{rem:hardening}(b). Under Assumption~\ref{ass:median}, Corollary~\ref{cor:median} gives $M=\Delta U_{\theta_M}$, so $\partial M/\partial\alpha_l=2(\alpha_M-\alpha_l)S$, and Lemma~\ref{lem:win} gives $\partial\pi/\partial\alpha_l=\phi\,\partial M/\partial\alpha_l$. The effect of a valence shift on $\pi$ is Remark~\ref{rem:hardening}(c).
\end{proof}

\section{Equilibrium: existence, uniqueness and regularity}\label{app:small-h}

Two results support the equilibrium analysis. The first covers the symmetric benchmark at every level of electoral sensitivity; the second covers general distributions and asymmetric parties when electoral sensitivity is moderate. Under Assumption~\ref{ass:small-h}, $M$ is bounded, since $|\Delta U_i|$ is bounded over the compact set $\operatorname{supp}(G)$ and over $S\in[0,\operatorname{Var}(\omega)]$, so the tie value lies inside the support of the noise for $\phi$ small enough.

\begin{lemma}[Global best responses in the symmetric benchmark]\label{lem:global}
Suppose Assumption~\ref{ass:median} holds, $S>0$, and the parties' ideal types are distinct and symmetric about the median point. For every $\phi>0$, the profile of Proposition~\ref{prop:symmetric} is an equilibrium, with the winning probability given by the clipped expression $\min\{1,\max\{0,\tfrac12+\phi M\}\}$, and it is the unique equilibrium.
\end{lemma}

\begin{proof}
Use the inner product $\nu\cdot\nu'=\nu_\beta\nu'_\beta+S\nu_\alpha\nu'_\alpha$, put the median point at zero, so that $\theta_L=-\theta_R$ and $\Lambda\equiv\Lambda_L^R=4\|\theta_R\|^2$. Fix $R$'s nominee at $\theta_r=w\theta_R$, where $w$ is given in Proposition~\ref{prop:symmetric}.

For a fixed norm, $L$ obtains its best policy by placing its nominee on the ray toward $\theta_L$. The winning probability depends only on that norm, since $M=\|\theta_r\|^2-\|\theta_l\|^2$. A candidate farther from zero than $\theta_L$ is dominated by $\theta_L$ whenever the total value of victory is positive: this deviation improves policy and weakly increases the winning probability. A candidate with nonpositive total value cannot improve on the proposed profile. It therefore suffices to consider nominees $\theta_l=w_L\theta_L$ with $0\le w_L\le1$, all of which are feasible by convexity. Then $M=\tfrac14\Lambda(w^2-w_L^2)$ and $\Delta U_L=\tfrac14\Lambda\big[(1+w)^2-(1-w_L)^2\big]$, and $L$'s gain relative to losing is $\Pi_L-U_L^r=\pi\,(B+\Delta U_L)$, where $\pi=\min\{1,\max\{0,\tfrac12+\phi M\}\}$. On the interval with interior $\pi$, its second derivative in $w_L$ is
\[
-\tfrac12\Lambda\pi-\tfrac12\phi\Lambda\,(B+\Delta U_L)-\tfrac12\phi\Lambda^2w_L(1-w_L)<0.
\]
Where $\pi=1$, the gain is increasing in $w_L$; where $\pi=0$, it is zero. At $w_L=w$, its derivative is $\tfrac14\Lambda\big[1-w-2\phi Bw-2\phi\Lambda w^2\big]=0$. Thus $w_L=w$ is a global best response. The argument for the other party is identical.

For uniqueness, let $(\theta_l,\theta_r)$ be any equilibrium and let $\pi$ be $L$'s winning probability. First, neither party can win with certainty. If $\pi=0$ and $B>0$, $L$ can copy $\theta_r$ and receive half the benefit, improving on certain defeat. If $B=0$ and $\theta_r\ne\theta_L$, nominating a type slightly closer to $\theta_L$ than $\theta_r$ improves $L$'s policy and gives it a positive chance of winning. If $B=0$ and $\theta_r=\theta_L$, $R$ can replace $\theta_r$ with $\theta_R$, keeping its certain victory and improving its policy. The case $\pi=1$ is symmetric.

Second, both total values of office are positive. With $B>0$, copying the rival guarantees a positive gain relative to losing. With $B=0$, moving slightly from the rival toward the party ideal does so unless the rival is that ideal. In the exceptional case, both nominees would have to equal that ideal; the other party could then improve by nominating its own ideal without changing the winning probability.

Third, the argument based on norms now implies $\theta_l=w_L\theta_L$ and $\theta_r=w_R\theta_R$, with $w_L,w_R\in[0,1]$. Neither endpoint is optimal. At $w_L=0$, the policy gain from moving outward is first order while the electoral cost is second order; at $w_L=1$, moving inward increases the chance of winning to first order and costs only a second-order policy loss. Equivalently, the derivative of $L$'s gain is $\tfrac12\Lambda\pi>0$ at $w_L=0$ and $-\tfrac12\phi\Lambda\big[B+\tfrac14\Lambda(1+w_R)^2\big]<0$ at $w_L=1$.

Finally, the interior first-order conditions are
\[
\pi(1-w_L)=\phi w_L(B+\Delta U_L),\qquad
(1-\pi)(1-w_R)=\phi w_R(B-\Delta U_R),
\]
with $\Delta U_L=\tfrac14\Lambda\big[(1+w_R)^2-(1-w_L)^2\big]$ and $-\Delta U_R=\tfrac14\Lambda\big[(1+w_L)^2-(1-w_R)^2\big]$. Subtracting and using $\pi=\tfrac12+\tfrac14\phi\Lambda(w_R^2-w_L^2)$ gives
\[
(w_L-w_R)\big[1+2\phi B+\phi\Lambda(w_L+w_R)(2-w_L-w_R)\big]=0.
\]
The bracket is positive, so $w_L=w_R$. The remaining equation, $(1+2\phi B)w+2\phi\Lambda w^2=1$, has exactly one root in $(0,1)$, namely the stated $w$.
\end{proof}

Lemma~\ref{lem:global} establishes that, with symmetric parties' ideal types and a median point, the symmetric profile is the unique equilibrium at every level of electoral sensitivity.

For general distributions, smoothness of the density does not by itself make the tie value $M$ differentiable in the candidates' types everywhere: for the uniform distribution on a triangle, the average type on a median line has kinks when the line passes through a vertex. The following condition is imposed only where it is needed.

\begin{assumption}[Local regularity]\label{ass:local}
The parties' ideal types are distinct and lie in the interior of $\operatorname{supp}(G)$, and for unpredictability in a compact interval $[\underline S,\overline S]$ with $\underline S>0$, the tie value $M$ is twice continuously differentiable jointly in the candidates' types and $S$ on a neighbourhood of $(\theta_L,\theta_R)$ times $[\underline S,\overline S]$.
\end{assumption}

\begin{lemma}[Moderate electoral sensitivity]\label{lem:smallh}
Under Assumption~\ref{ass:local}, there is $\hat \phi>0$ such that for all $\phi\le\hat \phi$ and all $S\in[\underline S,\overline S]$: every best response lies in the neighbourhood of Assumption~\ref{ass:local}; there is a unique equilibrium there; and the equilibrium candidates are continuously differentiable in $S$, with $\theta_j=\theta_J+O(\phi)$ and $\mathrm{d}\theta_j/\mathrm{d}S=O(\phi)$. The candidates are also continuously differentiable in the common benefit $B$, and $\partial\theta_j/\partial B=O(\phi)$.
\end{lemma}

\begin{proof}
Let $\mathcal{N}=\mathcal{N}_L\times\mathcal{N}_R$ be a closed convex neighbourhood of $(\theta_L,\theta_R)$ contained in the neighbourhood of Assumption~\ref{ass:local}. The functions $M$, $B+\Delta U_L$, and $B-\Delta U_R$ are bounded over candidates in $\operatorname{supp}(G)$ and $S\in[\underline S,\overline S]$, and $M(B+\Delta U_L)$ and $M(\Delta U_R-B)$ have bounded first and second derivatives on $\mathcal{N}\times[\underline S,\overline S]$.

\emph{Step 1: best responses lie near the ideal types.} By Lemma~\ref{lem:win}, $\Pi_L=\tfrac12\big(U_L^l+U_L^r+B\big)+\phi M(B+\Delta U_L)$. The first term depends on $\theta_l$ only through $\tfrac12U_L^l$, and by Lemma~\ref{lem:payoff}, $U_L^l(\theta_L)-U_L^l(\theta_l)=(\beta_l-\beta_L)^2+(\alpha_l-\alpha_L)^2S\ge\min\{1,\underline S\}\,\|\theta_l-\theta_L\|^2$. If $\tilde\theta_l$ maximizes $\Pi_L$ against any rival, comparing it with $\theta_L$ shows that $\|\tilde\theta_l-\theta_L\|^2=O(\phi)$, uniformly in the rival's candidate and in $S$. For $\phi$ small enough, every best response of $L$ therefore lies in $\mathcal{N}_L$, and likewise for $R$.

\emph{Step 2: uniqueness of best responses.} On $\mathcal{N}$, the Hessian of $\Pi_L$ with respect to $(\beta_l,\alpha_l)$ is $-\operatorname{diag}(1,S)$ plus $\phi$ times a bounded matrix, so for $\phi$ small enough it is negative definite, uniformly on $\mathcal{N}$. $\Pi_L$ is then strictly concave on $\mathcal{N}_L$, so its maximizer there is unique and interior, solves the first-order conditions of Proposition~\ref{prop:nomination}, and by Step 1 it is $L$'s unique best response against every rival $\theta_r\in \mathcal{N}_R$. The same holds for $R$.

\emph{Step 3: a contraction.} The term $\tfrac12\big(U_L^l+U_L^r\big)$ is additively separable in $\theta_l$ and $\theta_r$, so the cross derivative of $\Pi_L$ with respect to $\theta_l$ and $\theta_r$ is $\phi$ times a bounded matrix. By the implicit function theorem \citep{Sundaram1996} applied to the first-order conditions, each party's best response moves with the rival's nominee only at order $\phi$. For $\phi$ small enough, the map that sends a pair of nominees to the pair of best responses maps the closed set $\mathcal{N}$ into itself by Step 1 and is a contraction, so by the contraction mapping theorem \citep{Ok2007} it has a unique fixed point: there is a unique equilibrium, and by Step 1 every equilibrium lies in $\mathcal{N}$.

\emph{Step 4: regularity and bounds.} Using coordinate order $(\beta,\alpha)$ in this calculation, the first-order condition for $\theta_l$ reads $\operatorname{diag}(1,S)(\theta_l-\theta_L)=\phi\nabla_{\theta_l}\big(M(B+\Delta U_L)\big)$, so $\theta_l-\theta_L=O(\phi)$, and likewise for $R$. The equilibrium solves a system of first-order conditions whose Jacobian is invertible on $\mathcal{N}$ by Steps 2 and 3, so by the implicit function theorem it is continuously differentiable in $S$, $B$, and $\phi$. Differentiating the first-order conditions with respect to $S$, the derivative of the left-hand side at fixed candidates is $(0,\alpha_l-\alpha_L)=O(\phi)$ and the right-hand side is $\phi$ times a bounded term, so $\mathrm{d}\theta_j/\mathrm{d}S=O(\phi)$. Differentiating with respect to $B$ likewise gives $\partial\theta_j/\partial B=O(\phi)$. All bounds are uniform on $[\underline S,\overline S]$ and on the fixed compact interval of benefits.
\end{proof}

At $S=0$, responsiveness is irrelevant to payoffs and the tie-breaking rule in Section~\ref{sec:competition} applies; the selected nomination need not coincide with the limit as $S\to0$, which is why the interval of unpredictability is bounded away from zero.


\section{Proofs for Section~\ref{sec:competition}}







\begin{proof}[Proof of Proposition~\ref{prop:nomination}]
By Lemma~\ref{lem:win}, $\partial \Pi_L/\partial\theta_l=\big(\tfrac12+\phi M\big)\,\partial U_L^l/\partial\theta_l+\phi\,(\partial M/\partial\theta_l)\big(B+\Delta U_L\big)$. Differentiating the condition that defines $M$ in Lemma~\ref{lem:median} gives $\partial M/\partial\theta_l$ as the $g$-weighted average of $\partial\Delta U_i/\partial\theta_l$ over the voters on $\psi_M$. By Lemma~\ref{lem:payoff}, $\partial\Delta U_i/\partial\beta_l=-2(\beta_l-\beta_i)$ and $\partial\Delta U_i/\partial\alpha_l=-2(\alpha_l-\alpha_i)S$, so $\partial M/\partial\beta_l=-2(\beta_l-\overline{\beta})$ and $\partial M/\partial\alpha_l=-2(\alpha_l-\overline{\alpha})S$. The first-order condition for $\beta_l$ is therefore
\[
-2\Big(\tfrac12+\phi M\Big)(\beta_l-\beta_L)-2\phi(\beta_l-\overline{\beta})\big(B+\Delta U_L\big)=0,
\]
and the first-order condition for $\alpha_l$ is the same with $\alpha$ in place of $\beta$, multiplied by $S>0$. Solving each gives the expression for $\theta_l$; the argument for $\theta_r$ is symmetric. The first-order conditions hold because the candidates are interior and $M$ is differentiable, and since the total values of office are positive and the winning probabilities lie in $(0,1)$, $w_L$ and $w_R$ lie in $(0,1)$, so each candidate is a convex combination of her party's ideal type and $\overline{\theta}$. The $g$-weighted average along $\psi_M$ arises because the indifference line is affine, so the factor from the coarea formula \citep{EvansGariepy2015} is constant along it and cancels. Under Assumption~\ref{ass:median}, Corollary~\ref{cor:median} gives $M=\Delta U_{\theta_M}$, so the same steps hold with $\theta_M$ in place of $\overline{\theta}$.
\end{proof}

\section{Proofs for Section~\ref{sec:effects}}

\begin{proof}[Proof of Proposition~\ref{prop:symmetric}]
At symmetric nominations, the candidates' midpoint is $\theta_M$, so $M=0$ and both win with probability $\tfrac12$. Write $\Lambda=\Lambda_L^R$. The policy value of office for either party is $w\Lambda$, since the squared distances from a party ideal to its own and the rival nominee differ by $w\Lambda$. The total value of office is therefore $B+w\Lambda$. Because each candidate wins with probability $\tfrac12$, the winning probability is locally interior for every $\phi$. The first-order condition used in Proposition~\ref{prop:nomination} therefore gives
\[
1-w=2\phi w(B+w\Lambda),
\quad\text{that is,}\quad (1+2\phi B)\,w+2\phi\Lambda w^2=1,
\]
whose root in $(0,1)$ is the stated $w$. Lemma~\ref{lem:global} establishes global optimality and uniqueness at every $\phi>0$.

Differentiating this equation gives $\big(1+2\phi B+4\phi\Lambda w\big)\,\mathrm{d}w=-2\phi w^2\,\mathrm{d}\Lambda-2\phi w\,\mathrm{d}B$, and the equation itself implies $1+2\phi B+4\phi\Lambda w=(1+2\phi\Lambda w^2)/w$. Hence
\[
\frac{\partial w}{\partial\Lambda}=-\frac{2\phi w^3}{1+2\phi\Lambda w^2},\qquad
\frac{\partial w}{\partial B}=-\frac{2\phi w^2}{1+2\phi\Lambda w^2},\qquad
\frac{\partial(w\Lambda)}{\partial\Lambda}=w+\Lambda\frac{\partial w}{\partial\Lambda}=\frac{w}{1+2\phi\Lambda w^2}.
\]
Because $D^*=w^2\Lambda$, $\partial D^*/\partial\Lambda=w^2+2w\Lambda\,\partial w/\partial\Lambda=w^2(1-2\phi\Lambda w^2)/(1+2\phi\Lambda w^2)=(1+2\phi B)w^3/(1+2\phi\Lambda w^2)$, using $1-2\phi\Lambda w^2=(1+2\phi B)w$, and $\partial D^*/\partial B=2w\Lambda\,\partial w/\partial B=-4\phi\Lambda w^3/(1+2\phi\Lambda w^2)$. Multiplying the derivatives in $\Lambda$ by $\partial\Lambda/\partial S=(\alpha_R-\alpha_L)^2$ proves (a)--(c). Also $\partial(B+w\Lambda)/\partial B=1+\Lambda\,\partial w/\partial B=1/(1+2\phi\Lambda w^2)>0$, proving (d). Finally, $2\phi D^*=2\phi\Lambda w^2=1-(1+2\phi B)w<1$, and $w\to0$ as $\Lambda\to\infty$, giving the ceiling.

For (e), the voters' weight $P_V$ multiplies every voter's payoff gap and hence the margin, so $\pi=\tfrac12+\phi P_VM$ with $M$ computed from unweighted payoffs. The parties' weight $P_J$ multiplies their policy payoffs, and dividing a party's payoff by $P_J$ gives the benchmark payoff with electoral sensitivity $\phi P_V$ and benefit $B/P_J$. Substituting into $w$ gives the stated expression, and the nomination condition becomes $(1+2\phi P_VB/P_J)w+2\phi P_V\Lambda w^2=1$. A higher $P_J$ lowers $\phi P_VB/P_J$ and leaves $\phi P_V\Lambda$ unchanged, so $w$ rises when $B>0$; a higher $P_V$ raises both terms, so $w$ falls; $D^*=w^2\Lambda$ moves with $w$. With $P=P_J=P_V$, $(\phi P_V)(B/P_J)=\phi B$, and differentiating $(1+2\phi B)w+2\phi P\Lambda w^2=1$ gives $\partial w/\partial P=-2\phi\Lambda w^2/(1+2\phi B+4\phi P\Lambda w)=-2\phi\Lambda w^3/(1+2\phi P\Lambda w^2)$ and $\partial D^*/\partial P=2w\Lambda\,\partial w/\partial P$.
\end{proof}

\begin{proof}[Proof of Proposition~\ref{prop:asymmetric}]
By Lemma~\ref{lem:smallh}, for $\phi$ small enough there is a unique equilibrium near the parties' ideal types, continuous in $\phi$. By Proposition~\ref{prop:nomination} and Assumption~\ref{ass:median}, $\theta_j-\theta_M=w_J(\theta_J-\theta_M)$. Use the inner product $\nu\cdot\nu'=\nu_\beta\nu'_\beta+\nu_\alpha\nu'_\alpha S$, in which $\Lambda_M^J=\|\theta_J-\theta_M\|^2$. By Corollary~\ref{cor:median} and Proposition~\ref{prop:office},
\begin{equation}\label{eq:Masym}
M=\Lambda_M^r-\Lambda_M^l=w_R^2\Lambda_M^R-w_L^2\Lambda_M^L.
\end{equation}

(a) First, ties. Suppose $M=0$, so that $w_L^2\Lambda_M^L=w_R^2\Lambda_M^R$. The parties' policy values of office are then $\Delta U_L=2\big[w_L\Lambda_M^L-w_R(\theta_L-\theta_M)\cdot(\theta_R-\theta_M)\big]$ and $-\Delta U_R=2\big[w_R\Lambda_M^R-w_L(\theta_L-\theta_M)\cdot(\theta_R-\theta_M)\big]$. Both win with probability $\tfrac12$, so the conditions of Proposition~\ref{prop:nomination} reduce to $1-w_L=2\phi w_L(B+\Delta U_L)$ and $1-w_R=2\phi w_R(B-\Delta U_R)$, and
\[
1-(1+2\phi B)w_L=4\phi\big[w_L^2\Lambda_M^L-w_Lw_R(\theta_L-\theta_M)\cdot(\theta_R-\theta_M)\big]=1-(1+2\phi B)w_R .
\]
Hence $w_L=w_R$, and \eqref{eq:Masym} gives $\Lambda_M^L=\Lambda_M^R$. Conversely, if $\Lambda_M^L=\Lambda_M^R$, the profile with a common $w$ satisfies these conditions with $M=0$, and by uniqueness it is the equilibrium. So the equilibrium margin is zero if and only if $\Lambda_M^L=\Lambda_M^R$. Second, signs. Hold the primitives and $S$ fixed and vary $\phi$. As $\phi\to0$, the candidates converge to the parties' ideal types, so $M\to\Lambda_M^R-\Lambda_M^L$. Since the equilibrium margin is continuous in $\phi$ and can vanish only if $\Lambda_M^L=\Lambda_M^R$, it keeps the sign of $\Lambda_M^R-\Lambda_M^L$ for all $\phi$ small enough, and $l$ wins with probability above $\tfrac12$ if and only if $\Lambda_M^L<\Lambda_M^R$.

(b) The rate follows from the definition of $\Lambda_M^J$. Under the stated conditions, $\Lambda_M^R-\Lambda_M^L$ is positive at $S=0$, since $L$ is closer in baseline, and decreasing in $S$, since $R$ is closer in responsiveness; it is affine and zero at $S^\dagger$. Part (a) then gives the result whenever $S^\dagger$ lies in the attainable range of unpredictability.

(c) The conditions of Proposition~\ref{prop:nomination} give $1-w_L=2\phi(B+\Delta U_L)+O(\phi^2)$ and $1-w_R=2\phi(B-\Delta U_R)+O(\phi^2)$, and by Lemma~\ref{lem:smallh} the remainders have derivatives in $S$ of order $\phi^2$. By Proposition~\ref{prop:general}, $\mathrm{d}\Delta U_L/\mathrm{d}S$ and $-\mathrm{d}\Delta U_R/\mathrm{d}S$ equal $(\alpha_R-\alpha_L)^2+O(\phi)>0$. Both $w_L$ and $w_R$ therefore fall with $S$ for $\phi$ small enough, and each candidate moves toward $\theta_M$ along her segment. For the type-distance comparison in part (c), the norm and inner product are Euclidean. The distance between the candidates' types is $\|w_R(\theta_R-\theta_M)-w_L(\theta_L-\theta_M)\|$, which is increasing in $w_R$ and $w_L$ when $(\theta_L-\theta_M)\cdot(\theta_R-\theta_M)\le0$; so it falls as both fall.

(d) Part (d) is Proposition~\ref{prop:general}.
\end{proof}

\begin{proof}[Proof of Proposition~\ref{prop:welfare-unc}]
(a) By Lemma~\ref{lem:payoff} with $S=R^2\operatorname{Var}(\omega)$, $U_i^j=-\alpha_i^2\operatorname{Var}(\omega)-(\beta_j-\beta_i)^2+R^2\,\alpha_j(2\alpha_i-\alpha_j)\operatorname{Var}(\omega)$, so $\mathrm{d}U_i^j/\mathrm{d}\operatorname{Var}(\omega)=-\alpha_i^2+R^2\alpha_j(2\alpha_i-\alpha_j)=-(1-R^2)\alpha_i^2-R^2(\alpha_j-\alpha_i)^2$.

(b) Player $i$'s equilibrium policy payoff is $U_i^*=\pi U_i^l+(1-\pi)U_i^r$, where $\pi=\tfrac12+\phi M$ by Lemma~\ref{lem:win}. The derivative of $U_i^*$ with respect to $\operatorname{Var}(\omega)$ has three parts. The direct part, at fixed candidates and $\pi$, is $-(1-R^2)\alpha_i^2-R^2\big[\pi(\alpha_l-\alpha_i)^2+(1-\pi)(\alpha_r-\alpha_i)^2\big]\le-R^2 \pi(1-\pi)(\alpha_r-\alpha_l)^2$, since a weighted sum of squared distances to two points is smallest at their weighted mean. With $\alpha_r-\alpha_l=\alpha_R-\alpha_L+O(\phi)$ and $\pi=\tfrac12+O(\phi)$, this bound is negative and bounded away from zero. The part through $\pi$ is $(\mathrm{d}\pi/\mathrm{d}\operatorname{Var}(\omega))(U_i^l-U_i^r)$, which is $O(\phi)$ because $\pi=\tfrac12+\phi M$ with $M$ and its derivatives bounded. The part through the candidates is $O(\phi)$ because $\mathrm{d}\theta_j/\mathrm{d}S=O(\phi)$ (proof of Proposition~\ref{prop:general}). Hence the derivative is negative for $\phi$ small enough. Adding a fixed benefit changes each party's derivative by $B\,\mathrm{d}\pi/\mathrm{d}\operatorname{Var}(\omega)$ or its negative, also $O(\phi)$, so the strict sign persists for total party payoffs.

(c) Part (c) follows from the expression for $U_i^*$ in Proposition~\ref{prop:welfare}: holding $R^2$ fixed,
\[
\frac{\mathrm{d}U_i^*}{\mathrm{d}\operatorname{Var}(\omega)}=-(1-R^2)\alpha_i^2-R^2(\alpha_i-\alpha_M)^2-\tfrac{R^2}{4}\,\frac{\mathrm{d}D^*}{\mathrm{d}S}<0,
\]
because $\mathrm{d}D^*/\mathrm{d}S>0$ by Proposition~\ref{prop:symmetric}. The expected benefit is $B/2$ for either party, so it contributes no derivative.
\end{proof}

\begin{proof}[Proof of Proposition~\ref{prop:welfare}]
In the symmetric equilibrium each candidate wins with probability $\tfrac12$ and the candidates' midpoint is $\theta_M$. By Proposition~\ref{prop:office}, $U_i^*=U_i^i-\tfrac12\big(\Lambda_i^l+\Lambda_i^r\big)$, and since $\Lambda_i^j$ is a squared distance in the inner product $\nu\cdot\nu'=\nu_\beta\nu'_\beta+\nu_\alpha\nu'_\alpha S$, the parallelogram identity gives $\tfrac12\big(\Lambda_i^l+\Lambda_i^r\big)=\Lambda_i^{\theta_M}+\tfrac14\Lambda_l^r$, where $\Lambda_i^{\theta_M}=(\beta_M-\beta_i)^2+(\alpha_M-\alpha_i)^2S$ and $\Lambda_l^r=D^*$. Substituting $U_i^i=-\alpha_i^2\big(\operatorname{Var}(\omega)-S\big)$ gives the expression for $U_i^*$, and differentiating gives $\mathrm{d}U_i^*/\mathrm{d}S$. At $\alpha_M=0$, the derivative is negative for every $i$, since it then equals minus a quarter of $\mathrm{d}D^*/\mathrm{d}S$. The derivative is affine in $\alpha_i$ with slope $2\alpha_M$, which gives the ranking, since $\alpha_R-\alpha_M=\delta$ and $\alpha_L-\alpha_M=-\delta$; solving for a positive derivative gives the threshold.
\end{proof}

\paragraph{Information without symmetry.}
Without symmetry, the candidates need not win with equal probability, and the midpoint response no longer separates from the policy lottery. Holding candidates and the winning probability $\pi$ fixed, the direct effect of information is
\[
\pi\alpha_l(2\alpha_i-\alpha_l)+(1-\pi)\alpha_r(2\alpha_i-\alpha_r).
\]
As electoral sensitivity vanishes, $\pi$ tends to $\tfrac12$ and the candidates tend to the parties' ideal types. The leading-order effect is therefore $\alpha_i(\alpha_L+\alpha_R)-(\alpha_L^2+\alpha_R^2)/2$. The leading-order effect depends on the parties' responses, and the decisive voter's response does not enter. The effects through nominations and winning probabilities add terms of order $\phi$. Voters with responsiveness near zero still lose to leading order. A party can also benefit from a change in its winning probability when its responsiveness is closer to that of the decisive voter. Without a median point, that decisive voter can change with unpredictability, so the identity of the players who benefit may change too.

\section{Different assumptions on information}

In this subsection, I give the details behind Corollaries~\ref{cor:endolearn} and~\ref{cor:timing} and their proofs.

\paragraph{Endogenous learning in office.}
Suppose that the winner chooses the precision of her signal and bears the cost herself. In the normal example, she chooses signal precision $\rho\ge0$, so $\sigma_s^2=1/\rho$, at cost $\zeta \rho^2/2$ with $\zeta>0$. Candidate $j$ maximizes
\[
-\alpha_j^2\big(\sigma_\omega^2-S(\rho)\big)-\tfrac{\zeta}{2}\rho^2,
\qquad
S(\rho)=\frac{\sigma_\omega^4\rho }{\sigma_\omega^2\rho+1}.
\]
When the candidates would choose different precisions, their signals share a common noise, $s_j=\omega+\varepsilon/\sqrt{\rho_j}$ with the same $\varepsilon\sim\operatorname{Norm}(0,1)$ independent of $\omega$, so that equal precisions give a common signal.
Let $S_j^*$ be the unpredictability that her choice of precision produces. Other players anticipate this choice and evaluate her by
\[
U_i^j=-\alpha_i^2\sigma_\omega^2-(\beta_j-\beta_i)^2+\alpha_j(2\alpha_i-\alpha_j)\,S_j^*.
\]
Responsiveness now affects both the candidate's use of information and the amount of information that she acquires. Parties account for both effects when nominating the candidate.

In Corollary~\ref{cor:endolearn}, learning becomes cheaper when $\zeta$ falls, and part (b) means
\[
\frac{\mathrm{d}S_j^*}{\mathrm{d}\sigma_\omega^2}=\frac{\partial S}{\partial \sigma_\omega^2}+\frac{\partial S}{\partial \rho }\,\frac{\mathrm{d}\rho_j^*}{\mathrm{d}\sigma_\omega^2}>\frac{\partial S}{\partial \sigma_\omega^2}.
\]
Part (c) has $\alpha_R=-\alpha_L>0$ and assumes that the limiting nominations, with endogenous learning and with precisions frozen at their limiting values, lie in the interior of $\operatorname{supp}(G)$, and that the electoral margins in both cases are twice continuously differentiable, jointly in the nominations, uncertainty, and the cost scale, near those nominations. The benchmark in part (c) freezes each party's precision at its own equilibrium value and lets the parties reoptimize their nominations. The limiting nominations are given in closed form in the proof.

\begin{proof}[Proof of Corollary~\ref{cor:endolearn}]
(a) Candidate $j$'s first-order condition is $\alpha_j^2\,\partial S/\partial \rho=\zeta \rho$, with $\partial S/\partial \rho=\sigma_\omega^4/(\sigma_\omega^2\rho+1)^2$, which is positive and decreasing in $\rho$ for $\sigma_\omega^2>0$. The objective is therefore strictly concave, and its maximizer is positive if and only if $\alpha_j\ne0$. For $\alpha_j\ne0$, the implicit function theorem gives
\[
\frac{\mathrm{d}\rho_j^*}{\mathrm{d}\alpha_j^2}=\frac{\partial S/\partial \rho}{\zeta-\alpha_j^2\,\partial^2S/\partial \rho^2}>0,
\qquad
\frac{\mathrm{d}\rho_j^*}{\mathrm{d}\sigma_\omega^2}=\frac{\alpha_j^2\,\partial^2S/\partial \rho\,\partial \sigma_\omega^2}{\zeta-\alpha_j^2\,\partial^2S/\partial \rho^2}>0,
\qquad
\frac{\mathrm{d}\rho_j^*}{\mathrm{d}\zeta}=\frac{-\rho_j^*}{\zeta-\alpha_j^2\,\partial^2S/\partial \rho^2}<0,
\]
because $\partial^2S/\partial \rho^2<0$ and $\partial^2S/\partial \rho\,\partial \sigma_\omega^2=2\sigma_\omega^2/(\sigma_\omega^2\rho+1)^3>0$. Weak monotonicity also follows without differentiability from monotone comparative statics \citep{MilgromShannon1994,Topkis1998}: for $\alpha_j\ne0$, the objective has strictly increasing differences in $(\rho,\alpha_j^2)$ and in $(\rho,\sigma_\omega^2)$ and strictly decreasing differences in $(\rho,\zeta)$.

(b) Part (b) follows from (a), since $\partial S/\partial \rho >0$.

(c) We first show that, as $\phi\to0$, each party nominates its own baseline and a candidate whose responsiveness has magnitude
\[
\widehat\alpha_r=\frac{3\alpha_R}{2+\widehat R^2}\in\Big(\alpha_R,\tfrac32\alpha_R\Big),
\qquad\text{where }\widehat R^2\in(0,1)\text{ solves }\frac{\sqrt{\widehat R^2}\,\big(2+\widehat R^2\big)}{\big(1-\widehat R^2\big)^{3/2}}=\frac{3\alpha_R\,\sigma_\omega^{3}}{\sqrt\zeta},
\]
and that the expected squared policy gap is $\widehat D=(\beta_R-\beta_L)^2+36\alpha_R^2\sigma_\omega^2\widehat R^2/(2+\widehat R^2)^2$. Here $\widehat R^2$ is the revealed share at the chosen precision; in the normal example, $R^2/(1-R^2)=\sigma_\omega^2/\sigma_s^2=\sigma_\omega^2\rho$. With fixed finite $B$, as $\phi\to0$ the electoral terms in $\Pi_J$ vanish, and party $J$ chooses its candidate to maximize its payoff when the candidate wins, $U_J^j=-\alpha_J^2\sigma_\omega^2-(\beta_j-\beta_J)^2+\alpha_j(2\alpha_J-\alpha_j)S_j^*$, where $S_j^*$ depends only on $\alpha_j$. The baseline enters only through $(\beta_j-\beta_J)^2$, so $\beta_j=\beta_J$. Consider party $R$, and write $\alpha_r$ for its candidate's responsiveness. The candidate's first-order condition $\zeta \rho=\alpha_r^2\sigma_\omega^4/(1+\sigma_\omega^2\rho)^2$ gives $\alpha_r=\sqrt{\zeta \rho}\,(1+\sigma_\omega^2\rho)/\sigma_\omega^2$, so the party's problem can be written in terms of the induced precision. Since $\alpha_r S=\sigma_\omega^2\sqrt\zeta\,\rho^{3/2}$ and $\alpha_r^2S=\zeta \rho^2(1+\sigma_\omega^2\rho)$, the party maximizes
\[
2\alpha_R \sigma_\omega^2\sqrt\zeta\,\rho^{3/2}-\zeta \rho^2(1+\sigma_\omega^2\rho),
\]
whose first-order condition is $3\alpha_R \sigma_\omega^2/\sqrt\zeta=2\sqrt \rho+3\sigma_\omega^2\rho^{3/2}$. The right-hand side increases from zero to infinity, and the objective increases below the root and decreases above it, so the optimum is unique. Since $\sigma_\omega^2\rho=R^2/(1-R^2)$, the condition becomes the one stated, and substituting back gives $\widehat\alpha_r=3\alpha_R/(2+\widehat R^2)$, which lies in $(\alpha_R,\tfrac32\alpha_R)$. Party $L$ chooses $-\widehat\alpha_r$ by symmetry. The expected squared policy gap is then $(\beta_R-\beta_L)^2+4\widehat\alpha_r^2S(\widehat \rho)=(\beta_R-\beta_L)^2+36\alpha_R^2\sigma_\omega^2\widehat R^2/(2+\widehat R^2)^2$, where
\[
\frac{\mathrm{d}}{\mathrm{d}R^2}\,\frac{R^2}{(2+R^2)^2}=\frac{2-R^2}{(2+R^2)^3}>0 .
\]
The left side of the condition for $\widehat R^2$ increases in $\widehat R^2$, so $\widehat R^2$ rises with $\sigma_\omega^2$ and falls with $\zeta$; hence $\widehat\alpha_r$ falls toward $\alpha_R$ and $\widehat D$ rises as uncertainty rises or learning becomes cheaper. For the comparison, hold precision fixed at $\widehat \rho$ as $\sigma_\omega^2$ changes locally; with exogenous information, the parties nominate their own ideal types in the limit, so the gap is $(\beta_R-\beta_L)^2+4\alpha_R^2S(\widehat \rho,\sigma_\omega^2)$, with derivative $4\alpha_R^2\widehat R^2(2-\widehat R^2)$. Differentiating the condition for $\widehat R^2$ gives $\sigma_\omega^2\,\mathrm{d}\widehat R^2/\mathrm{d}\sigma_\omega^2=3\widehat R^2(2+\widehat R^2)(1-\widehat R^2)/(2+7\widehat R^2)$, and substituting,
\[
\frac{\mathrm{d}\widehat D}{\mathrm{d}\sigma_\omega^2}-4\alpha_R^2\widehat R^2\big(2-\widehat R^2\big)=\frac{4\alpha_R^2\widehat R^2\big(1-\widehat R^2\big)^2\big(7\widehat R^4+30\widehat R^2+56\big)}{\big(2+\widehat R^2\big)^2\big(2+7\widehat R^2\big)}>0 .
\]
For $\alpha_R=\sigma_\omega^2=\zeta=1$, $\widehat\alpha_r\approx1.261$, $\widehat \rho\approx0.612$, and the two derivatives are approximately $3.98$ and $2.46$.


For small positive $\phi$, the argument follows the proof of Lemma~\ref{lem:smallh}, with the limiting nominations above, $\widehat\theta_L=(-\widehat\alpha_r,\beta_L)$ and $\widehat\theta_R=(\widehat\alpha_r,\beta_R)$, in place of the ideal types. As $\phi\to0$, $\Pi_J$ converges uniformly to $\tfrac12\big(U_J^l+U_J^r+B\big)$, which depends on $J$'s candidate only through $\tfrac12U_J^j$. This limiting objective has a unique global maximizer at $\widehat\theta_J$: the baseline enters only through $-(\beta_j-\beta_J)^2$; for party $R$, the objective $\alpha_r(2\alpha_R-\alpha_r)S_j^*$ is negative for $\alpha_r<0$ and for $\alpha_r>2\alpha_R$, positive at $\widehat\alpha_r$, and maximized uniquely there by (c). The maximizer is nondegenerate: in terms of induced precision, the first-order condition is $\zeta \sqrt \rho$ times $3\alpha_R \sigma_\omega^2/\sqrt\zeta-2\sqrt \rho-3\sigma_\omega^2\rho^{3/2}$, a strictly decreasing function of $\rho$, so the second derivative is strictly negative at $\widehat \rho$, and since the map from $\alpha_r$ to induced precision is smooth and strictly increasing, the second derivative in $\alpha_r$ is strictly negative at $\widehat\alpha_r$. By compactness, for any neighbourhood of $\widehat\theta_J$ there is $\delta>0$ such that the limiting objective is at least $\delta$ below its maximum outside it, uniformly over the compact parameter intervals. Because $|M|$, $|B+\Delta U_L|$, and $|B-\Delta U_R|$ are bounded over the candidate set and those parameter intervals, the electoral terms are bounded by a constant times $\phi$, and every best response lies in that neighbourhood once $\phi$ is small enough. In a neighbourhood on which $M$ is twice continuously differentiable, the objective is strictly concave for $\phi$ small, best responses are unique and move with the rival only at order $\phi$, and the best-response map is a contraction, so the equilibrium near $(\widehat\theta_L,\widehat\theta_R)$ is unique. By the implicit function theorem, it is continuously differentiable in $\phi$, $\sigma_\omega^2$ and $\zeta$, and its derivatives converge to those of the limiting nominations as $\phi\to0$, uniformly on the compact parameter intervals. The strict inequalities in (c), $\widehat\alpha_r>\alpha_R$ and the positive derivatives of the policy gap, including the positive difference from the fixed-precision benchmark, whose equilibrium derivative also converges to its limiting value, therefore hold for $\phi$ small enough.
\end{proof}

\paragraph{Learning before the election.}
Some information is already public when parties nominate: a crisis may be under way, or a trend may already be visible in the data. Suppose that this news concerns the same state to which the winner will later respond. The infrastructure that produces the early news may also improve the winner's later information.

Before nomination, all players observe a public signal $s^1$. After the election, the winner observes $s^2$ and chooses $\beta_j+\alpha_j\mathbb{E}[\omega\mid s^1,s^2]$. Let
\[
m^1\equiv\mathbb{E}[\omega\mid s^1],
\qquad
S^2(s^1)\equiv\operatorname{Var}\big(\mathbb{E}[\omega\mid s^1,s^2]\mid s^1\big).
\]
The first quantity is the expected state after the news; the second measures the winner's remaining learning. Conditional on $s^1$, voter $i$'s expected payoff from candidate $j$ is
\[
U_i^j(s^1)=-\alpha_i^2\operatorname{Var}(\omega\mid s^1)-\big[(\beta_j+\alpha_jm^1)-(\beta_i+\alpha_im^1)\big]^2+\alpha_j(2\alpha_i-\alpha_j)\,S^2(s^1).
\]
The conditional payoff is the benchmark payoff with effective baseline $\beta_i+\alpha_im^1$ and unpredictability $S^2(s^1)$. With early news, part of the disagreement over future responses becomes disagreement over known positions. The remaining uncertainty concerns the officeholder's learning after the election.

Suppose that the electorate has a median point and that the parties' ideal types are symmetric about it. The transformation from original to effective types is linear and invertible. The transformation preserves the median-point property and party symmetry. The symmetric benchmark therefore applies after each realization of the news, and I select its symmetric continuation equilibrium. The conditional equilibrium policy gap is $D^*\big(\Lambda(s^1)\big)$, where
\[
D^*(\Lambda)\equiv\frac{4\Lambda}{\big(1+2\phi B+\sqrt{(1+2\phi B)^2+8\phi\Lambda}\big)^2},
\]
and
\[
\Lambda(s^1)\equiv\mathbb{E}\Big[\big(\beta_R-\beta_L+(\alpha_R-\alpha_L)\,\mathbb{E}[\omega\mid s^1,s^2]\big)^2\,\Big|\,s^1\Big].
\]
The quantity $\Lambda(s^1)$ measures the parties' disagreement given the news. For each fixed benefit $B$, the equilibrium gap $D^*(\Lambda)$ is increasing and strictly concave in the parties' disagreement $\Lambda$. It approaches $\Lambda$ as $\phi\to0$ and is bounded by $1/(2\phi)$.

In Corollary~\ref{cor:timing}, part (a) states that the expected squared policy gap equals $(\beta_r-\beta_l)^2+(\alpha_r-\alpha_l)^2\operatorname{Var}\big(\mathbb{E}[\omega\mid s^1,s^2]\big)$. Part (b) compares disclosures before the election $\mathcal H\subseteq\mathcal H'\subseteq\mathcal F$, where $\mathcal F$ is the officeholder's final information, and measures expected polarization by $\mathbb{E}\big[D^*(\Lambda)\big]$; the additional disclosure changes the parties' disagreement when $\Lambda$ changes with positive probability. The next subsection shows numerically how early disclosure interacts with improvements in the officeholder's own information.

\begin{proof}[Proof of Corollary~\ref{cor:timing}]
Conditional on $s^1$, the payoffs are those of Lemma~\ref{lem:payoff} with effective baselines $\beta_i+\alpha_im^1$ and unpredictability $S^2(s^1)$, by the computation in the text, which uses that $\mathbb{E}[\omega\mid s^1,s^2]-m^1$ has conditional variance and conditional covariance with $\omega$ both equal to $S^2(s^1)$. The map $(\alpha,\beta)\mapsto(\alpha,\beta+\alpha m^1)$ is linear and invertible, so it maps every line through $\theta_M$ to a line through its image and preserves the median-point property, as well as the symmetry of the parties' ideal types about the median point; it preserves central symmetry when $G$ is centrally symmetric. When $S^2(s^1)>0$, Proposition~\ref{prop:symmetric} and Lemma~\ref{lem:global} give the selected continuation equilibrium. If $S^2(s^1)=0$ and $\Lambda(s^1)>0$, payoffs depend on candidates only through their effective baselines, and the same argument applies to the one-dimensional effective-baseline game. Select the symmetric nominees on the segments from the party ideals to the median point. If $S^2(s^1)=0$ and $\Lambda(s^1)=0$, both parties' effective ideal policies coincide with the median's. The selected nominees implement this common policy and each wins with probability $\tfrac12$; a unilateral deviation can neither improve the party's policy payoff nor increase its winning probability. This profile is therefore also an equilibrium. In every case, the selected continuation equilibrium has gap $D^*\big(\Lambda(s^1)\big)$.

(a) With candidates fixed at $\theta_l$ and $\theta_r$, the expected squared gap is
\[
\mathbb{E}\Big[\big(\beta_r-\beta_l+(\alpha_r-\alpha_l)\,\mathbb{E}[\omega\mid s^1,s^2]\big)^2\Big],
\]
which by the law of iterated expectations depends only on the distribution of $\mathbb{E}[\omega\mid s^1,s^2]$ and equals the stated expression because its mean is zero.

(b) Let $\Lambda(\mathcal F)=\big(\beta_R-\beta_L+(\alpha_R-\alpha_L)\,\mathbb{E}[\omega\mid\mathcal F]\big)^2$ be the parties' disagreement given the final information $\mathcal F$. The parties' disagreement given the information disclosed before the election is $\Lambda=\mathbb{E}[\Lambda(\mathcal F)\mid\mathcal H]$. If $\mathcal H\subseteq\mathcal H'\subseteq\mathcal F$, then $\Lambda=\mathbb{E}[\Lambda'\mid\mathcal H]$ with $\Lambda'=\mathbb{E}[\Lambda(\mathcal F)\mid\mathcal H']$, and since $D^*$ is concave in $\Lambda$, conditional Jensen's inequality \citep{Billingsley1995} gives $\mathbb{E}[D^*(\Lambda')]\le\mathbb{E}[D^*(\Lambda)]$, strictly if $\Lambda'$ differs from $\Lambda$ with positive probability, because $D^*$ is strictly concave in $\Lambda$. A Blackwell comparison \citep{Blackwell1953} of signals about $\omega$ alone does not deliver this comparison, which concerns nested information about the squared policy gap with the final information held fixed.

(c) Since $0\le D^*(\Lambda)\le\Lambda$ and $\mathbb{E}[\Lambda]<\infty$, dominated convergence gives $\mathbb{E}[D^*(\Lambda)]\to\mathbb{E}[\Lambda]$ as $\phi\to0$; this handles unbounded news, as in the normal example. The limit is the expression in (a) with candidates at the parties' ideal types, which increases with $\operatorname{Var}\big(\mathbb{E}[\omega\mid s^1,s^2]\big)$ when $\alpha_R\ne\alpha_L$. For the rate, $1-w=2\phi w(B+w\Lambda)\le2\phi(B+\Lambda)$ gives $0\le\Lambda-D^*(\Lambda)=\Lambda(1-w^2)\le4\phi\Lambda(B+\Lambda)$. At fixed finite $B$, $D^*(\Lambda)=\Lambda-4\phi\Lambda(B+\Lambda)+o(\phi)$. When $\mathbb{E}[\Lambda^2]$ is finite, $\mathbb{E}[\Lambda]-\mathbb{E}[D^*(\Lambda)]\le4\phi\mathbb{E}[\Lambda(B+\Lambda)]$, and for nested disclosures with $\mathbb{E}[(\Lambda')^2]$ also finite $\mathbb{E}[D^*(\Lambda)]-\mathbb{E}[D^*(\Lambda')]=4\phi\big(\mathbb{E}[(\Lambda')^2]-\mathbb{E}[\Lambda^2]\big)+o(\phi)$. A finite fourth moment of the state is sufficient for this condition. At $B=0$, the selection matters when the news leaves no further learning and the parties' effective ideal policies coincide: a candidate who loses with certainty is then payoff-irrelevant, and other continuation equilibria exist with a positive counterfactual gap.
\end{proof}

\section{Learning before the election: improvements in both signals}
When the same infrastructure informs parties before nomination and the officeholder after the election, better information has two effects. Additional learning in office increases policy divergence. Additional news before nomination gives parties more scope to moderate in response to the circumstances.

Figure~\ref{fig:sticky} compares these effects in the normal example, with $\beta_R-\beta_L=\alpha_R-\alpha_L=\sigma_\omega^2=1$ and $B=0$. If only the signal in office improves, expected polarization rises at every $\phi$. If both signals improve together, polarization rises less under moderate electoral sensitivity. Under high electoral sensitivity, polarization can rise initially and then fall below its no-information level, or remain below that level throughout.

\begin{figure}[t]
\centering
\begin{subfigure}[b]{0.35\linewidth}
\centering
\begin{tikzpicture}
\begin{axis}[
  width=5.6cm, height=5.6cm,
  xmin=0, xmax=1, ymin=0.85, ymax=1.65,
  xtick={0,0.5,1}, ytick={0.9,1,1.2,1.4,1.6},
  xlabel={Share revealed by each signal},
  ylabel={Polarization relative to no information},
  axis lines=left, every axis plot/.append style={thick},
  tick label style={font=\scriptsize}, label style={font=\scriptsize}, title style={font=\small},
  title={(a) $\phi=0.1$}]
  \addplot[gray, thin, domain=0:1] {1};
  \addplot[black] coordinates {(0.00,1.0000) (0.10,1.0728) (0.20,1.1424) (0.30,1.2089) (0.40,1.2726) (0.50,1.3338) (0.60,1.3926) (0.70,1.4491) (0.80,1.5036) (0.90,1.5561) (0.95,1.5817) (0.98,1.5968)};
  \addplot[black, dashed] coordinates {(0.00,1.0000) (0.10,1.0697) (0.20,1.1243) (0.30,1.1660) (0.40,1.1969) (0.50,1.2188) (0.60,1.2331) (0.70,1.2408) (0.80,1.2427) (0.90,1.2394) (0.95,1.2360) (0.98,1.2334)};
\end{axis}
\end{tikzpicture}
\end{subfigure}\hfill
\begin{subfigure}[b]{0.31\linewidth}
\centering
\begin{tikzpicture}
\begin{axis}[
  width=5.6cm, height=5.6cm,
  xmin=0, xmax=1, ymin=0.85, ymax=1.65,
  xtick={0,0.5,1}, ytick={0.9,1,1.2,1.4,1.6},
  xlabel={Share revealed by each signal},
  yticklabels={},
  axis lines=left, every axis plot/.append style={thick},
  tick label style={font=\scriptsize}, label style={font=\scriptsize}, title style={font=\small},
  title={(b) $\phi=0.5$}]
  \addplot[gray, thin, domain=0:1] {1};
  \addplot[black] coordinates {(0.00,1.0000) (0.10,1.0427) (0.20,1.0818) (0.30,1.1177) (0.40,1.1510) (0.50,1.1818) (0.60,1.2106) (0.70,1.2375) (0.80,1.2628) (0.90,1.2866) (0.95,1.2980) (0.98,1.3046)};
  \addplot[black, dashed] coordinates {(0.00,1.0000) (0.10,1.0065) (0.20,1.0172) (0.30,1.0273) (0.40,1.0337) (0.50,1.0357) (0.60,1.0330) (0.70,1.0257) (0.80,1.0133) (0.90,0.9955) (0.95,0.9843) (0.98,0.9767)};
\end{axis}
\end{tikzpicture}
\end{subfigure}\hfill
\begin{subfigure}[b]{0.31\linewidth}
\centering
\begin{tikzpicture}
\begin{axis}[
  width=5.6cm, height=5.6cm,
  xmin=0, xmax=1, ymin=0.85, ymax=1.65,
  xtick={0,0.5,1}, ytick={0.9,1,1.2,1.4,1.6},
  xlabel={Share revealed by each signal},
  yticklabels={},
  axis lines=left, every axis plot/.append style={thick},
  tick label style={font=\scriptsize}, label style={font=\scriptsize}, title style={font=\small},
  title={(c) $\phi=2$}]
  \addplot[gray, thin, domain=0:1] {1};
  \addplot[black] coordinates {(0.00,1.0000) (0.10,1.0229) (0.20,1.0433) (0.30,1.0617) (0.40,1.0784) (0.50,1.0936) (0.60,1.1076) (0.70,1.1205) (0.80,1.1324) (0.90,1.1436) (0.95,1.1488) (0.98,1.1519)};
  \addplot[black, dashed] coordinates {(0.00,1.0000) (0.10,0.9902) (0.20,0.9882) (0.30,0.9894) (0.40,0.9895) (0.50,0.9871) (0.60,0.9814) (0.70,0.9716) (0.80,0.9566) (0.90,0.9341) (0.95,0.9186) (0.98,0.9072)};
\end{axis}
\end{tikzpicture}
\end{subfigure}
\caption{Information timing and expected polarization. Polarization is measured relative to its no-information level in the symmetric normal example, with $\beta_R-\beta_L=\alpha_R-\alpha_L=\sigma_\omega^2=1$ and $B=0$. Each panel uses a different electoral sensitivity $\phi$. The horizontal axis is the variance share revealed by each signal, $\sigma_\omega^2/(\sigma_\omega^2+\sigma_s^2)$. Along the solid curves, only the officeholder's signal improves. Along the dashed curves, the signals before nomination and after the election have equal precision and improve together. Values below one indicate less polarization than without information.}
\label{fig:sticky}
\end{figure}

The comparison suggests a distinction between shocks whose consequences become visible during a term and shocks whose consequences are already visible when candidates are selected.


\section{Proofs for Section~\ref{sec:pre-election}}

\begin{proof}[Proof of Proposition~\ref{prop:strategic}]
Write $S=R^2 \sigma_\omega^2$. The challenger party maximizes $U_{\mathcal C}^{\mathcal I}+(1-\pi)\big(\Delta U_{\mathcal C}+B_{\mathcal C}\big)$, where $\Delta U_{\mathcal C}=\Lambda_{\mathcal C}^{\mathcal I}-\Lambda_{\mathcal C}^c$ does not depend on the incumbent's benefit. Under Assumption~\ref{ass:median}, $1-\pi=\tfrac12-\phi\big(\Lambda_M^c-\Lambda_M^{\mathcal I}\big)$ is concave in $\theta_c$, and $\Delta U_{\mathcal C}+B_{\mathcal C}$ is strictly concave and, for $\phi$ small, positive, so the second term is strictly log-concave where positive and the best response is unique. Its first-order condition gives $\theta_c=w_c\theta_{\mathcal C}+(1-w_c)\theta_M$ with
\[
w_c=\frac{1-\pi}{1-\pi+\phi\,(B_{\mathcal C}+\Delta U_{\mathcal C})}.
\]
As $\phi\to0$, $\pi\to\tfrac12$ uniformly, so $w_c\to1/(1+2\phi B_{\mathcal C})$ and $\alpha_c\to\overline{\alpha}_c$; since $\pi$ and $\phi\Delta U_{\mathcal C}$ depend on $\sigma_\omega^2$ only at order $\phi$, the implicit function theorem gives $\mathrm{d}\theta_c/\mathrm{d}\sigma_\omega^2=O(\phi)$.

The incumbent's payoff is $\Pi_{\mathcal I}=\pi\,(U_{\mathcal I}^{\mathcal I}+B_{\mathcal I})+(1-\pi)U_{\mathcal I}^c$, with $U_{\mathcal I}^{\mathcal I}=-\alpha_{\mathcal I}^2(1-R^2)\sigma_\omega^2$ and $U_{\mathcal I}^c=U_{\mathcal I}^{\mathcal I}-\Lambda_{\mathcal I}^c$. Differentiating,
\[
\frac{\mathrm{d}\Pi_{\mathcal I}}{\mathrm{d}\sigma_\omega^2}=\frac{\mathrm{d}\pi}{\mathrm{d}\sigma_\omega^2}\big(\Lambda_{\mathcal I}^c+B_{\mathcal I}\big)-(1-R^2)\alpha_{\mathcal I}^2-(1-\pi)\,R^2(\alpha_c-\alpha_{\mathcal I})^2-(1-\pi)\,\nabla_{\theta_c}\Lambda_{\mathcal I}^c\cdot\frac{\mathrm{d}\theta_c}{\mathrm{d}\sigma_\omega^2},
\]
with $\mathrm{d}\pi/\mathrm{d}\sigma_\omega^2=\phi\big[R^2\big((\alpha_c-\alpha_M)^2-(\alpha_{\mathcal I}-\alpha_M)^2\big)+\nabla_{\theta_c}M\cdot\mathrm{d}\theta_c/\mathrm{d}\sigma_\omega^2\big]$. Substituting the limits above,
\[
\frac{\mathrm{d}\Pi_{\mathcal I}}{\mathrm{d}\sigma_\omega^2}=R^2\,\phi B_{\mathcal I}\big[(\overline{\alpha}_c-\alpha_M)^2-(\alpha_{\mathcal I}-\alpha_M)^2\big]-(1-R^2)\alpha_{\mathcal I}^2-\tfrac12R^2(\overline{\alpha}_c-\alpha_{\mathcal I})^2+O(\phi),
\]
uniformly on the compact parameter sets. The leading term is positive if and only if $\phi B_{\mathcal I}>\overline{B}$, so there is $\overline{\phi}$ below which the derivative is positive; $\mathrm{d}\pi/\mathrm{d}\sigma_\omega^2=\phi R^2\big[(\overline{\alpha}_c-\alpha_M)^2-(\alpha_{\mathcal I}-\alpha_M)^2\big]+O(\phi^2)$ is positive as well, which proves (a). For (b), the expected squared gap is $(\beta_c-\beta_{\mathcal I})^2+(\alpha_c-\alpha_{\mathcal I})^2R^2 \sigma_\omega^2$, whose derivative is $R^2(\overline{\alpha}_c-\alpha_{\mathcal I})^2+O(\phi)>0$. For (c), player $i$'s policy payoff is $U_i^*=\pi\,U_i^{\mathcal I}+(1-\pi)U_i^c$; since $\mathrm{d}\pi/\mathrm{d}\sigma_\omega^2=O(\phi)$, $U_i^{\mathcal I}-U_i^c$ is bounded, and $\mathrm{d}\theta_c/\mathrm{d}\sigma_\omega^2=O(\phi)$,
\[
\frac{\mathrm{d}U_i^*}{\mathrm{d}\sigma_\omega^2}=-(1-R^2)\alpha_i^2-\tfrac12R^2\big[(\alpha_{\mathcal I}-\alpha_i)^2+(\overline{\alpha}_c-\alpha_i)^2\big]+O(\phi)\le-\tfrac14R^2(\overline{\alpha}_c-\alpha_{\mathcal I})^2+O(\phi),
\]
uniformly over players, because $(\alpha_{\mathcal I}-\alpha_i)^2+(\overline{\alpha}_c-\alpha_i)^2\ge\tfrac12(\overline{\alpha}_c-\alpha_{\mathcal I})^2$. For a higher revealed share at fixed $\sigma_\omega^2$, the same computation without the residual-risk term gives $\mathrm{d}\Pi_{\mathcal I}/\mathrm{d}S=\phi B_{\mathcal I}\big[(\overline{\alpha}_c-\alpha_M)^2-(\alpha_{\mathcal I}-\alpha_M)^2\big]+\alpha_{\mathcal I}^2-\tfrac12(\overline{\alpha}_c-\alpha_{\mathcal I})^2+O(\phi)$, which is positive if and only if $\phi B_{\mathcal I}>\underline B$, up to the $O(\phi)$ term.
\end{proof}

\section{General distributions}

Suppose Assumption~\ref{ass:local} holds. Let $D^*$ and $\Delta U_J^*$ denote the expected squared policy gap and party $J$'s payoff gap at the equilibrium of Lemma~\ref{lem:smallh}, which is unique near the parties' ideal types and differentiable in $S$ for $\phi$ small enough.

\begin{proposition}[General distributions]\label{prop:general}
Under Assumption~\ref{ass:local}, for all $S\in[\underline S,\overline S]$,
\begin{equation}\label{eq:general}
\frac{\mathrm{d}D^*}{\mathrm{d}S}=(\alpha_R-\alpha_L)^2+O(\phi),
\qquad
\frac{\mathrm{d}\Delta U_L^*}{\mathrm{d}S}=(\alpha_R-\alpha_L)^2+O(\phi),
\qquad
-\frac{\mathrm{d}\Delta U_R^*}{\mathrm{d}S}=(\alpha_R-\alpha_L)^2+O(\phi),
\end{equation}
where the two remainders for the values of office may differ. Hence, if $\alpha_R\ne\alpha_L$, for $\phi$ small enough the policy gap in equilibrium and both parties' values of office rise with uncertainty and with the revealed share.
\end{proposition}

\begin{proof}
By Lemma~\ref{lem:smallh}, $\theta_j=\theta_J+O(\phi)$ and $\mathrm{d}\theta_j/\mathrm{d}S=O(\phi)$, uniformly on $[\underline S,\overline S]$. Differentiating $D=(\beta_r-\beta_l)^2+(\alpha_r-\alpha_l)^2S$ at the equilibrium candidates gives $(\alpha_r-\alpha_l)^2$ plus terms proportional to $\mathrm{d}\theta_j/\mathrm{d}S$, which are $O(\phi)$, and $(\alpha_r-\alpha_l)^2=(\alpha_R-\alpha_L)^2+O(\phi)$. The same argument applies to $\Delta U_L=\Lambda_L^r-\Lambda_L^l$, whose derivative at $\theta_l=\theta_L$, $\theta_r=\theta_R$ is $(\alpha_R-\alpha_L)^2$, and to $-\Delta U_R=\Lambda_R^l-\Lambda_R^r$, with separate remainders.
\end{proof}

The gap between the candidates' types is
\begin{equation}\label{eq:typegap}
\theta_r-\theta_l=w_R\,\theta_R-w_L\,\theta_L+(w_L-w_R)\,\overline{\theta} .
\end{equation}
The conditions of Proposition~\ref{prop:nomination} give $1-w_L=2\phi(B+\Delta U_L)+O(\phi^2)$ and $1-w_R=2\phi(B-\Delta U_R)+O(\phi^2)$, and the two policy values of office differ by $O(\phi)$, so $w_L-w_R=O(\phi^2)$. When electoral sensitivity is moderate, the movement of $\overline{\theta}$ therefore enters the gap only at order $\phi^2$, and falling $w_L$ and $w_R$ pull the candidates together, as in Proposition~\ref{prop:symmetric}. When electoral sensitivity is high and $w_L\ne w_R$, the movement of $\overline{\theta}$ can widen or narrow the gap. The electoral advantage also moves with the decisive point: as uncertainty rises, the winning probability shifts toward the party whose responsiveness is closer to that of the voters around $\overline{\theta}$. Under high electoral sensitivity, outside the range of Proposition~\ref{prop:general}, the policy gap can fall with unpredictability, even with a median point. Let $G$ be uniform on $[-1.1,1.1]^2$, so that $\theta_M=(0,0)$ and both ideal types below are interior, with $\theta_L=(-1,-1)$, $\theta_R=(-\tfrac12,-\tfrac13)$, $\phi=\tfrac14$, $B=0$, and $S=\operatorname{Var}(\omega)=\tfrac1{10}$. In equilibrium, $l$ wins with probability about $0.37$, the candidates are approximately $(-0.77,-0.77)$ and $(-0.46,-0.31)$, both are global best responses, and $\mathrm{d}D^*/\mathrm{d}S\approx-0.0061$. The decisive point does not move; the trailing party's candidate converges toward it by more than information pushes the policies apart. Assumption~\ref{ass:small-h} holds in this example, since $\phi\overline{M}\approx0.33<\tfrac12$, so interior winning probabilities alone do not deliver Propositions~\ref{prop:asymmetric}(d) and~\ref{prop:welfare-unc}(b); both require $\phi$ to be sufficiently small.

\begin{corollary}[Office benefits under moderate electoral sensitivity]\label{cor:benefits}
Under Assumption~\ref{ass:local}, with the common office benefit in a fixed compact interval, the equilibrium of Lemma~\ref{lem:smallh} satisfies
\begin{align*}
\frac{\partial}{\partial B}\|\theta_r-\theta_l\|^2
&=-4\phi\|\theta_R-\theta_L\|^2+O(\phi^2)<0,\\
\frac{\partial D^*}{\partial B}
&=-4\phi\Lambda_L^R+O(\phi^2)<0
\end{align*}
for all sufficiently small $\phi>0$. Thus higher common office benefits reduce both type and policy gaps even without a median point.
\end{corollary}

\begin{proof}
Let $H=\operatorname{diag}(S,1)$, following the coordinate order $(\alpha,\beta)$, and let $\overline{\theta}_0$ be the average type on the median line for the ideal candidates. The first-order conditions and Lemma~\ref{lem:smallh} give
\[
\theta_j=\theta_J+2\phi(B+\Lambda_L^R)(\overline{\theta}_0-\theta_J)+O(\phi^2).
\]
To differentiate in $B$, use the first-order conditions directly. Their Jacobian is block diagonal with blocks $-H$, up to $O(\phi)$, while their direct $B$ derivatives are $2\phi H(\overline{\theta}_0-\theta_J)+O(\phi^2)$. Hence
\[
\frac{\partial\theta_j}{\partial B}=2\phi(\overline{\theta}_0-\theta_J)+O(\phi^2),
\qquad
\frac{\partial(\theta_r-\theta_l)}{\partial B}=-2\phi(\theta_R-\theta_L)+O(\phi^2).
\]
Differentiating the squared Euclidean gap and the squared gap in the metric $H$ proves the result. The leading coefficients are strictly negative because the party ideals are distinct and $S>0$.
\end{proof}

\begin{corollary}\label{cor:leadgeneral}
Under Assumption~\ref{ass:local}, for $\phi$ small enough, in the equilibrium of Lemma~\ref{lem:smallh}, candidate $l$ wins with probability above $\tfrac12$ if and only if $M(\theta_L,\theta_R)>0$, where $M(\theta_L,\theta_R)$ is the tie value of the noise when each party nominates its own ideal type.
\end{corollary}

\begin{proof}
The argument follows the proof of Proposition~\ref{prop:asymmetric}(a), with the median point replaced by the relevant average type on the median line. Let $\overline{\theta}_0$ be the average type of the voters on the median line for candidates $\theta_L$ and $\theta_R$. Candidates of the form $\theta_j=w\,\theta_J+(1-w)\overline{\theta}_0$ with a common $w$ leave the direction of the candidates' difference, and hence the slope of the line of indifferent voters, the median line with that slope and $\overline{\theta}_0$, unchanged. If $M(\theta_L,\theta_R)=0$, the parties' midpoint lies on that median line, and so does the candidates' midpoint, so $M=0$; with $\overline{\theta}_0$ in place of $\theta_M$, the computation in the proof of Proposition~\ref{prop:asymmetric}(a) shows that the common benefit still forces both parties to choose the same $w$, so this profile is the equilibrium. Conversely, if the equilibrium margin is zero, both parties win with probability $\tfrac12$ and the same computation forces $w_L=w_R$, so the ideal pair and the nominated pair share the median line and $M(\theta_L,\theta_R)=0$. Holding the primitives fixed and varying $\phi$, the equilibrium margin is continuous, tends to $M(\theta_L,\theta_R)$ as $\phi\to0$, and can vanish only if $M(\theta_L,\theta_R)=0$, so it keeps the sign of $M(\theta_L,\theta_R)$.
\end{proof}

\section{Several states and several policies}\label{app:multi}

Let there be policies $k=1,\dots,K$, chosen by the officeholder, and a state $\omega\in\mathbb{R}^N$ with mean zero and covariance matrix $\Sigma\equiv\operatorname{Var}(\omega)$. Player $i$'s target is the vector $\tau(\omega,\theta_i)=\beta_i+A_i\,\omega$. The vector $\beta_i\in\mathbb{R}^K$ collects her baselines, and the $K\times N$ matrix $A_i$ collects her responsiveness: its entry $\alpha_{ikn}$ is her responsiveness of policy $k$ to state $n$. Her policy payoff is $-\|x-\tau(\omega,\theta_i)\|^2$. A weighted quadratic loss reduces to this case by a linear change of the policy coordinates, which transforms the types linearly. The states may be correlated, and a state may bear on several policies. Unpredictability is the matrix $S\equiv\operatorname{Var}(\mathbb{E}_\mu[\omega])$, which satisfies $0\preceq S\preceq\Sigma$ in the positive semidefinite order. Information rises when $S$ rises in this order, and uncertainty rises at a fixed revealed share when $\Sigma$ and $S$ rise in proportion.

\begin{proposition}[Several states and several policies]\label{prop:multistate}
In this setting:
\begin{enumerate}
  \item[(a)] For any types and distribution of posteriors,
  \[
  \begin{aligned}
  U_i^j={}&-\operatorname{tr}\!\big(A_i(\Sigma-S)A_i'\big)-\|\beta_j-\beta_i\|^2\\
  &-\operatorname{tr}\!\big((A_j-A_i)\,S\,(A_j-A_i)'\big).
  \end{aligned}
  \]
  \item[(b)] With candidates fixed, voter $i$'s payoff gap $U_i^l-U_i^r$ is affine in her type, with gradient $2(\beta_l-\beta_r)$ in her baselines and $2(A_l-A_r)S$ in her responsiveness. When $S$ rises in proportion, the gradient in responsiveness scales up and the gradient in baselines is unchanged, so the hyperplane of indifferent voters rotates, as in Proposition~\ref{prop:realign}. The cost of delegation
  \[
  \Lambda_i^j\equiv\|\beta_j-\beta_i\|^2+\operatorname{tr}\big((A_j-A_i)S(A_j-A_i)'\big)
  \]
  is nondecreasing in $S$, and it rises strictly when $S$ rises to $S'$ with $(A_j-A_i)(S'-S)(A_j-A_i)'\ne0$. The expected squared gap between the candidates' policies is $\Lambda_l^r$, and the expected signed gap stays $\beta_r-\beta_l$. Propositions~\ref{prop:office} and~\ref{prop:auto-polar} and Corollary~\ref{cor:office} hold in this sense.
  \item[(c)] If $S\succ0$ and every hyperplane through $\theta_M$ divides $G$ in half, Lemmas~\ref{lem:median} and~\ref{lem:win}, Propositions~\ref{prop:nomination} and~\ref{prop:symmetric}, and Lemma~\ref{lem:global} hold with the inner product $\langle\nu,\tilde\nu\rangle_S\equiv\nu_\beta^\top\tilde\nu_\beta+\operatorname{tr}\big(\nu_A S\tilde\nu_A^\top\big)$ and $\Lambda$ as in (b). Each comparative static in $S$ in these results holds for a rise in $S$ in the positive semidefinite order that raises $\Lambda_L^R$ strictly, which replaces the condition $\alpha_R\ne\alpha_L$. Parts (a) and (b) hold for any positive semidefinite $S$; part (c) needs $S\succ0$, because with a singular $S$, responsiveness in a null direction of $S$ changes neither policies nor comparisons, and payoff-equivalent nominations multiply.
\end{enumerate}
\end{proposition}

\begin{proof}
Write $m=\mathbb{E}_\mu[\omega]$. Officeholder $j$ sets $x=\beta_j+A_jm$. Because $\omega-m$ has mean zero given the posterior, $\mathbb{E}\big[m(\omega-m)'\big]=0$ and $\operatorname{Var}(\omega-m)=\Sigma-S$, and $\mathbb{E}[m]=0$. Hence
\[
\begin{aligned}
\mathbb{E}\|x-\tau(\omega,\theta_i)\|^2
&=\|\beta_j-\beta_i\|^2+\mathbb{E}\big\|(A_j-A_i)m-A_i(\omega-m)\big\|^2\\
&=\|\beta_j-\beta_i\|^2+\operatorname{tr}\big((A_j-A_i)S(A_j-A_i)'\big)+\operatorname{tr}\big(A_i(\Sigma-S)A_i'\big),
\end{aligned}
\]
which is (a). Expanding the squares in (a), the terms that involve both the voter's type and the candidates' types are linear in the voter's type, with the gradients in (b). For any matrix $D$ and any positive semidefinite $P$, $\operatorname{tr}(DPD')\ge0$, with equality if and only if $DPD'=0$; applied to $P=S'-S$, this gives the monotonicity in (b). The signed gap is $\mathbb{E}[x_r-x_l]=\beta_r-\beta_l$ because $\mathbb{E}[m]=0$. For (c), the proofs of these results use payoffs only through the quadratic form in (a) and the inner product that it defines, and none uses the dimension of the type space. In Proposition~\ref{prop:symmetric}, the weight $w$ and the policy gap depend on $S$ only through $\Lambda_L^R$.
\end{proof}

Suppose now that $N=K$, that each policy $n$ has its own state $\omega_n$, so that $A_i$ is diagonal with entries $\alpha_{in}$, and that the states are independent, so that $\Sigma$ is diagonal with entries $\operatorname{Var}(\omega_n)$. Independence does not make $S$ diagonal, because a signal about a combination of states can correlate their posterior means. With diagonal responsiveness matrices, however, payoffs depend on $S$ only through its diagonal entries $S_n\equiv S_{nn}=R_n^2\operatorname{Var}(\omega_n)$. A type is then a baseline and a responsiveness on each issue, $\theta_i=(\beta_{in},\alpha_{in})_{n=1}^N\in\mathbb{R}^{2N}$, and Proposition~\ref{prop:multistate}(a) gives
\[
\Lambda_i^j=\sum_n\Big[(\beta_{jn}-\beta_{in})^2+(\alpha_{jn}-\alpha_{in})^2S_n\Big].
\]

\begin{proposition}[Independent issues]\label{prop:issues}
In this setting, suppose $S_n>0$ for every $n$, and the parties' ideal types are distinct and lie in the interior of $\operatorname{supp}(G)$.
\begin{enumerate}
  \item[(a)] Under Assumptions~\ref{ass:median} and~\ref{ass:small-h}, for all sufficiently small $\phi>0$, candidate $l$ wins with probability above $\tfrac12$ if and only if $\Lambda_M^L<\Lambda_M^R$, where
  \[
  \Lambda_M^J=\sum_n\Big[(\beta_{Jn}-\beta_{Mn})^2+(\alpha_{Jn}-\alpha_{Mn})^2S_n\Big].
  \]
   As $S_n$ rises, $\Lambda_M^R-\Lambda_M^L$ changes at rate $(\alpha_{Rn}-\alpha_{Mn})^2-(\alpha_{Ln}-\alpha_{Mn})^2$, so the lead moves toward the party whose responsiveness on issue $n$ is closer to $\alpha_{Mn}$, and the lead switches at the value of $S_n$ at which $\Lambda_M^L=\Lambda_M^R$, if that value is positive and attainable. Parts (c) and (d) of Proposition~\ref{prop:asymmetric} hold for a rise in each $S_n$ with $\alpha_{Rn}\ne\alpha_{Ln}$.
  \item[(b)] Under Assumption~\ref{ass:local} on $\mathbb{R}^{2N}$ in place of Assumption~\ref{ass:median}, for all sufficiently small $\phi>0$, candidate $l$ wins with probability above $\tfrac12$ if and only if $M(\theta_L,\theta_R)>0$ (Corollary~\ref{cor:leadgeneral}). Let $\overline{\alpha}_{0n}$ be the average responsiveness on issue $n$ of the voters on the median hyperplane at the parties' ideal types, weighted by the density of $G$ on that hyperplane. Then
  \[
  \frac{\mathrm{d}\pi^*}{\mathrm{d}S_n}=\phi\Big[(\alpha_{Rn}-\overline{\alpha}_{0n})^2-(\alpha_{Ln}-\overline{\alpha}_{0n})^2\Big]+O(\phi^2),
  \]
  so if the bracket is nonzero, then for all sufficiently small $\phi$ a rise in $S_n$ moves the lead toward the party whose responsiveness on issue $n$ is closer to $\overline{\alpha}_{0n}$. Because this average changes with the slope of the hyperplane, the statement is local.
  \item[(c)] Under Assumption~\ref{ass:median}, suppose that the incumbent of Section~\ref{sec:pre-election} can raise uncertainty about issue $n$, with $R_n^2$ and the other issues fixed, and let $\overline{\alpha}_{cn}\equiv\alpha_{Mn}+(\alpha_{\mathcal{C}n}-\alpha_{Mn})/(1+2\phi B_{\mathcal C})$. Proposition~\ref{prop:strategic} holds for this rise with $\alpha_{\mathcal{I}n}$, $\overline{\alpha}_{cn}$, $\alpha_{Mn}$, and $R_n^2$ in place of $\alpha_{\mathcal I}$, $\overline{\alpha}_c$, $\alpha_M$, and $R^2$, whatever the incumbent's positions on the other issues.
  \item[(d)] Propositions~\ref{prop:welfare-unc} and~\ref{prop:welfare} hold issue by issue under their original assumptions: party symmetry and a median point for Proposition~\ref{prop:welfare} and for Proposition~\ref{prop:welfare-unc}(c), and local regularity with sufficiently small $\phi$ for Proposition~\ref{prop:welfare-unc}(b). For information about issue $n$, $\alpha_{Mn}(2\alpha_{in}-\alpha_{Mn})$ replaces $\alpha_M(2\alpha_i-\alpha_M)$, and the strict conclusions about issue $n$ require $\alpha_{Rn}\ne\alpha_{Ln}$.
\end{enumerate}
\end{proposition}

\begin{proof}
Each part applies Proposition~\ref{prop:multistate} to diagonal matrices. Every payoff and every margin is then a sum over issues, and the terms for the other issues do not depend on $S_n$. For (a), the proof of Proposition~\ref{prop:asymmetric} uses payoffs only through $\Lambda_M^J$ and the inner product of Proposition~\ref{prop:multistate}(c), and differentiating $\Lambda_M^R-\Lambda_M^L$ in $S_n$ gives the rate. For (b), the gap $U_i^l-U_i^r$ is affine in $\theta_i$, so the voters who are indifferent at the tie value form a hyperplane $H$. As $S_n$ rises, the gap changes per unit by
\[
(\alpha_{Rn}-\alpha_{in})^2-(\alpha_{Ln}-\alpha_{in})^2=\alpha_{Rn}^2-\alpha_{Ln}^2-2\alpha_{in}(\alpha_{Rn}-\alpha_{Ln}),
\]
which is affine in $\alpha_{in}$. Its average over $H$, weighted by the density of $G$, is therefore $(\alpha_{Rn}-\overline{\alpha}_{0n})^2-(\alpha_{Ln}-\overline{\alpha}_{0n})^2$, which is positive if and only if $L$ is closer to $\overline{\alpha}_{0n}$. By the implicit function theorem applied to the condition that defines the tie value, $\mathrm{d}M(\theta_L,\theta_R)/\mathrm{d}S_n$ equals this average. By Lemma~\ref{lem:smallh}, the equilibrium nominees are within $O(\phi)$ of the parties' ideal types, so the equilibrium margin is $M^*=M(\theta_L,\theta_R)+O(\phi)$, and by Lemma~\ref{lem:win}, $\pi^*=\tfrac12+\phi M^*$. Hence $\mathrm{d}\pi^*/\mathrm{d}S_n=\phi\,\mathrm{d}M(\theta_L,\theta_R)/\mathrm{d}S_n+O(\phi^2)$, which gives the displayed expression; when the bracket is nonzero, its sign determines the direction for all sufficiently small $\phi$. For (c), the proof of Proposition~\ref{prop:strategic} differentiates the incumbent's payoff and the margin in uncertainty, and the challenger's nominee does not depend on uncertainty in the limit $\phi\to0$. With independent issues, the derivatives in $\operatorname{Var}(\omega_n)$ involve only the terms for issue $n$, which gives the stated substitution. Part (d) follows in the same way from Proposition~\ref{prop:multistate}(c) restricted to issue $n$.
\end{proof}

\section{Log-concave electoral noise}

\begin{proposition}[Log-concave noise]\label{prop:logconcave}
Let $\eta=\iota/\phi$, where $\iota$ has a log-concave, continuously differentiable density $\xi$ with full support on $\mathbb R$, mean and median zero, and $\xi(0)=1$. Let $\Xi$ be its distribution function. Then $\Pr(l\text{ wins})=\Xi(\phi M)$, and:
\begin{enumerate}
  \item[(a)] both winning probabilities lie in $(0,1)$ for every pair of candidates;
  \item[(b)] under Assumption~\ref{ass:median} and $S>0$, each party's gain relative to losing is strictly log-concave in its own nominee on the set where its total value of office is positive. This set contains its best response, except when $B=0$ and the rival is the party's ideal type. Best responses are unique and continuous, and an equilibrium exists for every $\phi>0$;
  \item[(c)] under Assumption~\ref{ass:median}, $S>0$, and distinct party ideals symmetric about the median point, the profile of Proposition~\ref{prop:symmetric}, including the common benefit $B$, is the unique equilibrium for every $\phi>0$;
  \item[(d)] in any equilibrium with interior candidates and positive total values of office,
  \[
  w_L=\frac{\Xi(\phi M)}{\Xi(\phi M)+\phi\xi(\phi M)(B+\Delta U_L)},
  \qquad
  w_R=\frac{1-\Xi(\phi M)}{1-\Xi(\phi M)+\phi\xi(\phi M)(B-\Delta U_R)}.
  \]
  If $\xi$ is symmetric, $w_L$ and $w_R$ are weakly larger than under uniform noise at the same $\phi$, margin, and total value of office, wherever the uniform winning probability is interior;
  \item[(e)] the results for sufficiently low electoral sensitivity continue to hold, with $\Xi(\phi M)=\tfrac12+\phi M+O(\phi^2)$ in place of $\tfrac12+\phi M$ and fixed finite benefits.
\end{enumerate}
\end{proposition}

\begin{proof}
(a) Full support makes both winning probabilities positive.

(b) Party $L$'s payoff is $\Pi_L=U_L^r+\Xi(\phi M)(B+\Delta U_L)$. If $B>0$ or $\theta_r\ne\theta_L$, nominating $\theta_L$ gives a positive total value of office. A best response must therefore lie in the convex set where $B+\Delta U_L>0$. On this set,
\[
\log(\Pi_L-U_L^r)=\log \Xi(\phi M)+\log(B+\Delta U_L).
\]
Under the median-point assumption, $M=\|\theta_r-\theta_M\|^2-\|\theta_l-\theta_M\|^2$ in the policy metric, and $M$ is concave in $\theta_l$. The function $\log \Xi$ is increasing and concave because $\xi$ is log-concave \citep{BagnoliBergstrom2005}, so the first term is concave. The total value $B+\Delta U_L=B+\Lambda_L^r-\|\theta_l-\theta_L\|^2$ is positive and strictly concave on the relevant set, so its logarithm is strictly concave. The maximizer is unique. If $B=0$ and $\theta_r=\theta_L$, the unique best response is $\theta_L$. The argument for $R$ uses $1-\Xi(\phi M)$ and is otherwise identical. The maximum theorem \citep{Berge1963} gives continuous best responses on the compact convex strategy set, and Brouwer's theorem \citep{Border1985} gives an equilibrium.

(c) At symmetric nominations, $M=0$, $\Xi(0)=\tfrac12$, and $\xi(0)=1$. The first-order conditions therefore coincide with those in Proposition~\ref{prop:symmetric}; by (b), the resulting nominees are global best responses. To establish uniqueness, use the norm argument in Lemma~\ref{lem:global}. Normalize $\theta_M=0$, so that $\theta_L=-a$ and $\theta_R=a$ for some $a\in\mathbb{R}^2$. Positive winning probabilities and positive total values of office imply that every equilibrium has nominees $\theta_l=-w_La$ and $\theta_r=w_Ra$ with $0<w_L,w_R<1$. Their first-order conditions are
\[
\Xi(\phi M)(1-w_L)=\phi\xi(\phi M)w_L(B+\Delta U_L),\qquad
\big(1-\Xi(\phi M)\big)(1-w_R)=\phi\xi(\phi M)w_R(B-\Delta U_R),
\]
with $\Delta U_L$ and $\Delta U_R$ as in the proof of Lemma~\ref{lem:global}. Subtracting gives
\[
\big[1+2\phi\xi(\phi M)B\big](w_L-w_R)
=(2-w_L-w_R)\big[2\Xi(\phi M)-1+2\phi\xi(\phi M)M\big].
\]
The bracket on the right has the sign of $M=\tfrac14\Lambda(w_R^2-w_L^2)$, opposite to the sign of $w_L-w_R$ unless $w_L=w_R$. Both other factors are positive, so $w_L=w_R$. The unique root is the one in Proposition~\ref{prop:symmetric}. This argument does not require symmetry of $\xi$.

(d) Differentiating $\Pi_L$ gives
\[
\Xi(\phi M)\nabla U_L^l+\phi\xi(\phi M)(B+\Delta U_L)\nabla M=0,
\]
which yields $w_L$; $w_R$ follows in the same way. Then
\[
w_L=\left[1+\frac{\phi\xi(\phi M)(B+\Delta U_L)}{\Xi(\phi M)}\right]^{-1},
\qquad
w_R=\left[1+\frac{\phi\xi(\phi M)(B-\Delta U_R)}{1-\Xi(\phi M)}\right]^{-1}.
\]
For a symmetric log-concave density, $\xi'\le0$ on $(0,\infty)$ and $\xi'\ge0$ on $(-\infty,0)$. Thus $(\Xi/\xi)'=1-\Xi\xi'/\xi^2$ is at least one to the right of zero and at most one to its left. Since $\Xi(0)/\xi(0)=\tfrac12$, $\Xi(y)/\xi(y)\ge\tfrac12+y$. Applying the same argument to $\big((1-\Xi)/\xi\big)'=-1-(1-\Xi)\xi'/\xi^2$ gives $(1-\Xi(y))/\xi(y)\ge\tfrac12-y$. These are the corresponding ratios under uniform noise, which proves the comparison. The comparison holds margins and values of office fixed and therefore does not order equilibrium comparative statics. For example, with $G$ uniform on $[-1.1,1.1]^2$, $\theta_L=(-1,-1)$, $\theta_R=(-\tfrac34,-\tfrac58)$, $S=\tfrac1{10}$, $B=0$, and $\phi=\tfrac32$, the policy gap rises with unpredictability under uniform noise but falls under logistic noise $\Xi(y)=1/(1+e^{-4y})$.

(e) The arguments for moderate electoral sensitivity use bounds on winning probabilities and their derivatives near the relevant nominees, together with the probability and density at a tie. Since $\xi$ is continuously differentiable near zero, $\Xi(\phi M)=\tfrac12+\phi M+O(\phi^2)$ with the required derivative bounds. At a tie, both parties still face probability $\tfrac12$ and density $\phi$, so the common-benefit tie identity is unchanged.
\end{proof}

\paragraph{Adding idiosyncratic noise.}
Retain the aggregate noise $\eta$ and let voter $i$ vote for $l$ if and only if $\Delta U_i\ge\eta+\varepsilon_i$. Suppose that $G$ is centrally symmetric about $\theta_M$ and that the $\varepsilon_i$ are independent draws, independent of types and $\eta$, from a common distribution with continuous, strictly increasing cdf $F_\varepsilon$ satisfying $F_\varepsilon(z)+F_\varepsilon(-z)=1$. Assume that idiosyncratic draws aggregate to their population probabilities, so the conditional vote share is
\[
\mathcal S_l(\eta)=\int F_\varepsilon\big(\Delta U_i-\eta\big)\,\mathrm{d}G(\theta_i).
\]
Then the winning probability is exactly the same as with aggregate noise alone.

\begin{proof}
By Lemma~\ref{lem:payoff}, the payoff gap is affine in the voter's type. Reflecting a type through $\theta_M$ therefore gives
\[
\Delta U_{2\theta_M-\theta_i}-\Delta U_{\theta_M}
=-\big(\Delta U_i-\Delta U_{\theta_M}\big).
\]
Central symmetry of $G$ and symmetry of $F_\varepsilon$ imply
\[
\mathcal S_l\big(\Delta U_{\theta_M}\big)
=\tfrac12\int\left[F_\varepsilon\big(\Delta U_i-\Delta U_{\theta_M}\big)
+F_\varepsilon\big(-\Delta U_i+\Delta U_{\theta_M}\big)\right]\mathrm{d}G(\theta_i)
=\tfrac12.
\]
Since $F_\varepsilon$ is strictly increasing, $\mathcal S_l(\eta)$ is strictly decreasing in $\eta$. Thus $l$ wins if and only if $\eta<\Delta U_{\theta_M}$, the same threshold as in Corollary~\ref{cor:median}. Keeping the distribution of aggregate noise fixed therefore leaves every pair of candidates' winning probability unchanged. Party payoffs and all equilibrium implications under their original conditions follow unchanged. Central symmetry is a sufficient condition for this argument; the median-point condition alone is not being used to assert invariance to idiosyncratic noise.
\end{proof}

\section{Beyond affine targets}

\begin{lemma}\label{lem:targets}
For any targets $\tau(\cdot,\theta_i)$ with finite second moments, with $x^*_j(\mu)=\mathbb{E}_\mu[\tau(\omega,\theta_j)]$,
\begin{align}
U_i^j&=-\int\operatorname{Var}_\mu\big(\tau(\omega,\theta_i)\big)\,\mathrm{d}F_\mu(\mu)-\Lambda_i^j,\label{eq:targets}\\
\Lambda_i^j&=\int\Big(\mathbb{E}_\mu\big[\tau(\omega,\theta_j)\big]-\mathbb{E}_\mu\big[\tau(\omega,\theta_i)\big]\Big)^2\mathrm{d}F_\mu(\mu).\nonumber
\end{align}
\end{lemma}

\begin{proof}
For each posterior, $\mathbb{E}_\mu\big[(x-\tau(\omega,\theta_i))^2\big]=\operatorname{Var}_\mu\big(\tau(\omega,\theta_i)\big)+\big(x-\mathbb{E}_\mu[\tau(\omega,\theta_i)]\big)^2$. Setting $x=x^*_j(\mu)$ and integrating over $F_\mu$ gives \eqref{eq:targets}.
\end{proof}

\paragraph{Responses with a common shape.}
Suppose that $\alpha(\omega,\theta_i)=\alpha_ih(\omega)$ for a function $h$ with $h(\mathbb{E}[\omega])=0$ and finite variance, and write $\bar h\equiv\mathbb{E}[h(\omega)]$ and $\tilde\beta_i\equiv\beta_i+\alpha_i\bar h$ for player $i$'s expected policy. Officeholder $j$ then chooses $x^*_j(\mu)=\tilde\beta_j+\alpha_j\big(\mathbb{E}_\mu[h(\omega)]-\bar h\big)$. Write $S_h\equiv\operatorname{Var}_{F_\mu}\big(\mathbb{E}_\mu[h(\omega)]\big)$. The law of iterated expectations gives $\mathbb{E}\big[(\mathbb{E}_\mu[h(\omega)]-\bar h)(h(\omega)-\bar h)\big]=S_h$, so the computation in the proof of Lemma~\ref{lem:payoff} gives
\[
U_i^j=-\alpha_i^2\operatorname{Var}\big(h(\omega)\big)-(\tilde\beta_j-\tilde\beta_i)^2+\alpha_j(2\alpha_i-\alpha_j)\,S_h .
\]
Every result in the main text depends on the prior and the information environment only through Lemma~\ref{lem:payoff}, so each holds with $\tilde\beta_i$ in place of $\beta_i$, $S_h$ in place of $S$, and $\operatorname{Var}(h(\omega))$ in place of $\operatorname{Var}(\omega)$. When $h$ is nonlinear, a change in the prior can also change $\bar h$ and hence the expected policies, an effect on levels that linear responses rule out. The exception is the normal example used for endogenous learning, which sets $h(\omega)=\omega$.

The first term in \eqref{eq:targets} is a player's own residual uncertainty, which she would bear even if she governed herself; the second is her disagreement with $j$. Since $\Lambda_i^j$ is the expectation of a convex function of the posterior, and a Blackwell improvement spreads the distribution of posteriors \citep{Blackwell1953}, every Blackwell improvement raises $\Lambda_i^j$, as it does the expectation of any convex function of the policy gap. The effects of information in Propositions~\ref{prop:office} and~\ref{prop:auto-polar} therefore hold for any target. The effects of uncertainty depend on the shape of the targets. With full information, $\Lambda_i^j=\mathbb{E}[\Lambda_i^j(\omega)]$, where $\Lambda_i^j(\omega)\equiv\big(\tau(\omega,\theta_j)-\tau(\omega,\theta_i)\big)^2$ is the cost of delegation given the state, so greater uncertainty raises disagreement when it moves probability mass toward states in which $\Lambda_i^j(\omega)$ is large. With affine targets, $\Lambda_i^j(\omega)$ is convex in $\omega$, and every mean-preserving spread \citep{RothschildStiglitz1970} of $\omega$ raises $\Lambda_i^j$. With saturating targets, $\Lambda_i^j(\omega)$ can be largest near the prior and vanish in extreme states, as when two players want drastic action at different thresholds but agree once a crisis is severe. For such targets, the form of the dispersion matters more than its variance. For two prior distributions $F_0^1$ and $F_0^2$ of the state, say that $F_0^2$ is a \emph{uniform dispersion} of $F_0^1$ if $F_0^2(\omega)\ge F_0^1(\omega)$ for $\omega<0$ and $F_0^2(\omega)\le F_0^1(\omega)$ for $\omega>0$, so that every quantile of the state moves weakly away from the prior mean, as when the state is scaled up. A dispersion that fails this condition is \emph{local}: it can move mass toward the prior mean as well as into the tails, as when ordinary times become calmer and rare crises more severe.

\begin{proposition}[Uniform and local dispersions]\label{prop:dispersion}
Suppose that the officeholder learns the state, and that $\Lambda_i^j(\omega)$ is bounded, continuous, and single-peaked in $\omega$ at the prior mean: nondecreasing on $(-\infty,0]$ and nonincreasing on $[0,\infty)$.
\begin{enumerate}
  \item[(a)] If $F_0^2$ is a uniform dispersion of $F_0^1$, then $\Lambda_i^j$ is weakly lower under $F_0^2$ than under $F_0^1$.
  \item[(b)] If $F_0^2$ is not a uniform dispersion of $F_0^1$, then there are targets whose squared gap is bounded, continuous, and single-peaked at the prior mean for which $\Lambda_i^j$ is strictly higher under $F_0^2$ than under $F_0^1$.
  \item[(c)] With affine targets, every mean-preserving spread of $\omega$ raises $\Lambda_i^j$.
\end{enumerate}
\end{proposition}

\begin{proof}
(a) Since $\Lambda_i^j(\omega)$ is continuous and of bounded variation in $\omega$, and $F_0^2-F_0^1$ vanishes at $\pm\infty$, integration by parts gives
\[
\mathbb{E}_{F_0^2}\big[\Lambda_i^j(\omega)\big]-\mathbb{E}_{F_0^1}\big[\Lambda_i^j(\omega)\big]=\int\Lambda_i^j(\omega)\,\mathrm{d}(F_0^2-F_0^1)(\omega)=-\int\big(F_0^2(\omega)-F_0^1(\omega)\big)\,\mathrm{d}\Lambda_i^j(\omega).
\]
On $(-\infty,0)$, $F_0^2-F_0^1\ge0$ and $\Lambda_i^j(\omega)$ is nondecreasing; on $(0,\infty)$, $F_0^2-F_0^1\le0$ and $\Lambda_i^j(\omega)$ is nonincreasing. The integrand is therefore nonnegative on both half-lines, and $\Lambda_i^j$ is weakly lower under $F_0^2$.

(b) If $F_0^2$ is not a uniform dispersion of $F_0^1$, then $F_0^2(t)>F_0^1(t)$ for some $t>0$, or $F_0^2(t)<F_0^1(t)$ for some $t<0$. In the first case, let $\delta=F_0^2(t)-F_0^1(t)>0$ and choose $t'<0$ with $F_0^2(t')<\delta/4$. For $\lambda>0$, choose targets $\tau(\omega,\theta_i)=0$ and $\tau(\omega,\theta_j)\ge0$ such that $\Lambda_i^j(\omega)=\tau(\omega,\theta_j)^2$ equals one on $[t',t]$, zero outside $[t'-\lambda,t+\lambda]$, and is linear in between. Because $t'<0<t$, $\Lambda_i^j(\omega)$ is bounded, continuous, and single-peaked at zero. As $\lambda\to0$, $\Lambda_i^j(\omega)$ converges pointwise to the indicator of $[t',t]$, so by dominated convergence
\[
\mathbb{E}_{F_0^2}\big[\Lambda_i^j(\omega)\big]-\mathbb{E}_{F_0^1}\big[\Lambda_i^j(\omega)\big]\to\big(F_0^2(t)-F_0^1(t)\big)-\big(F_0^2(t'^-)-F_0^1(t'^-)\big)\ge\delta-F_0^2(t')>\tfrac34\delta>0.
\]
The difference is therefore positive for $\lambda$ small enough. The second case is symmetric.

(c) With affine targets, $\Lambda_i^j(\omega)=\big(\beta_j-\beta_i+(\alpha_j-\alpha_i)\omega\big)^2$ is convex, so a mean-preserving spread raises its expectation \citep{RothschildStiglitz1970}.
\end{proof}

Substantively, a uniform dispersion makes every period more turbulent, so players spend more time in the extreme states in which saturating targets agree. A local dispersion leaves most periods calm and makes rare crises more severe, so it can raise disagreement. The variance rises by the same amount in both cases. Figure~\ref{fig:targets} illustrates both cases.

\begin{figure}[t]
\centering
\begin{tikzpicture}
\pgfplotsset{tgt/.style={width=7.8cm, height=5.2cm, title style={font=\small}, label style={font=\small}, tick label style={font=\small}, every axis plot/.append style={thick}, legend style={font=\small, draw=none, at={(0.5,-0.32)}, anchor=north, legend columns=-1, column sep=6pt}, legend cell align=left}}
\begin{axis}[tgt, name=pa, xmin=-3.5, xmax=3.5, ymin=0, ymax=3, axis lines=left, xlabel={state $\omega$}, ytick={0,1,2,3}, title={(a) Disagreement conditional on the state}, clip=true]
\addplot[black, domain=-3.5:3.5, samples=120] {x^2};
\addplot[black, dashed, domain=-3.5:3.5, samples=200] {(tanh(x+1)-tanh(x-1))^2};
\legend{affine, saturating}
\end{axis}
\begin{axis}[tgt, name=pb, at={(pa.east)}, anchor=west, xshift=1.2cm, xmin=-5, xmax=5, ymin=0, ymax=0.6, axis lines=left, xlabel={state $\omega$}, ytick=\empty, title={(b) Distributions of the state}]
\addplot[gray!80!black, domain=-5:5, samples=150] {exp(-x^2/2)/sqrt(2*pi)};
\addplot[gray!80!black, dashed, domain=-5:5, samples=150] {exp(-x^2/4.5)/(1.5*sqrt(2*pi))};
\addplot[gray!80!black, dotted, domain=-5:5, samples=200] {0.8*exp(-x^2/0.72)/(0.6*sqrt(2*pi))+0.2*exp(-x^2/19.62)/(3.132*sqrt(2*pi))};
\legend{prior, uniform, local}
\end{axis}
\begin{axis}[tgt, name=pc, at={(pa.south)}, anchor=north, yshift=-3.0cm, xmin=0, xmax=1, ymin=0, ymax=1.6, axis lines=left, xlabel={revealed share $R^2$}, ylabel={expected disagreement}, title={(c) Better information}]
\addplot[black, domain=0:1, samples=2] {x};
\addplot[black, dashed] coordinates {(0.0,1.2118) (0.09,1.2129) (0.18,1.2164) (0.27,1.2222) (0.36,1.2305) (0.45,1.2415) (0.54,1.2553) (0.63,1.2723) (0.72,1.2927) (0.81,1.3171) (0.9,1.3461) (0.99,1.3804)};
\legend{affine, saturating}
\end{axis}
\begin{axis}[tgt, name=pd, at={(pc.east)}, anchor=west, xshift=1.2cm, ybar, bar width=9pt, ymin=0, ymax=2.6, symbolic x coords={prior,uniform,local}, xtick=data, xticklabels={prior, uniform, local}, axis lines*=left, ylabel={expected disagreement}, title={(d) Greater uncertainty}, enlarge x limits=0.25]
\addplot[fill=black, draw=black] coordinates {(prior,1) (uniform,2.25) (local,2.25)};
\addplot[fill=white, draw=black, pattern=north east lines] coordinates {(prior,1.385) (uniform,1.028) (local,1.556)};
\legend{affine, saturating}
\end{axis}
\end{tikzpicture}
\caption{Affine and saturating targets. The affine targets are $\tau(\omega,\theta)=\alpha\omega$ with a gap in responsiveness of one. The saturating targets are $\tau(\omega,\theta_i)=\tanh(\omega+1)$ and $\tau(\omega,\theta_j)=\tanh(\omega-1)$: the two players want drastic action at different thresholds, so their gap is largest at the prior and vanishes in extreme states. Panel (b) shows the prior $\operatorname{Norm}(0,1)$, a uniform dispersion $\operatorname{Norm}(0,2.25)$, and a local dispersion with the same variance, which puts more mass near the prior and in the tails. In panel (c), better information raises expected disagreement for both targets. In panel (d), with full information, both dispersions raise disagreement for affine targets; for saturating targets, the uniform dispersion lowers it and the local dispersion raises it.}
\label{fig:targets}
\end{figure} Proposition~\ref{prop:realign} uses the two-dimensional representation of types, and so the affine form.


\section{Measuring responsiveness and testing the model}\label{app:tests}

\paragraph{A vote equation.} With independent issues (Proposition~\ref{prop:issues}), voter $i$ votes for candidate $r$ when $U_i^r-U_i^l$ exceeds the electoral noise. Expanding the costs of delegation gives
\[
\begin{aligned}
\Pr(\text{vote } r)&=\phi\Big(\kappa+2\sum_n(\beta_{rn}-\beta_{ln})\,\beta_{in}+2\sum_n(\alpha_{rn}-\alpha_{ln})\,S_n\,\alpha_{in}\Big),\\
\kappa&\equiv\sum_n\Big[(\beta_{ln}^2-\beta_{rn}^2)+(\alpha_{ln}^2-\alpha_{rn}^2)\,S_n\Big],
\end{aligned}
\]
where, with an abuse of notation, $\phi(\cdot)$ is the distribution function of the noise. With uniform noise, $\phi(z)=\tfrac12+\phi z$ on its support. In the model, the noise is common to all voters. An empirical version adds idiosyncratic taste shocks, which makes $\phi(\cdot)$ logistic or normal, and absorbs the common shock into an intercept for each election. The constant $\kappa$ moves with $S_n$, so tests should rely on the slopes, with fixed effects for each election and issue absorbing $\kappa$. Because the scale of $\phi(\cdot)$ is not identified, the prediction concerns a ratio. With the candidates fixed, the ratio of the coefficient on $\alpha_{in}$ to the coefficient on $\beta_{in}$, $(\alpha_{rn}-\alpha_{ln})S_n/(\beta_{rn}-\beta_{ln})$, rises in proportion to $S_n$. When the candidates change between elections, their responsiveness must be measured separately, for example with expert surveys, to separate the candidates' gap in responsiveness from $S_n$.

Surveys sometimes ask about responses, and recent work comes close to the distinction without naming it. Some surveys ask explicitly, by specifying a projected state, such as a rise in temperatures, in migration, or in the automation of jobs, and asking which policy the respondent would support. Questions about uncertain domains often do so implicitly. Information-provision experiments tell respondents about a state, such as intergenerational mobility or the number of immigrants, and measure how their policy preferences change \citep{AlesinaStantchevaTeso2018,AlesinaMianoStantcheva2023,HaalandRothWohlfart2023}. Heterogeneity in those effects is heterogeneity in responsiveness. In the debate over moderates, scholars ask whether voters who appear centrist hold centrist views on each issue \citep{Fowler2023} or extreme but cross-cutting ones \citep{Broockman2016,AhlerBroockman2018,Broockman2025}. In the model, both descriptions can fit the same voter: a voter with a moderate baseline and a strong desired response to one state looks centrist in ordinary times and extreme once that state becomes uncertain. Conjoint estimates combine the direction and the intensity of preferences \citep{AbramsonKocakMagazinnik2022}, and the model implies that the intensity attached to a candidate's response rises with uncertainty. These approaches lack a framework that identifies which component a measure captures and when each component governs the vote. The model supplies such a framework. The decomposition can be made explicit by presenting states $s_1,\dots,s_M$ within one domain and eliciting respondent $i$'s preferred policy $x_{im}$ at each state $s_m$. Fitting $x_{im}$ to $\beta_i+\alpha_is_m$ by least squares estimates $\beta_i$ by the intercept and $\alpha_i$ by the slope. With common policy units and reference states, the same scenarios can elicit the policies that respondents expect each candidate to choose, which estimates $(\beta_j,\alpha_j)$. The model's predictions can then be checked. As $S$ rises, $\Pr(\text{vote } r)$ should rise for voters with $(\alpha_i-\alpha_r)^2<(\alpha_i-\alpha_l)^2$ and fall for the others (Lemma~\ref{lem:payoff}). The candidates' types should converge, $\theta_r-\theta_l=w(\theta_R-\theta_L)$ with $w$ falling in $S$, while $D^*$ rises (Proposition~\ref{prop:symmetric}). Voters' $\alpha_{in}$ should predict their realignment after a shock to $\omega_n$, and the new alignment should cut across the old one when $\alpha_{in}$ is weakly correlated with $\beta_{in}$ across voters (Remark~\ref{rem:crosscut}).

The predictions vary across regions because regions differ in their exposure to a shock. Suppose that the voters in a region care about a local state $e\omega$. The parameter $e$ is the region's exposure to the shock, such as its dependence on an industry that automation threatens or its exposure to trade or to climate risk. A voter's effective responsiveness is then $\alpha_ie$, and her comparison of the candidates depends on exposure exactly as it depends on responsiveness. With several states, region $q$'s local conditions are $\sum_ne_{qn}\omega_n$, and the region can be highly exposed to one state and barely exposed to another (Proposition~\ref{prop:multistate}). Because the states are independent,
\[
\operatorname{Var}\Big(\sum_ne_{qn}\omega_n\Big)=\sum_ne_{qn}^2\operatorname{Var}(\omega_n).
\]
This sum measures the total uncertainty that region $q$ faces, and it can be compared across shocks. A shock raises the weight on responsiveness most in the regions most exposed to it. Whether those regions see the most switching depends on their initial alignment and on the distribution of switching thresholds: where many voters have already passed their thresholds, a further rise in exposure produces fewer additional switches. With voters uniform on a square and a dividing line $\beta=-t\alpha$, for example, the share who switch when $t$ doubles is $1/(8t)$ for $t\ge1$, which falls from 12.5 percent between $t=1$ and $t=2$ to about 3 percent between $t=4$ and $t=8$. Greater switching in exposed regions is therefore a conditional prediction.

\paragraph{Separating uncertainty from levels.}
Most shocks change both the level of local conditions and their uncertainty. Existing estimates of the political effect of a shock include both changes. Exposure to Chinese imports cost jobs, and exposed workers also changed jobs and employers more often \citep{AutorDornHansonSong2014}. The shock began when the United States granted China permanent normal trade relations in 2000. The decision removed the risk of tariff increases for importers \citep{PierceSchott2016,HandleyLimao2017} and made conditions more volatile for exposed workers. In the model, $\omega$ is the state that the officeholder faces after the election, with $\mathbb{E}[\omega]=0$, and conditions realized before the election are absorbed into baselines; only the timing extension lets information arrive before nomination. A shock such as the China shock therefore involves two states. Let $\omega_1$ be manufacturing employment, whose fall $\omega_{01}$ was realized before the election, and let $\omega_2$ be the future of alternative employment in the region, which remains uncertain. The realized fall moves voter $i$'s target for policy~1 by $\alpha_{i1}\omega_{01}$ and candidate $j$'s post-election policy $x^*_{j1}(\mu)$ by $\alpha_{j1}\omega_{01}$: it changes the baselines from $\beta_{i1}$ to $\beta_{i1}+\alpha_{i1}\omega_{01}$ and from $\beta_{j1}$ to $\beta_{j1}+\alpha_{j1}\omega_{01}$. The uncertainty about $\omega_2$ raises $S_2$, and with it $\operatorname{Var}\big(x^*_{j2}(\mu)\big)=\alpha_{j2}^2S_2$, the extent to which the post-election policy moves with the state. In voter $i$'s comparison, the realized fall moves the baselines and leaves the slope in $\alpha_{i2}$, $2(\alpha_{r2}-\alpha_{l2})S_2$, unchanged. The rise in $S_2$ raises that slope and leaves the baselines unchanged. Because baselines move in proportion to $\alpha_{i1}$ and $\alpha_{j1}$, the realized fall can also move voters in both directions and change the distance between the candidates' baselines, so switching in both directions and hardening do not separate the two changes by themselves. The coefficient on $\alpha_{i2}$, holding fixed the baselines measured after the shock, separates them.

Three comparisons can isolate the change in that coefficient. The first comparison involves three groups: residents of exposed regions whose own job or industry was hit, residents of the same regions whose incomes were not affected, such as public employees and retirees, and residents of unexposed regions. Residents who were hit indirectly face a higher $S_2$ for their region's future, and their own baseline on policy~1 changes little. The model predicts that the coefficient on $\alpha_{i2}$ rises for them, and a pure level account predicts that it does not change. Local spillovers and sociotropic voting move baselines and leave that coefficient unchanged, so they do not confound the comparison. With panel data, a researcher can assign residents to groups by their residence and occupation before the shock and so distinguish switching from migration. The second comparison uses timing, because $S_2$ often rises before $\omega_{01}$ is realized. Uncertainty rose at the announcements of the 2018--2019 trade war. It stayed high among firms in the three and a half years between the Brexit referendum and the United Kingdom's departure \citep{BloomEtAl2019}. The release of ChatGPT in November 2022 changed the prospects of exposed occupations \citep{EloundouEtAl2024,FeltenRajSeamans2021}. The measure of total uncertainty defined above can be computed from exposure before the shock and from national indices of uncertainty \citep{BakerBloomDavis2016,CaldaraEtAl2020}. This measure varies over time in a way that realized losses do not. Exposure to generative AI is concentrated among educated professionals, so the model predicts realignment in a different group from the one that the China shock affected. The third comparison holds the level fixed. Regions with similar realized losses can differ in how much $e_{q2}^2\operatorname{Var}(\omega_2)$ rose, for example because their remaining industries differ in diversity. The coefficient on $\alpha_{i2}$ should rise more where that term rose more.

Existing panels do not measure responsiveness, so they can test only the weaker implications. With survey panels that record occupation and past votes, such as the Cooperative Election Study and the British Election Study Internet Panel, a researcher can test whether switching in both directions was larger in exposed areas. The model predicts this pattern, and a realized loss can also produce it when $\alpha_{r1}\ne\alpha_{l1}$. Tests of the coefficient on $\alpha_{i2}$ require the scenario questions described above. In the United States, the two directions are from Obama to Trump and from Romney to Clinton; in Britain, they are from Labour to the Conservatives and from the Conservatives to the Liberal Democrats. Leave and Remain identities became strong affective attachments after the referendum \citep{HoboltLeeperTilley2021}, as the model predicts.

\renewcommand{\refname}{References for the Supplementary Information}
\begin{thebibliography}{99}
\setlength{\itemsep}{2pt}

\bibitem[{Abramson, Ko\c{c}ak and Magazinnik(2022)}]{AbramsonKocakMagazinnik2022}
Abramson, S. F., K. Ko\c{c}ak, and A. Magazinnik. 2022. ``What Do We Learn about Voter Preferences from Conjoint Experiments?'' \textit{American Journal of Political Science} 66 (4): 1008--1020.

\bibitem[{Ahler and Broockman(2018)}]{AhlerBroockman2018}
Ahler, D. J., and D. E. Broockman. 2018. ``The Delegate Paradox: Why Polarized Politicians Can Represent Citizens Best.'' \textit{Journal of Politics} 80 (4): 1117--1133.

\bibitem[{Alesina, Miano, and Stantcheva(2023)}]{AlesinaMianoStantcheva2023}
Alesina, A., A. Miano, and S. Stantcheva. 2023. ``Immigration and Redistribution.'' \textit{Review of Economic Studies} 90 (1): 1--39.

\bibitem[{Alesina, Stantcheva, and Teso(2018)}]{AlesinaStantchevaTeso2018}
Alesina, A., S. Stantcheva, and E. Teso. 2018. ``Intergenerational Mobility and Preferences for Redistribution.'' \textit{American Economic Review} 108 (2): 521--554.

\bibitem[{Autor, Dorn, Hanson, and Song(2014)}]{AutorDornHansonSong2014}
Autor, D. H., D. Dorn, G. H. Hanson, and J. Song. 2014. ``Trade Adjustment: Worker-Level Evidence.'' \textit{Quarterly Journal of Economics} 129 (4): 1799--1860.

\bibitem[{Bagnoli and Bergstrom(2005)}]{BagnoliBergstrom2005}
Bagnoli, M., and T. Bergstrom. 2005. ``Log-Concave Probability and Its Applications.'' \textit{Economic Theory} 26 (2): 445--469.

\bibitem[{Baker, Bloom, and Davis(2016)}]{BakerBloomDavis2016}
Baker, S. R., N. Bloom, and S. J. Davis. 2016. ``Measuring Economic Policy Uncertainty.'' \textit{Quarterly Journal of Economics} 131 (4): 1593--1636.

\bibitem[{Berge(1963)}]{Berge1963}
Berge, C. 1963. \textit{Topological Spaces}. Edinburgh: Oliver and Boyd.

\bibitem[{Billingsley(1995)}]{Billingsley1995}
Billingsley, P. 1995. \textit{Probability and Measure}, 3rd ed. New York: Wiley.

\bibitem[{Blackwell(1953)}]{Blackwell1953}
Blackwell, D. 1953. ``Equivalent Comparisons of Experiments.'' \textit{Annals of Mathematical Statistics} 24 (2): 265--272.

\bibitem[{Bloom et al.(2019)}]{BloomEtAl2019}
Bloom, N., P. Bunn, S. Chen, P. Mizen, P. Smietanka, and G. Thwaites. 2019. ``The Impact of Brexit on UK Firms.'' NBER Working Paper 26218.

\bibitem[{Border(1985)}]{Border1985}
Border, K. C. 1985. \textit{Fixed Point Theorems with Applications to Economics and Game Theory}. Cambridge: Cambridge University Press.

\bibitem[{Broockman(2016)}]{Broockman2016}
Broockman, D. E. 2016. ``Approaches to Studying Policy Representation.'' \textit{Legislative Studies Quarterly} 41 (1): 181--215.

\bibitem[{Broockman(2025)}]{Broockman2025}
Broockman, D. E. 2025. ``\,`Moderates'\,.'' \textit{American Political Science Review}, published online January 20.

\bibitem[{Caldara et al.(2020)}]{CaldaraEtAl2020}
Caldara, D., M. Iacoviello, P. Molligo, A. Prestipino, and A. Raffo. 2020. ``The Economic Effects of Trade Policy Uncertainty.'' \textit{Journal of Monetary Economics} 109: 38--59.

\bibitem[{Eloundou et al.(2024)}]{EloundouEtAl2024}
Eloundou, T., S. Manning, P. Mishkin, and D. Rock. 2024. ``GPTs Are GPTs: Labor Market Impact Potential of LLMs.'' \textit{Science} 384 (6702): 1306--1308.

\bibitem[{Evans and Gariepy(2015)}]{EvansGariepy2015}
Evans, L. C., and R. F. Gariepy. 2015. \textit{Measure Theory and Fine Properties of Functions}, revised ed. Boca Raton: CRC Press.

\bibitem[{Felten, Raj, and Seamans(2021)}]{FeltenRajSeamans2021}
Felten, E., M. Raj, and R. Seamans. 2021. ``Occupational, Industry, and Geographic Exposure to Artificial Intelligence: A Novel Dataset and Its Potential Uses.'' \textit{Strategic Management Journal} 42 (12): 2195--2217.

\bibitem[{Fowler et~al.(2023)}]{Fowler2023}
Fowler, A., S. J. Hill, J. B. Lewis, C. Tausanovitch, L. Vavreck, and C. Warshaw. 2023. ``Moderates.'' \textit{American Political Science Review} 117 (2): 643--660.

\bibitem[{Haaland, Roth, and Wohlfart(2023)}]{HaalandRothWohlfart2023}
Haaland, I., C. Roth, and J. Wohlfart. 2023. ``Designing Information Provision Experiments.'' \textit{Journal of Economic Literature} 61 (1): 3--40.

\bibitem[{Handley and Lim\~{a}o(2017)}]{HandleyLimao2017}
Handley, K., and N. Lim\~{a}o. 2017. ``Policy Uncertainty, Trade, and Welfare: Theory and Evidence for China and the United States.'' \textit{American Economic Review} 107 (9): 2731--2783.

\bibitem[{Milgrom and Shannon(1994)}]{MilgromShannon1994}
Milgrom, P., and C. Shannon. 1994. ``Monotone Comparative Statics.'' \textit{Econometrica} 62 (1): 157--180.

\bibitem[{Ok(2007)}]{Ok2007}
Ok, E. A. 2007. \textit{Real Analysis with Economic Applications}. Princeton: Princeton University Press.

\bibitem[{Pierce and Schott(2016)}]{PierceSchott2016}
Pierce, J. R., and P. K. Schott. 2016. ``The Surprisingly Swift Decline of US Manufacturing Employment.'' \textit{American Economic Review} 106 (7): 1632--1662.

\bibitem[{Rothschild and Stiglitz(1970)}]{RothschildStiglitz1970}
Rothschild, M., and J. E. Stiglitz. 1970. ``Increasing Risk: I. A Definition.'' \textit{Journal of Economic Theory} 2 (3): 225--243.

\bibitem[{Sundaram(1996)}]{Sundaram1996}
Sundaram, R. K. 1996. \textit{A First Course in Optimization Theory}. Cambridge: Cambridge University Press.

\bibitem[{Topkis(1998)}]{Topkis1998}
Topkis, D. M. 1998. \textit{Supermodularity and Complementarity}. Princeton: Princeton University Press.

\end{thebibliography}

\end{document}